\section{Surjectivity I: runs and index bounds}\label{sec:pf:onto1}

\providecommand{\ontooneidx}{\operatorname{idx}}
\providecommand{\ontooneql}{\operatorname{ql}}
\providecommand{\ontoonesize}{\operatorname{size}}
\providecommand{\ontoonemls}{\operatorname{mls}}
\providecommand{\ontooneMsup}{M^{\mathrm{sup}}}
\providecommand{\ontoonePt}{\operatorname{Pt}}
\providecommand{\ontooneDF}{\operatorname{DF}}
\providecommand{\ontoonedepth}{\operatorname{depth}}
\providecommand{\ontoonerk}{\operatorname{rk}}
\providecommand{\ontooneord}{\operatorname{ord}}
\providecommand{\ontoonelg}{\operatorname{lg}}
\providecommand{\ontoonePP}{\mathsf{P}}
\providecommand{\ontoonePPz}{\mathsf{P}_0}
\providecommand{\ontoonebin}{\binom{T-1}{\delta}}
\providecommand{\ontooneIB}{\mathrm{IB}}
\providecommand{\ontooneet}{\eta}

This section proves two results of the formal development \cite{LeanSource2026} that enter the
proof of the surjectivity estimate $\sfname{SurjectiveOntoX}$ of \cref{prop:families:points}: an
\emph{index bound along runs} (\cref{pf:onto1:runholds}) and the existence of \emph{admissible generating
families of scaled theta lattices} (\cref{pf:onto1:genok}). Both concern a \emph{record}: a number field $K$, finite
dimensional $K$-spaces $\mathcal F_w$ of cusp forms indexed by weights $w$, lattices $\mathcal M_w\subseteq\mathcal
F_w$ of integral forms, and, at every finite place $\mathfrak l$ of $K$, a family of valuations of forms
attached to points. The index bound compares two lattices of forms of a large weight $\kappa$: the
lattice $\mathcal M_\kappa$ of integral forms and a \emph{theta lattice} $\Theta$ generated by products
$W^T\cdot\mathrm{Th}_F$ of a fixed lattice $W$ of forms of a small weight $w_L$ (scaled theta forms) with forms
of a complementary weight. It asserts that the $\mathfrak l$-adic index of $\Theta$ in $\mathcal M_\kappa$,
measured on any subspace $V'$, is at most $e\dim V'+r(\mathfrak l)$ for every integer $e\ge T\nu(\mathfrak l)$,
where $\nu(\mathfrak l)$ is an explicit local depth and the local defects $r(\mathfrak l)\ge0$ satisfy the
global bound $\sum_{\mathfrak l}r(\mathfrak l)\log\Nm\mathfrak l/[K:\Q]\le\rank(\mathcal M_\kappa)\,f\,m\,b\,\delta_2$.

The proof has an abstract core, Theorem Zero (\cref{pf:onto1:thmzero}). Locally at $\mathfrak l$ one passes to
the Gauss extension of the completed ring $\OO_{\mathfrak l}$ in indeterminates $X_{i,k}$, forms $\delta+1$
generic linear combinations $s_0,\dots,s_\delta$ of generators of $W$ ($\delta$ is the dimension of the
relevant moduli space), and studies sections modulo the $a_0$-th power of the ideal of the
zero-dimensional scheme $s_1=\dots=s_\delta=0$ in the chart $s_0\ne0$, where $a_0=T-\delta$. A chain of
submodules and a length computation over discrete valuation rings (\cref{pf:onto1:levelsQDesc}) bound
the index by four lengths: a pole length governed by the depths at the points of the zero-dimensional
section, a base length governed by a local height $I_{\mathfrak l}$ of auxiliary forms $\theta'_i$, a Gauss
content of an auxiliary determinant, and the length of a quotient by the forms whose low
Fourier--Jacobi coefficients vanish. Globally, the sum of these defects is controlled by height bounds, by
the product formula applied to the auxiliary determinant, and by explicit numerical inequalities. The
remaining work consists in producing, for the record attached to a run, all the inputs of Theorem Zero; the main inputs
are proved at the level of their main steps in the last subsections, deeper ingredients being cited.

\paragraph{Organisation.} \Cref{sec:pf:onto1:setting} introduces weights, form providers, valued form
algebras, theta lattices and the index bound. \Cref{sec:pf:onto1:dvr} proves the length computations
over discrete valuation rings. \Cref{sec:pf:onto1:local} sets up the generic linear section at a prime and
proves the local bound, and \cref{sec:pf:onto1:lengths} the base and pole length bounds which it uses.
\Cref{sec:pf:onto1:thmzero} proves Theorem Zero, and \cref{sec:pf:onto1:gausssum} the bound on the Gauss contents.
\Cref{sec:pf:onto1:net} replaces the $L^2$-norm by a discrete norm on a finite net and converts the inputs into the form
required by Theorem Zero. \Cref{sec:pf:onto1:record} treats the record of a run and proves the index bound along
runs, and \cref{sec:pf:onto1:gen} constructs the admissible generating families. \Cref{sec:pf:onto1:inputs} proves the
inputs at the record used in \cref{sec:pf:onto1:record}, and \cref{sec:pf:onto1:frames} the lower bounds at frames used in
\cref{sec:pf:onto1:gen}.

\paragraph{Conventions.} Lengths of modules take values in $\N\cup\{\infty\}$; for $n\in\N\cup\{\infty\}$ we
write $n^\flat=n$ if $n<\infty$ and $n^\flat=0$ if $n=\infty$ (this is how the development reads a length
as a real number; in all uses below the length is first shown to be finite). Valuations of forms take
values in the extended real line $[-\infty,+\infty]$; a value $v\ne\pm\infty$ is used as a real number. As
in the development, $\log0=0$. For a prime $\mathfrak l$ of a number field $K$ we write
$\ontooneord_{\mathfrak l}$ for the normalised additive valuation, $\Nm\mathfrak l$ for the absolute norm,
$e_{\mathfrak l}$ for the ramification index over $\Q$ and
\[
\ell(\mathfrak l)=\frac{\log\Nm\mathfrak l}{[K:\Q]}.
\]
For natural numbers, $T-\delta$ denotes truncated subtraction. Sums $\sum_{\mathfrak l}$ over the primes of
$K$ are sums of summable families of reals. Conditions are named by short symbols, listed in the
definitions where they are introduced.

%----------------------------------------------------------------------------------------------------
\subsection{Weights, form providers and valued form algebras}\label{sec:pf:onto1:setting}

\begin{definition}[Datum and size quantities]\label{pf:onto1:sizeq}
Throughout, $D$ is a datum of the formal development. We use only the following quantities attached to
it: natural numbers $f,r,m,k,t,n_V$ (ranks); for each $j$ a signature $(p_j,q_j)$; a finite set $\mathcal T$ of
types; a CM field $E$ (the host field) with discriminant $d_E$ and a real number $\mathrm{rd}_E$; natural
numbers $g_0,g_1$ (two auxiliary primes) and $q,d,Q_S,p_{\mathrm{de}},p'_{\mathrm{de}}$. An index $j$ is \emph{in
range} if $1\le j\le m$. For $j$ in range put
\[
\delta=\delta_j=f\Bigl(p_jq_j+\frac{p_j(p_j-1)}2+\frac{q_j(q_j-1)}2\Bigr)\qquad(\text{integral divisions}),
\qquad c_0=2\cdot5^4\,r .
\]
A \emph{size datum} $L$ provides a scale $x>1$, a real number $\eps\in(0,\tfrac1{52})$ and real numbers
$\delta_2$, $\log\Omega$, $\log\mathrm{Der}$ (among others). Put $\mathrm L=x^{\eps}$, $b=e^{\mathrm L}$, $z_B=fr\mathrm L$
and, for real $C,A$,
\begin{gather*}
\ontoonePP(C,A)=C\,z_B^{A},\qquad \ontoonePPz(C,A)=C\bigl(fr+\mathrm L+\log|d_E|+|\log\Omega|\bigr)^{A},\qquad
R(C)=\mathrm{rd}_E^{\,C},\\
\lambda^*=\log(g_0g_1)+\log\bigl(4(f/2)^2qd\bigr),\qquad \mathrm{PN}=2\mathrm L+\lambda^*+4,\qquad
\mathrm{PO}=19fr^2\mathrm L+198f^2\mathrm L .
\end{gather*}
\formal{Formalism.BasicNotions.D7}
\formal{Formalism.BasicNotions.SizeData₂}
\formal{Formalism.MainNotions.WeightFamilies.dimY}
\formal{Formalism.BasicNotions.zBase}
\formal{Formalism.MainNotions.WeightFamilies.PPv}
\formal{Formalism.MainNotions.WeightFamilies.PP0v}
\formal{Formalism.MainNotions.WeightAlgebra.lamStar}
\formal{Formalism.MainNotions.WeightAlgebra.PO}
\end{definition}

These are the only features of $D$ and $L$ that appear in the statements of this section; further
explicit numerical conditions on them are packaged in named hypotheses, described where they are used.

\begin{definition}[Weights]\label{pf:onto1:weights}
Let $j$ be in range. A \emph{weight} $w$ consists of non-increasing functions
$a^+_\tau\colon\{0,\dots,p_j-1\}\to\N$ and $a^-_\tau\colon\{0,\dots,q_j-1\}\to\N$ for $\tau\in\mathcal T$, and an
even natural number $n_w$. Weights are added and multiplied by natural numbers componentwise. Put
\[
\ontooneDF(w)=\sum_{\tau}\Bigl(\sum_ia^+_\tau(i)+\sum_ia^-_\tau(i)\Bigr),\qquad
\ontoonesize(w)=\ontooneDF(w)+f\max(q_j,1)\,r\,n_w+1\ \ (\ge1).
\]
Given weights $w_L,w_F$ put $w_0=w_F+\delta w_L$, and for $i\in\N$ let $\epsilon_i$ be the weight with all
$a^\pm_\tau=0$ and $n=c_0i$.
\formal{Formalism.MainNotions.WeightFamilies.Wt}
\formal{Formalism.MainNotions.WeightFamilies.Wt.size}
\formal{Formalism.MainNotions.WeightAlgebra.DF}
\end{definition}

\begin{definition}[Form providers]\label{pf:onto1:prov}
A \emph{form provider} $B$ over $D$ consists of:
\begin{enumerate}[label=(\roman*)]
\item a number field $K$ together with an identification of $\OO_K$ with a ring attached to $D$;
\item for each $j$ and each weight $w$ a finite dimensional $K$-vector space $\mathcal F_w$ (the \emph{forms}
of weight $w$; the index $j$ is suppressed) and an $\OO_K$-submodule $\mathcal M_w\subseteq\mathcal F_w$ (the
\emph{integral forms});
\item for every embedding $\sigma\colon K\to\C$ two norms $\|\cdot\|_\sigma$ (the \emph{sup norm}) and
$\|\cdot\|_{2,\sigma}$ (the \emph{$L^2$-norm}) on each $\mathcal F_w$, $\sigma$-homogeneous
($\|aF\|=|\sigma(a)|\,\|F\|$), vanishing only at $0$, with $\|F\|_{2,\sigma}\le\|F\|_\sigma$;
\item a set of \emph{cusps}, for each cusp $u$ an $\OO_E$-lattice $X_u\subseteq E^{q_j}$, and for each cusp $u$,
each hermitian matrix $\chi$ positive definite at all embeddings and such that $g_0g_1\,\xi\chi\bar\xi'^{\,t}\in\OO_E$
for all $\xi,\xi'$ in the $\OO_E$-dual of $X_u$, and each weight $w$, a $K$-linear map
$\mathrm{FJ}_{u,\chi}\colon\mathcal F_w\to\mathcal F_{u,\chi,w}$ (the \emph{Fourier--Jacobi coefficient} of index $\chi$).
\end{enumerate}
For real $C_R,T,c$ call $(u,\chi)$ \emph{killed} if some non-zero $\xi\in E^{q_j}$ with
$\Tr_{E/\Q}(\xi\cdot X_u)\subseteq\Z$ satisfies $\sum_{\tau\in\mathcal T}\Re\,\tau(\xi\chi\bar\xi^{t})\le
R(C_R)\,T\,c\,(2fq_j)^2q_j/2$, and put
\[
\mathcal M^{\mathrm{gg}}_w(C_R,T,c)=\bigcap_{(u,\chi)\ \text{killed}}\ker\mathrm{FJ}_{u,\chi},
\qquad Q_w=\mathcal M_w/(\mathcal M_w\cap\mathcal M^{\mathrm{gg}}_w),\qquad \rho_P=\rank_{\OO_K}Q_w ,
\]
the forms whose low Fourier--Jacobi coefficients vanish and the corresponding quotient. The \emph{mean log sup}
of $F\in\mathcal F_w$ is $\ontoonemls(F)=[K:\Q]^{-1}\sum_\sigma\log\|F\|_\sigma$, and for a prime $\mathfrak l$ of $K$ and an
$\OO_K$-module $N$ we write $\ontoonelg_{\mathfrak l}(N)$ for the length over $\OO_{K,\mathfrak l}$ of the localisation
$N_{\mathfrak l}$.
\formal{Formalism.MainNotions.WeightFamilies.Prov}
\formal{Formalism.MainNotions.WeightFamilies.MggK}
\formal{Formalism.MainNotions.WeightFamilies.rhoP}
\formal{Formalism.MainNotions.WeightFamilies.meanLogSup}
\formal{Formalism.MainNotions.WeightFamilies.lenAt}
\end{definition}

\begin{definition}[Valued form algebras]\label{pf:onto1:PJ}
Let $B$ be a form provider and $j$ in range. A \emph{valued form algebra} $X$ over $B$ at $j$ consists of:
\begin{enumerate}[label=(\roman*)]
\item a commutative $K$-algebra $\mathcal A$ and, for every weight $w$, an injective $K$-linear map
$\iota_w\colon\mathcal F_w\to\mathcal A$, the images $\iota_w(\mathcal F_w)$ being independent, such that for
$F\in\mathcal F_w$, $G\in\mathcal F_{w'}$ the product $\iota_w(F)\iota_{w'}(G)$ lies in $\iota_{w+w'}(\mathcal F_{w+w'})$;
\item for every prime $\mathfrak l$ of $K$ a set $\ontoonePt(\mathfrak l)$ of \emph{points} and for each point $y$
a function $v_{\mathfrak l,y}\colon\mathcal A\to[-\infty,\infty]$ with $v_{\mathfrak l,y}(0)=\infty$,
$v_{\mathfrak l,y}\ne-\infty$, $v_{\mathfrak l,y}(F+G)\ge\min(v_{\mathfrak l,y}(F),v_{\mathfrak l,y}(G))$,
$v_{\mathfrak l,y}(aF)=\ontooneord_{\mathfrak l}(a)+v_{\mathfrak l,y}(F)$ for $a\in K^\times$, and
$v_{\mathfrak l,y}(\iota_w(F)\iota_{w'}(G))=v_{\mathfrak l,y}(\iota_wF)+v_{\mathfrak l,y}(\iota_{w'}G)$;
\item the \emph{integrality condition}: $F\in\mathcal M_w$ if and only if $v_{\mathfrak l,y}(\iota_wF)\ge0$ for
all $\mathfrak l$ and all $y\in\ontoonePt(\mathfrak l)$;
\item $\OO_K$-submodules $\mathcal T(w)$ and $\mathcal T^{\mathrm{sh}}(w)$ of $\mathcal F_w$ (the \emph{theta forms}
and the \emph{short theta forms} of weight $w$);
\item chart data relating points to points of a moduli chart, of rank $r$ when $j<m$, from which a real
function $\phi_C(\mathfrak l,y)$ on points is derived.
\end{enumerate}
We write $\iota$ for all the maps $\iota_w$.
\formal{Formalism.MainNotions.WeightAlgebra.PJ}
\formal{Formalism.MainNotions.WeightAlgebra.PJ.phiC}
\end{definition}

\begin{definition}[Theta lattices]\label{pf:onto1:theta}
Let $X$ be a valued form algebra. A \emph{scalar quadruple} is a quadruple $S=(s_W,s_F,s_0,s_K)$ of elements
of $\OO_K$. For $s\in\OO_K$ and a weight $w$ put
\[
W_s(w)=s\cdot\mathcal T(w),\qquad \mathrm{Th}_s(w)=s\,(g_0g_1)^{frn_w}\cdot\mathcal T^{\mathrm{sh}}(w).
\]
For weights $\kappa,w_L,w_F$ and $T\in\N$ the \emph{theta lattice} is the $\OO_K$-submodule of $\mathcal A$
\[
\Theta=\Theta(S,\kappa,w_L,w_F,T)=\iota W_{s_W}(w_L)^{T}\cdot\iota\mathrm{Th}_{s_F}(w_F)
+\iota W_{s_W}(w_L)^{T-\delta}\cdot\iota\mathrm{Th}_{s_0}(w_0)+\iota\mathrm{Th}_{s_K}(\kappa)
\]
(products and powers of $\OO_K$-submodules of $\mathcal A$), and $\ontooneMsup_\kappa=\iota(\mathcal M_\kappa)$. For a prime
$\mathfrak l$ and $y\in\ontoonePt(\mathfrak l)$ the \emph{depth} is $\ontoonedepth_{\mathfrak l,y}=\inf\{v_{\mathfrak
l,y}(\iota F):F\in W_{s_W}(w_L)\}$, and $\ontoonerk_\kappa=\dim_K(K\mathcal M_\kappa)$.
\formal{Formalism.MainNotions.WeightAlgebra.Scal}
\formal{Formalism.MainNotions.WeightAlgebra.PJ.WWc}
\formal{Formalism.MainNotions.WeightAlgebra.PJ.Thc}
\formal{Formalism.MainNotions.WeightAlgebra.PJ.Theta}
\formal{Formalism.MainNotions.WeightAlgebra.PJ.Msup}
\formal{Formalism.MainNotions.WeightAlgebra.PJ.depth}
\end{definition}

\begin{definition}[Local lengths and the index bound]\label{pf:onto1:concl}
Let $\mathfrak l$ be a prime of $K$, $K_{\mathfrak l}$ the completion and $\OO_{\mathfrak l}$ its ring of integers,
and $V_{\mathfrak l}=K_{\mathfrak l}\otimes_K\mathcal A$. For an $\OO_K$-submodule $N\subseteq\mathcal A$ let
$N_{\mathfrak l}$ be the $\OO_{\mathfrak l}$-span and $K_{\mathfrak l}N$ the $K_{\mathfrak l}$-span of its image in
$V_{\mathfrak l}$. For $\OO_{\mathfrak l}$-submodules $A,A'\subseteq V_{\mathfrak l}$ put $\ontooneql_{\mathfrak l}(A,A')=
\len_{\OO_{\mathfrak l}}(A/(A\cap A'))$. Let $\nu$ be a real function on the primes of $K$ and $\beta\in\R$. The
\emph{index bound} $\ontooneIB(\beta)$ for $(X,S,\kappa,w_L,w_F,T,\nu)$ is the statement: there is a function
$r\ge0$ on the primes of $K$ such that $\sum_{\mathfrak l}r(\mathfrak l)\ell(\mathfrak l)$ converges with sum $\le\beta$,
and for every prime $\mathfrak l$, every integer $e\ge T\nu(\mathfrak l)$ and every $K_{\mathfrak l}$-subspace
$V'\subseteq K_{\mathfrak l}\ontooneMsup_\kappa$,
\[
\ontooneql_{\mathfrak l}\bigl(\ontooneMsup_{\kappa,\mathfrak l}\cap V',\ \Theta_{\mathfrak l}\cap V'\bigr)<\infty
\quad\text{and}\quad
\ontooneql_{\mathfrak l}\bigl(\ontooneMsup_{\kappa,\mathfrak l}\cap V',\ \Theta_{\mathfrak l}\cap V'\bigr)\le
e\dim_{K_{\mathfrak l}}V'+r(\mathfrak l).
\]
\formal{Formalism.MainNotions.WeightAlgebra.Concl}
\formal{Formalism.MainNotions.WeightAlgebra.PJ.qlen}
\end{definition}

\begin{remark}\label{pf:onto1:rem-concl}
The index bound does not involve the points, the valuations or the norms: it depends only on $\mathcal A$,
$\iota$, $\mathcal M_\kappa$ and the theta forms. In particular it is unchanged when the set of points is
enlarged or the $L^2$-norm is replaced by another norm (\cref{pf:onto1:transport}).
\end{remark}

%----------------------------------------------------------------------------------------------------
\subsection{Lengths and indices over discrete valuation rings}\label{sec:pf:onto1:dvr}

In this subsection $\OO$ is a discrete valuation ring with fraction field $K$ and uniformiser $\varpi$, and $V$
is a finite dimensional $K$-vector space. A \emph{lattice} in $V$ is a finitely generated $\OO$-submodule
spanning $V$ over $K$. For $\OO$-submodules $A,B\subseteq V$ put
\[
\ontooneql(A,B)=\len_{\OO}\bigl(A/(A\cap B)\bigr)\in\N\cup\{\infty\},\qquad
\ontooneidx(A,B)=\ontooneql(A,B)^\flat-\ontooneql(B,A)^\flat\in\Z .
\]
For lattices both lengths are finite, and $\ontooneidx(A,B)$ is the $\varpi$-adic valuation of the generalised
index of $B$ in $A$. We use the following standard facts, valid for submodules $A,A',B,B',C$ of an
$\OO$-module, each a consequence of the isomorphism theorems and the additivity of length:
\begin{enumerate}[label=(Q\arabic*)]
\item\label{pf:onto1:Q1} $\ontooneql(A,B)\le\ontooneql(A',B)$ if $A\subseteq A'$, and $\ontooneql(A,B')\le\ontooneql(A,B)$
if $B\subseteq B'$;
\item\label{pf:onto1:Q2} $\ontooneql(A+B,B)=\ontooneql(A,B)$, and $\ontooneql(A+C,B)\le\ontooneql(A,B)$ if
$C\subseteq B$;
\item\label{pf:onto1:Q3} $\ontooneql(A,C)\le\ontooneql(A,B)+\ontooneql(B,C)$ (filter $A\supseteq A\cap B\supseteq
A\cap B\cap C$ and note $(A\cap B)/(A\cap B\cap C)\hookrightarrow B/(B\cap C)$);
\item\label{pf:onto1:Q4} if $\phi$ is $\OO$-linear and $\phi(S)\subseteq P$, then $\ontooneql(\phi(A),P)\le
\ontooneql(A,S)$ (the module $A/(A\cap S)$ surjects onto $\phi(A)/(\phi(A)\cap P)$);
\item\label{pf:onto1:Q5} for lattices $A,B,C$ in $V$: $\ontooneidx(A,C)=\ontooneidx(A,B)+\ontooneidx(B,C)$,
$\ontooneidx(A,\varpi^eA)=e\dim_KV$ for $e\in\Z$, and $\ontooneidx(A,B)=-\ontooneql(B,A)\le0$ if $A\subseteq B$; if
$W\subseteq V$ is a subspace then $A\cap W$ is a lattice in $W$ \cite[Ch.~III, \S1]{Serre1979}.
\end{enumerate}

\begin{lemma}[Reduction to a length]\label{pf:onto1:lastred}
Let $M,\Theta$ be lattices in $V$, $e\in\Z$ and $W\subseteq V$ a $K$-subspace. Then
\[
\ontooneidx(M\cap W,\Theta\cap W)\le e\dim_KW+\ontooneql(\varpi^eM+\Theta,\Theta)^\flat .
\]
\formallabel{Formalism.IncomingBounds.index_last_reduction_holds}
\end{lemma}

\begin{proof}
Put $N=\varpi^eM+\Theta$, a lattice. The intersections of $M$, $\varpi^eM$, $N$, $\Theta$ with $W$ are lattices
in $W$ by \ref{pf:onto1:Q5}, and $\varpi^e(M\cap W)=\varpi^eM\cap W\subseteq N\cap W$. By \ref{pf:onto1:Q5},
\[
\ontooneidx(M\cap W,\Theta\cap W)=\ontooneidx(M\cap W,\varpi^e(M\cap W))+\ontooneidx(\varpi^e(M\cap W),N\cap W)
+\ontooneidx(N\cap W,\Theta\cap W),
\]
where the first term is $e\dim W$ and the second is $\le0$. The third is at most $\ontooneql(N\cap W,\Theta\cap
W)\le\ontooneql(N,\Theta)$, because $(N\cap W)/(N\cap\Theta\cap W)$ injects into $N/(N\cap\Theta)$.
\end{proof}

\begin{lemma}[Splitting]\label{pf:onto1:reduction}
Let $X,\Theta,N_0,T',N_K$ be $\OO$-submodules of $V$ with $N_0\subseteq\Theta$ and $T'\subseteq\Theta$. Then
\[
\ontooneql(X+\Theta,\Theta)\le\ontooneql\bigl((X\cap N_K)+N_0,\,N_0\bigr)+\ontooneql\bigl(X,\,(X\cap N_K)+T'\bigr).
\]
\formallabel{Formalism.LatticeBaseChangeAlt.BQuot.reduction}
\end{lemma}

\begin{proof}
By \ref{pf:onto1:Q2} and \ref{pf:onto1:Q3}, $\ontooneql(X+\Theta,\Theta)=\ontooneql(X,\Theta)\le\ontooneql(X,(X\cap
N_K)+T')+\ontooneql((X\cap N_K)+T',\Theta)$. Since $T'\subseteq\Theta$, \ref{pf:onto1:Q2} gives
$\ontooneql((X\cap N_K)+T',\Theta)\le\ontooneql(X\cap N_K,\Theta)$, which is at most $\ontooneql(X\cap N_K,N_0)\le
\ontooneql((X\cap N_K)+N_0,N_0)$ by \ref{pf:onto1:Q1}.
\end{proof}

\begin{lemma}[Levels]\label{pf:onto1:levels}
Let $N',M_1,N_1,Z,Z_m$ be $\OO$-submodules of $V$, let $A,A',N_1'$ be submodules of an $\OO$-module $W'$ and
$\rho\colon N'\to W'$ an $\OO$-linear map such that $M_1\subseteq N'$, $\rho(N')\subseteq A'$, $\ker\rho=Z_m\cap N'$
and $\rho(N_1\cap N')=N_1'$. Then
\[
\ontooneql(M_1,N_1+Z)\le\ontooneql(A',A)+\ontooneql(A,N_1')+\ontooneql(Z_m,Z).
\]
\formallabel{Formalism.LatticeBaseChangeAlt.BQuot.levels}
\end{lemma}

\begin{proof}
Put $Y=\{x\in N':\rho(x)\in A\}$ and $Y'=\{x\in N':\rho(x)\in N_1'\}$. The map $M_1\to A'/(A'\cap A)$,
$x\mapsto\rho(x)$, has kernel $M_1\cap Y$, so $\ontooneql(M_1,Y)\le\ontooneql(A',A)$; likewise the map $Y\to
A/(A\cap N_1')$ has kernel $Y\cap Y'$, so $\ontooneql(Y,Y')\le\ontooneql(A,N_1')$. If $x\in Y'$ then
$\rho(x)=\rho(n)$ for some $n\in N_1\cap N'$, and $x-n\in\ker\rho\subseteq Z_m$; hence $Y'\subseteq N_1+Z_m$. The
map $Z_m\to(N_1+Z_m)/((N_1+Z_m)\cap(N_1+Z))$ is surjective and vanishes on $Z_m\cap Z$, so by \ref{pf:onto1:Q1}
$\ontooneql(Y',N_1+Z)\le\ontooneql(N_1+Z_m,N_1+Z)\le\ontooneql(Z_m,Z)$. Two applications of \ref{pf:onto1:Q3}
conclude.
\end{proof}

\begin{lemma}[Scaling]\label{pf:onto1:scaled}
Let $s\in\OO$, let $M,M^{\mathrm{gg}},N_K$ be $\OO$-submodules of $V$ with $M^{\mathrm{gg}}\subseteq M$ and
$M^{\mathrm{gg}}\subseteq N_K$, and $\theta_1,\dots,\theta_n\in V$. Then
\[
\ontooneql\Bigl(sM,\ (sM\cap N_K)+\textstyle\sum_i\OO s\theta_i\Bigr)\le
\ontooneql\Bigl(M,\ M^{\mathrm{gg}}+\textstyle\sum_i\OO\theta_i\Bigr).
\]
\formallabel{Formalism.MainNotions.Gaps.T0.ql_scaled_le}
\end{lemma}

\begin{proof}
Apply \ref{pf:onto1:Q4} to multiplication by $s$ and $S=M^{\mathrm{gg}}+\sum\OO\theta_i$: indeed
$sM^{\mathrm{gg}}\subseteq sM\cap N_K$.
\end{proof}

\begin{proposition}[The local step]\label{pf:onto1:localStepQ}
Let $W'$ be an $\OO$-module. Let $M,\Theta$ be lattices in $V$, let $N_1,Z,Z_m,N',M^{\mathrm{gg}}$ be $\OO$-submodules
of $V$, $\theta_1,\dots,\theta_n\in V$, $e,q\in\N$, let $A,A',N_1'$ be submodules of $W'$ and $\rho\colon N'\to W'$
$\OO$-linear. Assume:
\begin{enumerate}[label=(\alph*)]
\item $N_1+Z\subseteq\Theta$, and $\theta_i\in M\cap\Theta$ for all $i$;
\item $M^{\mathrm{gg}}\subseteq M$, $KM^{\mathrm{gg}}\subseteq K(N_1+Z)$, and $\ontooneql\bigl(M,M^{\mathrm{gg}}+\sum_i\OO\theta_i\bigr)\le q$;
\item $\varpi^eM\cap K(N_1+Z)\subseteq N'$, $N_1\subseteq N'$ and $Z\subseteq Z_m\subseteq N'$;
\item $\rho(N')\subseteq A'$, $\ker\rho=Z_m$, $\rho(N_1)=N_1'$, and $N_1'\subseteq A\subseteq A'$;
\item the lengths $\ontooneql(A',A)$, $\ontooneql(A,N_1')$ and $\ontooneql(Z_m,Z)$ are finite.
\end{enumerate}
Then for every $K$-subspace $V'\subseteq V$,
\[
\ontooneidx(M\cap V',\Theta\cap V')\le e\dim_KV'+\ontooneql(A',A)+\ontooneql(A,N_1')+\ontooneql(Z_m,Z)+q.
\]
\formallabel{Formalism.MainNotions.Gaps.T0.localStepQ}
\end{proposition}

\begin{proof}
Put $X=\varpi^eM$, $N_K=K(N_1+Z)$ (an $\OO$-module) and $T'=\sum_i\OO\varpi^e\theta_i\subseteq\Theta$. By
\cref{pf:onto1:lastred}, $\ontooneidx(M\cap V',\Theta\cap V')\le e\dim V'+\ontooneql(X+\Theta,\Theta)$. By
\cref{pf:onto1:reduction} with $N_0=N_1+Z\subseteq\Theta$,
\[
\ontooneql(X+\Theta,\Theta)\le\ontooneql\bigl((X\cap N_K)+N_1+Z,\ N_1+Z\bigr)+\ontooneql\bigl(X,(X\cap N_K)+T'\bigr).
\]
The module $M_1=(X\cap N_K)+N_1+Z$ is contained in $N'$ by (c), so \cref{pf:onto1:levels} bounds the first term by
$\ontooneql(A',A)+\ontooneql(A,N_1')+\ontooneql(Z_m,Z)$. Since $M^{\mathrm{gg}}\subseteq KM^{\mathrm{gg}}\subseteq N_K$,
\cref{pf:onto1:scaled} with $s=\varpi^e$ bounds the second term by $\ontooneql(M,M^{\mathrm{gg}}+\sum\OO\theta_i)\le q$.
All terms being finite, the inequality of lengths gives the asserted inequality of integers.
\end{proof}

\begin{remark}\label{pf:onto1:rem-localstep}
The hypotheses $\theta_i\in M$ in (a) and $N_1'\subseteq A\subseteq A'$ in (d) are not used in the proof. When
$\theta_i\in M$ for all $i$, the length in (b) is $\len_{\OO}\bigl(M/(M^{\mathrm{gg}}+\sum_i\OO\theta_i)\bigr)$.
\end{remark}

\begin{lemma}[Pole orders]\label{pf:onto1:levelsChain}
Let $Z_0$ be a finite set, and for $z\in Z_0$ let $e_z\ge1$ and $f_z\ge0$ be integers and $\nu_z$, $v_z\ge0$ real
numbers. Let $T\ge1$ and $e$ be integers and $\nu_s\in\R$, $c'\ge0$ with $\nu_z\le\nu_s+c'v_z$ for all $z$.
Assume that either (a) $T\nu_s\le e$ and $e_z\nu_z\in\Z$ for all $z$, or (b) $T\nu_s+1\le e$. Then
\[
\sum_{z\in Z_0}\max\bigl(0,\lceil e_z(T\nu_z-e)\rceil\bigr)f_z\le Tc'\sum_{z\in Z_0}e_zf_zv_z.
\]
\formallabel{Formalism.LatticeBaseChangeAlt.levelsChain}
\formallabel{Formalism.LatticeBaseChangeAlt.levelsChain_of_add_one_le}
\end{lemma}

\begin{proof}
It suffices to show $\max(0,\lceil e_z(T\nu_z-e)\rceil)\le e_zTc'v_z$ for each $z$; the right side is $\ge0$. In
case (a), $e_z(T\nu_z-e)=T(e_z\nu_z)-e_ze$ is an integer, hence equal to its ceiling, and it is at most
$e_z(T\nu_s-e)+e_zTc'v_z\le e_zTc'v_z$. In case (b), $T\nu_z-e\le Tc'v_z-1$, so
$\lceil e_z(T\nu_z-e)\rceil<e_z(T\nu_z-e)+1\le e_zTc'v_z-e_z+1\le e_zTc'v_z$.
\end{proof}

\begin{remark}\label{pf:onto1:rem-levelschain}
Of the two formal statements tagged in \cref{pf:onto1:levelsChain}, the first treats case (a) together with the
special case $e=\lceil T(\nu_s+1)\rceil$ of (b), and the second treats case (b).
\end{remark}

\begin{lemma}[Base change of indices]\label{pf:onto1:descent}
Let $\OO\to\OO'$ be an injective local homomorphism of discrete valuation rings such that the image of $\varpi$ is
a uniformiser of $\OO'$, with compatible fraction fields $K\to K'$, let $j\colon V\to V_1$ be a $K$-linear map
exhibiting $V_1$ as $K'\otimes_KV$, let $M,\Theta$ be lattices in $V$ and $V'\subseteq V$ a subspace, and write
$M'=\OO'j(M)$, $\Theta'=\OO'j(\Theta)$, $V''=K'j(V')$. Then $\dim_{K'}V''=\dim_KV'$ and
\[
\ontooneidx_{\OO'}(M'\cap V'',\Theta'\cap V'')=\ontooneidx_{\OO}(M\cap V',\Theta\cap V').
\]
Moreover, for an $\OO$-submodule $P\subseteq M$, $\len_{\OO'}\bigl(M'/\OO'j(P)\bigr)\le\len_{\OO}(M/P)$.
Consequently, if $d\cdot\ontooneidx_{\OO'}(M'\cap V'',\Theta'\cap V'')\le d(e\dim V''+Y)+Z$ for real $d,Y,Z$ and
$e\in\N$, then $d\cdot\ontooneidx_{\OO}(M\cap V',\Theta\cap V')\le d(e\dim V'+Y)+Z$.
\formallabel{Formalism.LatticeBaseChangeAlt.localStepDescent}
\end{lemma}

\begin{proof}
Since $j$ is a base change, it maps $K$-linearly independent families to $K'$-linearly independent families,
which gives the equality of dimensions. As $\OO$ is a principal ideal domain and $M/(M\cap V')$ is torsion
free, there is an $\OO$-basis $b_1,\dots,b_n$ of $M$ such that $b_1,\dots,b_s$ is a basis of $M\cap V'$ and a
$K$-basis of $V'$. Then $j(b_1),\dots,j(b_n)$ is an $\OO'$-basis of $M'$ and a $K'$-basis of $V_1$, and $V''$ is
spanned by $j(b_1),\dots,j(b_s)$; hence $M'\cap V''=\OO'j(M\cap V')$, and likewise $\Theta'\cap V''=\OO'j(\Theta\cap
V')$. Now let $A,B$ be lattices in a $K$-space $U$ with bases $a,b$ related by $b=aP$, $P\in\GL(K)$. By the theory
of elementary divisors, $\ontooneidx_{\OO}(A,B)=\ontooneord_\varpi\det P$. The bases $j(a),j(b)$ of $\OO'j(A)$ and
$\OO'j(B)$ are related by the same matrix, and the normalised valuation of $K'$ restricts to $\ontooneord_\varpi$ on
$K$ because $\varpi$ is a uniformiser of $\OO'$; this gives the equality of indices. For the length: a
composition series of $M/P$ has factors $\OO/\varpi$; tensoring with $\OO'$ (a right exact operation) gives a
filtration of $\OO'\otimes_{\OO}(M/P)$, which surjects onto $M'/\OO'j(P)$, with factors that are quotients of
$\OO'/\varpi\OO'$, the residue field of $\OO'$. The last assertion follows.
\end{proof}

\begin{proposition}[Levels and descent]\label{pf:onto1:levelsQDesc}
Let $\OO\to\OO'$, $K\to K'$, $j\colon V\to V_1$ and $\varpi$ be as in \cref{pf:onto1:descent}, and let $W'$ be an
$\OO'$-module. Let $Z_0$ be a finite set with integers $e_z\ge1$, $f_z\ge0$ and real numbers $\nu_z$, $v_z\ge0$
($z\in Z_0$), and let $T\ge1$, $R,I,e,n,q\in\N$, $\nu_s\in\R$ and $c'\ge0$ with $\nu_z\le\nu_s+c'v_z$. Let $M,\Theta$
be lattices in $V$, $M^{\mathrm{gg}}\subseteq V$ an $\OO$-submodule and $\theta_1,\dots,\theta_n\in V$; let
$N_1,Z,Z_m,N'$ be $\OO'$-submodules of $V_1$, $A,A',N_1'$ submodules of $W'$, and $\rho\colon N'\to W'$
$\OO'$-linear. Assume:
\begin{enumerate}[label=(\alph*)]
\item $N_1+Z\subseteq\OO'j(\Theta)$; $\theta_i\in M\cap\Theta$; $M^{\mathrm{gg}}\subseteq M$;
$j(M^{\mathrm{gg}})\subseteq K'(N_1+Z)$; $\len_{\OO}\bigl(M/(M^{\mathrm{gg}}+\sum\OO\theta_i)\bigr)\le q$;
\item $\varpi^e\OO'j(M)\cap K'(N_1+Z)\subseteq N'$, $N_1\subseteq N'$ and $Z\subseteq Z_m\subseteq N'$;
\item $\rho(N')\subseteq A'$, $\ker\rho=Z_m$, $\rho(N_1)=N_1'\subseteq A\subseteq A'$, and $\len(Z_m/Z)<\infty$;
\item either $T\nu_s\le e$ and $e_z\nu_z\in\Z$ for all $z$, or $T\nu_s+1\le e$;
\item $\len_{\OO'}(A'/A)\le R\sum_z\max(0,\lceil e_z(T\nu_z-e)\rceil)f_z$ and $\len_{\OO'}(A/N_1')\le RI$.
\end{enumerate}
Then for every $K$-subspace $V'\subseteq V$,
\[
\ontooneidx(M\cap V',\Theta\cap V')\le e\dim_KV'+RTc'\sum_{z}e_zf_zv_z+RI+\len_{\OO'}(Z_m/Z)+q.
\]
\formallabel{Formalism.MainNotions.Gaps.EtaleWeightGaps.levelsQDesc}
\end{proposition}

\begin{proof}
By \cref{pf:onto1:descent}, $\OO'j(M)$ and $\OO'j(\Theta)$ are lattices in $V_1$ (spanned by images of bases),
$\len_{\OO'}(\OO'j(M)/(\OO'j(M^{\mathrm{gg}})+\sum\OO'j(\theta_i)))\le q$, and $K'\OO'j(M^{\mathrm{gg}})\subseteq
K'(N_1+Z)$. Hence \cref{pf:onto1:localStepQ} applies over $\OO'$ to $\OO'j(M)$, $\OO'j(\Theta)$, $\OO'j(M^{\mathrm{gg}})$,
$j(\theta_i)$ and the subspace $V''=K'j(V')$ (the lengths in (e) are finite), and gives
$\ontooneidx(\OO'j(M)\cap V'',\OO'j(\Theta)\cap V'')\le e\dim V''+\len(A'/A)+\len(A/N_1')+\len(Z_m/Z)+q$. By (e)
and \cref{pf:onto1:levelsChain}, $\len(A'/A)\le RTc'\sum_ze_zf_zv_z$, and $\len(A/N_1')\le RI$. Finally
\cref{pf:onto1:descent} (with $d=1$) transfers the bound to $\ontooneidx(M\cap V',\Theta\cap V')$.
\end{proof}

%----------------------------------------------------------------------------------------------------
\subsection{The generic linear section at a prime and the local bound}\label{sec:pf:onto1:local}

Throughout this subsection we fix a form provider $B$ with field $K$, a valued form algebra $X$ at $j$, weights
$w_L,w_F$, an integer $T\ge1$ with $\delta\le T$, the weight $\kappa=w_F+Tw_L$, the integer $a_0=T-\delta$ and a
scalar quadruple $S$. We abbreviate $W=W_{s_W}(w_L)$, $\mathrm{Th}_F=\mathrm{Th}_{s_F}(w_F)$,
$\mathrm{Th}_0=\mathrm{Th}_{s_0}(w_0)$ and $\mathrm{Th}_K=\mathrm{Th}_{s_K}(\kappa)$.

\begin{definition}[Global family data]\label{pf:onto1:gfd}
A \emph{global family datum} $G$ for $(X,\kappa,w_L,w_F)$ consists of: forms $u_0,\dots,u_{N-1}\in\mathcal F_{w_L}$;
forms $\theta'_0,\dots,\theta'_{\Delta-1}\in\mathcal F_{w_F}$; for $k<r_z$ a multi-index $\alpha_k\in\N^{\delta}$, a form
$\theta^0_k\in\mathcal F_{w_0}$ and a finitely supported map $z_k$ from the monomials in variables $X_{i,k'}$
($0\le i\le\delta$, $k'<N$) to $\mathcal F_\kappa$; integral forms $\theta_0,\dots,\theta_{r_q-1}\in\mathcal M_\kappa$; an integer
$i_\eta\ge0$ and a form $\ontooneet\in\mathcal F_{\epsilon_{i_\eta}}$ (the \emph{auxiliary form}); and a real function $m_x$
on the embeddings $K\to\C$. For $\alpha\in\N^\delta$ and $\theta\in\mathcal F_{w}$ put
\[
z(\alpha,\theta)=\prod_{i=1}^{\delta}\Bigl(\sum_{k'<N}X_{i,k'}\,\iota(u_{k'})\Bigr)^{\alpha_i}\cdot\iota(\theta)\in\mathcal A[X].
\]
For $K$-linear forms $\phi_1,\dots,\phi_{r_z}$ on $\mathcal F_\kappa$ let $\det(\phi,z)\in K[X]$ be the determinant of the
matrix whose $(a,b)$ entry is the polynomial $\sum_\beta\phi_a(z_b(\beta))X^\beta$. For a prime $\mathfrak l$ the
\emph{Gauss content} $\gamma_{\mathfrak l}$ is the least $n\in\N$ such that there are $\phi_1,\dots,\phi_{r_z}$ with
$\ontooneord_{\mathfrak l}\phi_a(\mathcal M_\kappa)\ge0$ for all $a$ and a monomial $\beta$ whose coefficient in $\det(\phi,z)$ is
non-zero of $\mathfrak l$-adic valuation $n$ (and $\gamma_{\mathfrak l}=0$ if there is no such $n$).
\formal{Formalism.MainNotions.Gaps.EtaleWeightGaps.GlobalFamilyData}
\formal{Formalism.MainNotions.Gaps.SizeBoundsAlt.zPoly}
\formal{Formalism.MainNotions.Gaps.SizeBoundsAlt.detPoly}
\formal{Formalism.MainNotions.Gaps.SizeBoundsAlt.gaussC}
\end{definition}

\begin{definition}[The generic linear section]\label{pf:onto1:gauss}
Fix $G$ and a prime $\mathfrak l$ of $K$, and let $\varpi$ be a uniformiser of $\OO_{\mathfrak l}$.
\begin{enumerate}[label=(\roman*)]
\item The \emph{Gauss extension} $\OO'$ is the localisation of the polynomial ring $\OO_{\mathfrak
l}[X_{i,k}:0\le i\le\delta,\,k<N]$ at the ideal $\varpi\OO_{\mathfrak l}[X]$ of polynomials with all coefficients in
the maximal ideal, and $K'$ is its fraction field. Since $\OO_{\mathfrak l}[X]/\varpi\OO_{\mathfrak l}[X]$ is a
polynomial ring over a field, this ideal is prime; $\OO'$ is a noetherian local domain whose maximal ideal is
generated by $\varpi$, hence a discrete valuation ring with uniformiser $\varpi$, and it is flat over
$\OO_{\mathfrak l}$.
\item Let $R'=K'\otimes_{K_{\mathfrak l}}V_{\mathfrak l}=K'\otimes_K\mathcal A$ with the natural map
$\mathrm{emb}\colon\mathcal A\to R'$, and for a weight $u$ let $V'_u\subseteq R'$ be the $K'$-span of $\mathrm{emb}\,\iota(\mathcal
F_u)$; then $V'_uV'_{u'}\subseteq V'_{u+u'}$. The \emph{generic sections} are
\[
s_i=\sum_{k<N}X_{i,k}\,\mathrm{emb}\,\iota(u_k)\in V'_{w_L}\quad(0\le i\le\delta),\qquad
s^\alpha=\prod_{i=1}^{\delta}s_i^{\alpha_i}\quad(\alpha\in\N^\delta).
\]
\item Let $L'=R'[s_0^{-1}]$, and for a weight $u$ let $H(u)\subseteq L'$ be the $K'$-span of the fractions $g/s_0^k$
with $k\in\N$ and $g\in V'_{u+kw_L}$; then $H(u)H(u')\subseteq H(u+u')$, and the $K'$-algebra $\mathrm{Fr}$ generated by
$H(0)$ equals $H(0)$. Let $J\subseteq\mathrm{Fr}$ be the ideal generated by $s_1/s_0,\dots,s_\delta/s_0$, and put
\[
\OO_Z=\mathrm{Fr}/J,\qquad \mathrm{WK}(u,a)=H(u)/J^aH(u),\qquad \mathrm{SQ}(u)=H(u)/JH(u).
\]
Thus $\OO_Z$ is the ring of functions of the subscheme $s_1=\dots=s_\delta=0$ of the chart $s_0\ne0$,
$\mathrm{SQ}(u)$ is the module of restrictions to it of sections of weight $u$, and $\mathrm{WK}(u,a)$ is the module of
sections modulo vanishing to order $a$ along it. We write $\mathrm{res}\colon V'_u\to\mathrm{SQ}(u)$ for the
restriction and $\rho_T\colon V'_{u+Tw_L}\to\mathrm{WK}(u,a)$, $g\mapsto g/s_0^T\bmod J^aH(u)$. The Hasse derivatives
$\partial^\beta$ with respect to the variables $X_{i,k}$ act $K_{\mathfrak l}$-linearly on $L'$.
\item Let $\bar\OO$ be the integral closure of $\OO'$ in $\OO_Z$ and $\mathrm{Pts}$ the set of maximal ideals of
$\bar\OO$ (the \emph{points} of the section). For $\zeta\in\mathrm{Pts}$ such that $\bar\OO_\zeta$ is a discrete
valuation ring, $e_\zeta$ is the valuation of $\varpi$ and $f_\zeta$ the residue degree over $\OO'$; for
$x\in\OO_Z$, $\ontooneord_\zeta(x)=\ontooneord_\zeta(\varpi^nx)-ne_\zeta$ if $\varpi^nx\in\bar\OO$ for some $n$, and
$\ontooneord_\zeta(x)=\infty$ otherwise.
\end{enumerate}
\formal{Formalism.MainNotions.Gaps.EtaleWeightGaps.Op}
\formal{Formalism.MainNotions.Gaps.EtaleWeightGaps.Vw}
\formal{Formalism.MainNotions.Gaps.EtaleWeightGaps.sec}
\formal{Formalism.MainNotions.Gaps.EtaleWeightGaps.OsK}
\formal{Formalism.MainNotions.Gaps.EtaleWeightGaps.WK}
\formal{Formalism.MainNotions.Gaps.EtaleWeightGaps.SectionSheafQuotient}
\formal{Formalism.MainNotions.Gaps.EtaleWeightGaps.rhoK}
\formal{Formalism.MainNotions.Gaps.EtaleWeightGaps.hasse}
\formal{Formalism.MainNotions.Gaps.EtaleWeightGaps.ep}
\formal{Formalism.MainNotions.Gaps.EtaleWeightGaps.fdeg}
\formal{Formalism.DV.GaussP.isDVR}
\end{definition}

\begin{lemma}[Integral closures]\label{pf:onto1:intclos}
Let $\OO$ be a discrete valuation ring whose fraction field $K$ has characteristic $0$, and $A$ a finite
dimensional reduced $K$-algebra. Then the integral closure of $\OO$ in $A$ is a finite $\OO$-module and its
localisation at every maximal ideal is a discrete valuation ring.
\formallabel{Formalism.MainNotions.Gaps.EtaleWeightGaps.integralClosure_dvr'}
\end{lemma}

\begin{proof}
$A$ is artinian and reduced, hence a finite product of finite field extensions $A/\mathfrak m$ of $K$, separable
since $K$ has characteristic $0$; the idempotents separating the factors are integral over $\OO$. Hence the
integral closure of $\OO$ in $A$ is the product of the integral closures of $\OO$ in the factors, each a finite
$\OO$-module and a Dedekind domain which is not a field \cite[Ch.~I, \S4]{Serre1979}; the maximal ideals of the
product are those of the factors, and the localisations are discrete valuation rings.
\end{proof}

\begin{definition}[Local point data]\label{pf:onto1:lpd}
A \emph{local point datum} $Q$ for $G$ at $\mathfrak l$ consists of points $y_\zeta\in\ontoonePt(\mathfrak l)$ for
$\zeta\in\mathrm{Pts}$, a form $g_\circ\in\mathcal F_{w_F}$ and integers $c(\zeta)$. Put
$\varepsilon=\mathrm{res}\,\mathrm{emb}\,\iota(g_\circ)\in\mathrm{SQ}(w_F)$; when $\varepsilon$ generates $\mathrm{SQ}(w_F)$ freely over
$\OO_Z$, write $v=\mathrm{co}(v)\varepsilon$ and $\ontooneord^Q_\zeta(v)=\ontooneord_\zeta(\mathrm{co}(v))+c(\zeta)$. Define:
\begin{itemize}
\item the \emph{ball} $A_0\subseteq\mathrm{SQ}(w_F)$, the $\bar\OO$-span of $\{v:\ontooneord^Q_\zeta(v)\ge0\ \forall\zeta\}$; the
restrictions $t_i=\mathrm{res}\,\mathrm{emb}\,\iota(\theta'_i)$; and the \emph{local height}
$I_{\mathfrak l}=\len_{\OO'}\bigl(A_0/(A_0\cap\sum_i\OO't_i)\bigr)$;
\item the depths $\nu(\zeta)=\ontoonedepth_{\mathfrak l,y_\zeta}$; the sum
$\mathrm{Net}_{\mathfrak l}=\sum_\zeta e_\zeta f_\zeta\,v_{\mathfrak l,y_\zeta}(\iota\ontooneet)$; and
$v_\eta(\zeta)=v_{\mathfrak l,y_\zeta}(\iota\ontooneet)/(2i_\eta)$ if $j<m$, $v_\eta(\zeta)=0$ if $j\ge m$;
\item for $m\colon\mathrm{Pts}\to\N$, the $\OO'$-submodule $\mathrm{AA}^+(m)\subseteq\mathrm{WK}(w_F,a_0)$ spanned by the images of
those $x\in H(w_F)$ such that for every multi-index $\beta$ of degree $\le a_0-1$ the derivative $\partial^\beta x$ lies
in $H(w_F)$ and its restriction $v$ to $\mathrm{SQ}(w_F)$ satisfies $\ontooneord^Q_\zeta(v)\ge-m(\zeta)$ for all $\zeta$; and
$\mathrm{AA}=\mathrm{AA}^+(0)$;
\item for $e\in\N$, the pole orders $m_e(\zeta)=\max\bigl(0,\lceil e_\zeta(T\nu(\zeta)-e)\rceil\bigr)$.
\end{itemize}
\formal{Formalism.MainNotions.Gaps.EtaleWeightGaps.LocalPointData}
\formal{Formalism.MainNotions.Gaps.EtaleWeightGaps.LocalPointData.Ila}
\formal{Formalism.MainNotions.Gaps.EtaleWeightGaps.LocalPointData.Net}
\formal{Formalism.MainNotions.Gaps.EtaleWeightGaps.vEt}
\formal{Formalism.MainNotions.Gaps.EtaleWeightGaps.LocalPointData.AAp}
\formal{Formalism.MainNotions.Gaps.EtaleWeightGaps.LocalPointData.mm}
\end{definition}

\begin{definition}[Local conditions]\label{pf:onto1:localconds}
With $G,\mathfrak l,Q$ as above, consider the following conditions.
\begin{itemize}
\item[$(\mathrm B_{\mathfrak l})$] (basis) $t_0,\dots,t_{\Delta-1}$ is a $K'$-basis of $\mathrm{SQ}(w_F)$.
\item[$(\mathrm C_{\mathfrak l})$] (finiteness) $\OO_Z$ is a finite dimensional reduced $K'$-algebra and $\varepsilon$ generates
$\mathrm{SQ}(w_F)$ freely over $\OO_Z$.
\item[$(\mathrm T_{\mathfrak l})$] (ties) For every $\zeta$: $\nu(\zeta)\ne\pm\infty$; $v_{\mathfrak l,y_\zeta}(\iota\ontooneet)\ne\infty$ if
$j<m$; $e_\zeta\nu(\zeta)\in\Z$; $c(\zeta)=e_\zeta v_{\mathfrak l,y_\zeta}(\iota g_\circ)$; and for all $k\in\N$ and
$g\in\mathcal F_{kw_L}$ (resp.\ $g\in\mathcal F_{w_F+kw_L}$) the class $x$ of $\mathrm{emb}\,\iota(g)/s_0^k$ in $\OO_Z$ (resp.\ in
$\mathrm{SQ}(w_F)$) satisfies $\ontooneord_\zeta(x)=e_\zeta(v_{\mathfrak l,y_\zeta}(\iota g)-k\nu(\zeta))$ (resp.\
$\ontooneord^Q_\zeta(x)=e_\zeta(v_{\mathfrak l,y_\zeta}(\iota g)-k\nu(\zeta))$).
\item[$(\mathrm E_{\mathfrak l})$] (lifting) Every $a\in A_0$ is the restriction of some $x\in H(w_F)$ whose image in
$\mathrm{WK}(w_F,a_0)$ lies in $\mathrm{AA}$.
\item[$(\mathrm Z_{\mathfrak l})$] (zero containment) The $K'$-span of $\{s^\alpha\mathrm{emb}\,\iota\theta:|\alpha|=a_0,\
\theta\in\mathrm{Th}_0\}$ is contained in the $K'$-span of $\{s^{\alpha_k}\mathrm{emb}\,\iota\theta^0_k\}$.
\item[$(\mathrm M_{\mathfrak l})$] (Fourier--Jacobi containment, with parameters $C_R,c$) Write
$\mathcal M^{\mathrm{gg}}=\mathcal M^{\mathrm{gg}}_\kappa(C_R,T,c)$. The $K'$-span of the set
$\mathrm{emb}\,\iota(\mathcal M^{\mathrm{gg}})$ is contained in the $K'$-span of the elements
$s^\alpha s_0^{T-|\alpha|}\mathrm{emb}\,\iota\theta'_i$ with $|\alpha|<a_0$, $i<\Delta$, and $s^{\alpha_k}\mathrm{emb}\,\iota\theta^0_k$
with $k<r_z$.
\item[$(\mathrm M'_{\mathfrak l})$] (weak containment) The same, with the set $\{s^{\alpha_k}\mathrm{emb}\,\iota\theta^0_k\}$
replaced by the set of all elements $s^\alpha\mathrm{emb}\,\iota\theta$ with $|\alpha|=a_0$ and $\theta\in\mathrm{Th}_0$.
\end{itemize}
\formal{Formalism.MainNotions.Gaps.EtaleWeightGaps.InBasis}
\formal{Formalism.MainNotions.Gaps.EtaleWeightGaps.LocalPointData.InCurve}
\formal{Formalism.MainNotions.Gaps.EtaleWeightGaps.LocalPointData.InTie}
\formal{Formalism.MainNotions.Gaps.EtaleWeightGaps.LocalPointData.InLift}
\formal{Formalism.MainNotions.Gaps.EtaleWeightGaps.InContainZ}
\formal{Formalism.MainNotions.Gaps.EtaleWeightGaps.InContainM}
\formal{Formalism.MainNotions.Gaps.EtaleWeightGaps.InContainMR}
\end{definition}

\begin{definition}[Local defect and local modules]\label{pf:onto1:rLa}
For $C_R,c\in\R$ let $Q_\kappa=Q_\kappa(C_R,T,c)$ and $\lambda_{\mathfrak l}=\ontoonelg_{\mathfrak l}\bigl(Q_\kappa/\sum_i\OO_K\bar
\theta_i\bigr)^\flat$, where $\bar\theta_i$ is the class of $\theta_i$. The \emph{local defect} is
\[
r(\mathfrak l)=\ontoonebin\Bigl(I_{\mathfrak l}^\flat+[j<m]\,T\,\frac{\max(0,c)}{2i_\eta}\,\mathrm{Net}_{\mathfrak l}\Bigr)
+\gamma_{\mathfrak l}+\lambda_{\mathfrak l},
\]
where $[j<m]$ is $1$ if $j<m$ and $0$ otherwise. The \emph{local modules} are the $\OO'$-submodules of $V'_\kappa$
\begin{gather*}
M_1=\OO'\,\mathrm{emb}\,\iota(\mathcal M_\kappa),\qquad
N_1=\textstyle\sum_{|\alpha|<a_0,\,i<\Delta}\OO'\,s^\alpha s_0^{T-|\alpha|}\mathrm{emb}\,\iota\theta'_i,\\
Z_l=\textstyle\sum_{k<r_z}\OO'\,s^{\alpha_k}\mathrm{emb}\,\iota\theta^0_k,\qquad Z_m=M_1\cap K'Z_l,\\
N'_e=\bigl(\varpi^eM_1\cap K'(N_1+Z_l)\bigr)+N_1+Z_m\quad(e\in\N),
\end{gather*}
together with the map $\rho_e=\rho_T|_{N'_e}\colon N'_e\to\mathrm{WK}(w_F,a_0)$ and the module $N_1'=\rho_T(N_1)$.
\formal{Formalism.MainNotions.Gaps.EtaleWeightGaps.rLa}
\formal{Formalism.MainNotions.Gaps.EtaleWeightGaps.N1}
\formal{Formalism.MainNotions.Gaps.EtaleWeightGaps.Zl}
\formal{Formalism.MainNotions.Gaps.EtaleWeightGaps.Zm}
\formal{Formalism.MainNotions.Gaps.EtaleWeightGaps.Np}
\formal{Formalism.MainNotions.Gaps.EtaleWeightGaps.rhoN}
\formal{Formalism.MainNotions.Gaps.EtaleWeightGaps.firstLength_le_prime_bound}
\end{definition}

\begin{lemma}[Independence of monomials]\label{pf:onto1:ind_claim}
Let $K'$ be a field of characteristic $0$, $L'$ a commutative $K'$-algebra, $y_1,\dots,y_\delta\in L'$, $t'_1,\dots,t'_\Delta\in
L'$, $(J_a)_{a\in\N}$ a decreasing sequence of $K'$-subspaces of $L'$, and $D_1,\dots,D_\delta$ derivations of $L'$ (over a
subring) such that $D_l(y_{l'})=\delta_{ll'}$, $D_l(t'_i)=0$, $D_l(K')\subseteq K'$, $D_l(J_{a+1})\subseteq J_a$, and
$y^\mu t'_i\in J_{|\mu|}$ for every monomial $y^\mu$ and every $i$. Assume that $\sum_ib_it'_i\in J_1$ with $b_i\in K'$ implies
$b=0$. Then for every $a\in\N$ and all polynomials $c_1,\dots,c_\Delta\in K'[Y_1,\dots,Y_\delta]$ of degree $<a$ with
$\sum_ic_i(y)t'_i\in J_a$, all $c_i$ vanish.
\formallabel{Formalism.MainNotions.Gaps.EtaleWeightGaps.ind_claim}
\end{lemma}

\begin{proof}
Induction on $a$. For $a=0$ there is no non-zero polynomial of degree $<0$. Let $a\ge0$ and let $c_i$ have degree $\le a$
with $\sum_ic_i(y)t'_i\in J_{a+1}$. If $a=0$ the $c_i$ are constants and the hypothesis gives $c=0$. Let $a\ge1$ and write
$c_i=c_i^{<}+c_i^{(a)}$ with $c_i^{(a)}$ homogeneous of degree $a$. Then $\sum_ic_i^{(a)}(y)t'_i\in J_a$ by the monomial
condition, and $\sum_ic_i(y)t'_i\in J_{a+1}\subseteq J_a$, so $\sum_ic_i^{<}(y)t'_i\in J_a$ and the induction hypothesis
gives $c_i^<=0$: the $c_i$ are homogeneous of degree $a$. For a monomial $y^\mu$ with coefficient $\gamma\in K'$ one has
$D_l(\gamma y^\mu t'_i)=D_l(\gamma)y^\mu t'_i+\gamma(\partial_l Y^\mu)(y)t'_i$, because $D_l$ is a derivation with $D_l(y_{l'})
=\delta_{ll'}$ and $D_l(t'_i)=0$; the first term lies in $J_a$ since $D_l(\gamma)\in K'$. Hence
$D_l(\sum_ic_i(y)t'_i)-\sum_i(\partial_lc_i)(y)t'_i\in J_a$, and since $D_l$ maps $J_{a+1}$ into $J_a$,
$\sum_i(\partial_lc_i)(y)t'_i\in J_a$ with $\partial_lc_i$ of degree $a-1<a$. By induction $\partial_lc_i=0$ for all $l$, and
Euler's identity $\sum_lY_l\partial_lc_i=a\,c_i$ with $a\ne0$ in $K'$ gives $c_i=0$.
\end{proof}

\begin{lemma}[Injectivity on the section monomials]\label{pf:onto1:in64}
Assume $(\mathrm B_{\mathfrak l})$. If $x$ lies in the $K'$-span of the elements $s^\alpha s_0^{T-|\alpha|}\mathrm{emb}\,\iota\theta'_i$
($|\alpha|<a_0$, $i<\Delta$) and $\rho_T(x)=0$, then $x=0$.
\formallabel{Formalism.MainNotions.Gaps.EtaleWeightGaps.in64}
\end{lemma}

\begin{proof}
Put $y_l=s_l/s_0\in L'$ ($1\le l\le\delta$) and $t'_i=\mathrm{emb}\,\iota\theta'_i\in L'$, and let $J_a$ be the $K'$-span of the
products $y^\mu h$ with $|\mu|=a$ and $h\in H(w_F)$, so that $J_a=J^aH(w_F)$ and $\rho_T(g)=0$ iff $g/s_0^T\in J_{a_0}$. Since
$s^\alpha s_0^{T-|\alpha|}=s_0^Ty^\alpha$, we can write $x=s_0^T\sum_ic_i(y)t'_i$ with polynomials $c_i$ of degree $<a_0$, and
$\rho_T(x)=0$ means $\sum_ic_i(y)t'_i\in J_{a_0}$. Let $D_l=\sum_{k<N}X_{0,k}\,\partial/\partial X_{l,k}$, a derivation of $L'$. Since
$s_0$ does not involve the variables $X_{l,k}$ ($l\ge1$) and $\partial s_{l'}/\partial X_{l,k}=\delta_{ll'}\mathrm{emb}\,\iota(u_k)$, one
has $D_l(y_{l'})=\delta_{ll'}s_0/s_0=\delta_{ll'}$; $D_l$ kills $\mathrm{emb}(\mathcal A)$, maps $K'$ into $K'$, and maps $J_{a+1}$
into $J_a$; and $y^\mu t'_i\in J_{|\mu|}$. Finally, $\sum_ib_it'_i\in J_1$ means that $\sum_ib_it_i=0$ in $\mathrm{SQ}(w_F)$,
which forces $b=0$ by $(\mathrm B_{\mathfrak l})$. \Cref{pf:onto1:ind_claim} (with $K'$ of characteristic $0$) gives $c_i=0$,
hence $x=0$.
\end{proof}

\begin{lemma}[Gauss content bound]\label{pf:onto1:gaussC}
Let $G$ be a global family datum with $u_k\in\mathcal M_{w_L}$, $\theta^0_k\in\mathcal M_{w_0}$, $|\alpha_k|=a_0$, $\delta\le T$,
$\iota(z_k(\beta))$ the coefficient of $X^\beta$ in $z(\alpha_k,\theta^0_k)$ for all $k,\beta$, $\det(\phi,z)\ne0$ for some $\phi$,
and $\mathcal M_\kappa$ finitely generated. Then for every prime $\mathfrak l$, $\len_{\OO'}(Z_m/Z_l)\le\gamma_{\mathfrak l}$.
\formallabel{Formalism.MainNotions.Gaps.EtaleWeightGaps.len_Zm_le}
\end{lemma}

\begin{proof}
The coefficients $z_k(\beta)$ are sums of products of the integral forms $u_{k'}$ and $\theta^0_k$, of total weight
$a_0w_L+w_0=\kappa$, hence lie in $\mathcal M_\kappa$. Since $\mathcal M_\kappa$ is finitely generated, multiplying a family $\phi$
with $\det(\phi,z)\ne0$ by a suitable non-zero element of $\OO_K$ makes it $\mathfrak l$-integral on $\mathcal M_\kappa$; the
coefficients of the new determinant are non-zero multiples, integral at $\mathfrak l$. Hence the set in the definition of
$\gamma_{\mathfrak l}$ is non-empty, and there are $\phi_1,\dots,\phi_{r_z}$, integral on $\mathcal M_\kappa$ at $\mathfrak l$, and a
monomial $\beta$ such that the coefficient $d_\beta$ of $X^\beta$ in $\det(\phi,z)$ is non-zero with
$\ontooneord_{\mathfrak l}(d_\beta)=\gamma_{\mathfrak l}$.

Extend $\phi_a$ to $K'$-linear forms $\phi'_a$ on $V'_\kappa$ with $\phi'_a(\mathrm{emb}\,\iota f)=\phi_a(f)$; they map
$M_1$ into $\OO'$. Put $z'_k=\sum_\beta X^\beta\,\mathrm{emb}\,\iota z_k(\beta)$; by hypothesis $z'_k=s^{\alpha_k}\mathrm{emb}\,
\iota\theta^0_k$, so $Z_l=\sum_k\OO'z'_k$, and $z'_k\in M_1$. The matrix $P=(\phi'_a(z'_k))_{a,k}$ has entries
$\sum_\beta\phi_a(z_k(\beta))X^\beta\in\OO'$, and $\det P$ is the image of $\det(\phi,z)$ in $K'$; it is non-zero, so the $z'_k$
are linearly independent over $K'$, and it lies in $\OO'$ with valuation $\le\ontooneord_{\mathfrak l}(d_\beta)=\gamma_{\mathfrak l}$,
because the valuation of $\OO'$ on $\OO_{\mathfrak l}[X]$ is the minimum of the valuations of the coefficients (Gauss's
lemma).

Let $U=\sum_kK'z'_k$ and $\Phi=(\phi'_a)_a\colon U\to K'^{r_z}$. In the basis $(z'_k)$ the matrix of $\Phi$ is $P$, so $\Phi$ is
injective, $\Phi(Z_l)=P\,\OO'^{r_z}$, and $\Phi(M_1\cap U)\subseteq\OO'^{r_z}$. Hence $Z_m/Z_l=(M_1\cap U)/Z_l$ embeds into
$\OO'^{r_z}/P\,\OO'^{r_z}$, whose length is the valuation of $\det P$ by the theory of elementary divisors. Therefore
$\len(Z_m/Z_l)\le\ontooneord(\det P)\le\gamma_{\mathfrak l}$.
\end{proof}

\begin{theorem}[The local bound]\label{pf:onto1:localBound}
Let $G$ be a global family datum, $\mathfrak l$ a prime, $Q$ a local point datum at $\mathfrak l$, $C_R,c,\nu_s\in\R$ and
$e\in\N$. Assume:
\begin{enumerate}[label=(\alph*)]
\item $\mathcal M_\kappa$ is a lattice in $\mathcal F_\kappa$, $\Theta\subseteq\ontooneMsup_\kappa$ and $K\mathrm{Th}_K=\mathcal F_\kappa$;
\item $u_k\in W\cap\mathcal M_{w_L}$, $\theta'_i\in\mathrm{Th}_F\cap\mathcal M_{w_F}$, $\theta^0_k\in\mathrm{Th}_0\cap\mathcal M_{w_0}$ and
$|\alpha_k|=a_0$ for all indices; $\iota(z_k(\beta))$ is the coefficient of $X^\beta$ in $z(\alpha_k,\theta^0_k)$ for all
$k,\beta$; and $\det(\phi,z)\ne0$ for some $\phi$;
\item $\theta_i\in\mathrm{Th}_K$ for all $i$, and $\sum_iK\theta_i+\mathcal M^{\mathrm{gg}}_\kappa(C_R,T,c)=\mathcal F_\kappa$;
\item $(\mathrm B_{\mathfrak l})$, $(\mathrm C_{\mathfrak l})$, $(\mathrm T_{\mathfrak l})$ and $(\mathrm M_{\mathfrak l})$ (with $C_R,c$) hold;
\item $v_\eta(\zeta)\ge0$ and $\nu(\zeta)\le\nu_s+\max(0,c)\,v_\eta(\zeta)$ for all $\zeta$, and $T\nu_s\le e$;
\item (base length) $\ontooneql(\mathrm{AA},N_1')\le\ontoonebin I_{\mathfrak l}^\flat$;
\item (pole length) $\ontooneql(\mathrm{AA}^+(m_e),\mathrm{AA})\le\ontoonebin\sum_\zeta m_e(\zeta)f_\zeta$.
\end{enumerate}
Then for every $K_{\mathfrak l}$-subspace $V'\subseteq K_{\mathfrak l}\ontooneMsup_\kappa$ the length
$\ontooneql_{\mathfrak l}(\ontooneMsup_{\kappa,\mathfrak l}\cap V',\Theta_{\mathfrak l}\cap V')$ is finite and at most
$e\dim V'+r(\mathfrak l)$.
\formallabel{Formalism.MainNotions.Gaps.EtaleWeightGaps.localBound}
\end{theorem}

\begin{proof}
Write $W_{\mathfrak l}=K_{\mathfrak l}\ontooneMsup_\kappa$, $M=\ontooneMsup_{\kappa,\mathfrak l}$ and
$M^{\mathrm{gg}}=\bigl(\iota(\mathcal M_\kappa\cap\mathcal M^{\mathrm{gg}}_\kappa)\bigr)_{\mathfrak l}$; these and $\Theta_{\mathfrak
l}$ are $\OO_{\mathfrak l}$-submodules of $W_{\mathfrak l}$.

\emph{Step 1: lattices.} Since $\mathcal M_\kappa$ is finitely generated and spans $\mathcal F_\kappa$, $M$ is a lattice in
$W_{\mathfrak l}$. Since $\Theta\subseteq\ontooneMsup_\kappa$, $\Theta_{\mathfrak l}\subseteq M$ is finitely generated, and it
contains $\iota\mathrm{Th}_K$, which spans $\iota\mathcal F_\kappa$; hence $\Theta_{\mathfrak l}$ is a lattice in $W_{\mathfrak
l}$. The map $j\colon W_{\mathfrak l}\to V'_\kappa$ induced by $\mathrm{emb}$ exhibits $V'_\kappa$ as $K'\otimes_{K_{\mathfrak
l}}W_{\mathfrak l}$, and $\OO'j(M)=M_1$.

\emph{Step 2: the points.} By $(\mathrm C_{\mathfrak l})$ and \cref{pf:onto1:intclos}, $\mathrm{Pts}$ is finite, every
$\bar\OO_\zeta$ is a discrete valuation ring, and $e_\zeta\ge1$ because $\zeta$ lies over the maximal ideal
$\varpi\OO'$. By $(\mathrm T_{\mathfrak l})$, $e_\zeta\nu(\zeta)\in\Z$.

\emph{Step 3: the hypotheses of \cref{pf:onto1:levelsQDesc}.} We apply the proposition with $\OO=\OO_{\mathfrak l}$,
the Gauss extension $\OO'$, $V=W_{\mathfrak l}$, $V_1=V'_\kappa$, $W'=\mathrm{WK}(w_F,a_0)$, the lattices $M,\Theta_{\mathfrak l}$,
the module $M^{\mathrm{gg}}$, the elements $\theta_i$, the modules $N_1$, $Z=Z_l$, $Z_m$, $N'=N'_e$, $A=\mathrm{AA}$,
$A'=\mathrm{AA}^+(m_e)$, $N_1'$, the map $\rho_e$, the set $\mathrm{Pts}$ with $e_\zeta,f_\zeta,\nu(\zeta),v_\eta(\zeta)$, and
$R=\ontoonebin$, $I=I_{\mathfrak l}^\flat$, $c'=\max(0,c)$, $q=\lambda_{\mathfrak l}$. We check the hypotheses.
\begin{itemize}
\item $N_1+Z_l\subseteq\OO'j(\Theta_{\mathfrak l})$. Each $s_i$ is an $\OO'$-combination of the elements
$\mathrm{emb}\,\iota(u_k)$ with $u_k\in W$. Hence for $|\alpha|<a_0$ the element $s^\alpha s_0^{T-|\alpha|}\mathrm{emb}\,\iota\theta'_i$
lies in the $\OO'$-span of $\mathrm{emb}(\iota W^T\cdot\iota\mathrm{Th}_F)$, and since $|\alpha_k|=a_0=T-\delta$, the element
$s^{\alpha_k}\mathrm{emb}\,\iota\theta^0_k$ lies in the $\OO'$-span of $\mathrm{emb}(\iota W^{T-\delta}\cdot\iota\mathrm{Th}_0)$; both are
contained in $\OO'j(\Theta_{\mathfrak l})$.
\item $\theta_i\in M\cap\Theta_{\mathfrak l}$ because $\theta_i\in\mathcal M_\kappa$ and $\iota\theta_i\in\iota\mathrm{Th}_K\subseteq\Theta$;
and $M^{\mathrm{gg}}\subseteq M$ trivially.
\item $j(M^{\mathrm{gg}})\subseteq K'(N_1+Z_l)$: this is $(\mathrm M_{\mathfrak l})$, since $K'(N_1+Z_l)$ is the $K'$-span of the
generators listed there.
\item $\len_{\OO_{\mathfrak l}}\bigl(M/(M^{\mathrm{gg}}+\sum_i\OO_{\mathfrak l}\theta_i)\bigr)\le\lambda_{\mathfrak l}$. The module on
the left is a quotient of $\OO_{\mathfrak l}\otimes_{\OO_K}(Q_\kappa/\sum_i\OO_K\bar\theta_i)$; since $\OO_{\mathfrak l}$ is flat
over $\OO_K$ and $\OO_{\mathfrak l}/\mathfrak l\OO_{\mathfrak l}=\OO_K/\mathfrak l$, the length of this tensor product equals
$\ontoonelg_{\mathfrak l}(Q_\kappa/\sum_i\OO_K\bar\theta_i)$, which is finite because by (c) this module is a finitely
generated torsion module.
\item $\varpi^eM_1\cap K'(N_1+Z_l)\subseteq N'_e$, $N_1\subseteq N'_e$ and $Z_m\subseteq N'_e$ hold by definition, and
$Z_l\subseteq Z_m$ because $Z_l\subseteq K'Z_l$ and $Z_l\subseteq\OO'j(\Theta_{\mathfrak l})\subseteq\OO'j(M)=M_1$.
\item $\rho_e(N'_e)\subseteq\mathrm{AA}^+(m_e)$. This is
\lref{Formalism.MainNotions.Gaps.EtaleWeightGaps.rhoN_mem}, which states that, under $(\mathrm C_{\mathfrak l})$,
$(\mathrm T_{\mathfrak l})$, $u_k\in W$, $\theta'_i\in\mathcal M_{w_F}$, $\delta\le T$ and $|\alpha_k|=a_0$, the map $\rho_T$ sends
$N'_e$ into $\mathrm{AA}^+(m_e)$. Its proof treats the three summands of $N'_e$: $\rho_T$ kills $K'Z_l$
(\lref{Formalism.MainNotions.Gaps.EtaleWeightGaps.ZK_le_ker}, which states exactly this, under $\delta\le T$
and $|\alpha_k|=a_0$; note that $s^{\alpha_k}/s_0^{a_0}\in J^{a_0}$ and $\mathrm{emb}\,\iota\theta^0_k/s_0^{\delta}\in H(w_F)$), it
maps $N_1$ into $\mathrm{AA}$ (\cref{pf:onto1:good}(c), which uses $(\mathrm C_{\mathfrak l})$, $(\mathrm T_{\mathfrak l})$, $u_k\in W$
and $\theta'_i\in\mathcal M_{w_F}$),
and it maps $\varpi^eM_1$ into $\mathrm{AA}^+(m_e)$
(\lref{Formalism.MainNotions.Gaps.EtaleWeightGaps.reduction_uniformizer_smul_mem}, which states this under
$(\mathrm C_{\mathfrak l})$, $(\mathrm T_{\mathfrak l})$ and $u_k\in W$; its proof shows that $\rho_T(M_1)\subseteq\mathrm{AA}^+(m_0)$
and that multiplication by $\varpi^e$ maps $\mathrm{AA}^+(m_0)$ into $\mathrm{AA}^+(m_e)$).
\item $\ker\rho_e=Z_m$. This is \lref{Formalism.MainNotions.Gaps.EtaleWeightGaps.ker_rhoN}, which states it
under the injectivity of \cref{pf:onto1:in64}, $N_1\subseteq M_1$, $\delta\le T$ and $|\alpha_k|=a_0$; we recall the
argument. We have $K'Z_l\subseteq\ker\rho_T$ and $N'_e\subseteq M_1\cap K'(N_1+Z_l)$, and by \cref{pf:onto1:in64} (which
uses $(\mathrm B_{\mathfrak l})$) an element of the $K'$-span of the generators of $N_1$ that is killed by $\rho_T$ is zero.
Hence an element $n+z$ of $N'_e$ with $n\in K'N_1$, $z\in K'Z_l$ is killed by $\rho_e$ if and only if $n=0$, i.e.\ if and
only if it lies in $M_1\cap K'Z_l=Z_m$.
\item $\rho_e(N_1)=N_1'$ by definition; $N_1'\subseteq\mathrm{AA}$ as above; and $\mathrm{AA}\subseteq\mathrm{AA}^+(m_e)$ because
$m_e\ge0$.
\item $\len_{\OO'}(Z_m/Z_l)\le\gamma_{\mathfrak l}$, in particular it is finite, by \cref{pf:onto1:gaussC} (whose hypotheses
are contained in (a) and (b)).
\item Hypothesis (d) of \cref{pf:onto1:levelsQDesc} holds in its first form by (e) and Step 2, and (e) of
\cref{pf:onto1:levelsQDesc} is (f) together with (g).
\end{itemize}

\emph{Step 4: conclusion.} \Cref{pf:onto1:levelsQDesc} gives, for every subspace $U\subseteq W_{\mathfrak l}$,
\[
\ontooneidx(M\cap U,\Theta_{\mathfrak l}\cap U)\le e\dim U+\ontoonebin\,T\max(0,c)\sum_\zeta e_\zeta f_\zeta
v_\eta(\zeta)+\ontoonebin I_{\mathfrak l}^\flat+\gamma_{\mathfrak l}+\lambda_{\mathfrak l}.
\]
By the definition of $v_\eta$, $T\max(0,c)\sum_\zeta e_\zeta f_\zeta v_\eta(\zeta)=[j<m]\,T\frac{\max(0,c)}{2i_\eta}
\mathrm{Net}_{\mathfrak l}$, so the right side is $e\dim U+r(\mathfrak l)$. Since $\Theta_{\mathfrak l}\subseteq M$, we have
$\ontooneql(\Theta_{\mathfrak l}\cap U,M\cap U)=0$, so the index equals $\ontooneql(M\cap U,\Theta_{\mathfrak l}\cap U)$,
which is finite as both modules are lattices in $U$. For $V'\subseteq W_{\mathfrak l}$ the length in the statement,
computed in $V_{\mathfrak l}$, equals this length for $U=V'$.
\end{proof}

\begin{remark}\label{pf:onto1:rem-local}
The two length bounds (f) and (g) are hypotheses of \cref{pf:onto1:localBound}. They are proved in
\cref{sec:pf:onto1:lengths}: (f) holds whenever $(\mathrm B_{\mathfrak l})$, $(\mathrm C_{\mathfrak l})$, $(\mathrm T_{\mathfrak l})$,
$(\mathrm Z_{\mathfrak l})$, $(\mathrm M_{\mathfrak l})$, $\theta'_i\in\mathcal M_{w_F}$, $u_k\in W$ and $I_{\mathfrak l}<\infty$ hold
(\cref{pf:onto1:lenBase}), and (g) holds for every $e$ whenever $(\mathrm B_{\mathfrak l})$, $(\mathrm C_{\mathfrak l})$, $(\mathrm
T_{\mathfrak l})$, $(\mathrm Z_{\mathfrak l})$, $(\mathrm M_{\mathfrak l})$, $(\mathrm E_{\mathfrak l})$ and $u_k\in W$ hold
(\cref{pf:onto1:lenPole}). The factor $\binom{T-1}{\delta}=\binom{a_0+\delta-1}{\delta}$ is the number of monomials of degree
$<a_0$ in $\delta$ variables.
\end{remark}

%----------------------------------------------------------------------------------------------------
\subsection{The base and pole lengths}\label{sec:pf:onto1:lengths}

In this subsection we prove the length bounds (f) and (g) of \cref{pf:onto1:localBound}. We keep the notation of
\cref{sec:pf:onto1:local}: $G$ is a global family datum, $\mathfrak l$ a prime and $Q$ a local point datum at $\mathfrak l$, and
we write $H=H(w_F)$, $\mathrm{SQ}=\mathrm{SQ}(w_F)$, $\mathrm{WK}=\mathrm{WK}(w_F,a_0)$ and $\mathrm{res}\colon H\to\mathrm{SQ}$ for the
reduction modulo $JH$. Put
\[
y_l=s_l/s_0\ (1\le l\le\delta),\qquad \omega_k=\mathrm{emb}\,\iota(u_k)/s_0\ (k<N),\qquad t'_i=\mathrm{emb}\,\iota\theta'_i .
\]
The $y_l$ and $\omega_k$ lie in $H(0)=\mathrm{Fr}$ and the $t'_i$ in $H$; we denote the image of $\omega_k$ in $\OO_Z$ again by
$\omega_k$, and $\mathrm{res}(t'_i)=t_i$. Since $s_0=\sum_kX_{0,k}\,\mathrm{emb}\,\iota(u_k)$,
\begin{equation}\label{pf:onto1:sumone}
\sum_{k<N}X_{0,k}\,\omega_k=1 .
\end{equation}
The ideal $J$ of $\mathrm{Fr}$ is generated by $y_1,\dots,y_\delta$, so $J^aH=\sum_{|\mu|=a}y^\mu H$ for $a\in\N$; the image of
$J^kH$ in $\mathrm{WK}=H/J^{a_0}H$ is denoted $F_k$. For $m\colon\mathrm{Pts}\to\Z$ we put
\[
B(m)=\{v\in\mathrm{SQ}:\ \ontooneord^Q_\zeta(v)\ge-m(\zeta)\ \text{for all }\zeta\},
\]
so that $A_0=B(0)$; under $(\mathrm C_{\mathfrak l})$, $B(m)$ is an $\bar\OO$-submodule of $\mathrm{SQ}$. An element $x\in H$ is
\emph{$m$-good} if $\mathrm{res}(\partial^\beta x)\in B(m)$ for every multi-index $\beta$ with $|\beta|\le a_0-1$. By definition,
$\mathrm{AA}^+(m)$ is the $\OO'$-span of the images in $\mathrm{WK}$ of the $m$-good elements, and $\mathrm{AA}=\mathrm{AA}^+(0)$. For a
column $k<N$ and $\alpha\in\N^\delta$ let $\beta(k,\alpha)$ be the multi-index with entry $\alpha_l$ at the variable $X_{l,k}$
($1\le l\le\delta$) and $0$ at all other variables, and let $\mathrm{mon}(a)=\{\mu\in\N^\delta:|\mu|=a\}$.

\paragraph{Hasse derivatives.} The Hasse derivatives $\partial^\beta$ in the variables $X_{i,k}$ act on $L'$ as
$K_{\mathfrak l}$-linear operators with $\partial^0=\mathrm{id}$; they act on $K'=K_{\mathfrak l}(X)$ as the Hasse derivatives of
rational functions, kill $\mathrm{emb}(\mathcal A)$ for $\beta\ne0$, and satisfy the Leibniz rule
$\partial^\beta(fg)=\sum_{\beta'+\beta''=\beta}\partial^{\beta'}f\,\partial^{\beta''}g$. We use the following consequences.
\begin{enumerate}[label=(\roman*)]
\item\label{pf:onto1:hs1} $\partial^\beta(\OO')\subseteq\OO'$: a Hasse derivative of $p/q$, $p,q\in\OO_{\mathfrak l}[X]$, has the form
$p'/q^{n+1}$ with $p'\in\OO_{\mathfrak l}[X]$.
\item\label{pf:onto1:hs2} The subring $\Z[y,\omega]\subseteq\mathrm{Fr}$ generated by the $y_l$ and the $\omega_k$ is stable under
all $\partial^\beta$. Indeed $s_0$ is linear in the variables $X_{0,k}$ and does not involve the others, so
$\partial^{(n)}_{X_{0,k}}s_0^{-1}=(-\mathrm{emb}\,\iota u_k)^ns_0^{-n-1}$, whence $\partial^{(n)}_{X_{0,k}}y_l=y_l(-\omega_k)^n$ and
$\partial^{(n)}_{X_{0,k}}\omega_{k'}=\omega_{k'}(-\omega_k)^n$; moreover $y_l$ is linear in $X_{l,k}$ with
$\partial_{X_{l,k}}y_{l'}=\delta_{ll'}\omega_k$ for $l\ge1$, and $\omega_{k'}$ does not involve $X_{l,k}$ for $l\ge1$.
\item\label{pf:onto1:hs3} $\partial^\beta(H)\subseteq H$, and $\partial^\beta(J^aH)\subseteq J^{a-|\beta|}H$ for $|\beta|\le a$ (by the
Leibniz rule, \ref{pf:onto1:hs2} and the same computation for $s_0^{-k}$).
\item\label{pf:onto1:hs4} For $k<N$ and $\gamma,\alpha\in\N^\delta$ with $|\gamma|\le|\alpha|$: $\partial^{\beta(k,\alpha)}(y^\gamma)=
\omega_k^{|\alpha|}$ if $\gamma=\alpha$, and $=0$ otherwise. Indeed, by \ref{pf:onto1:hs2},
$\partial^{\beta(k,\alpha)}(y^\gamma)=\prod_l\binom{\gamma_l}{\alpha_l}y_l^{\gamma_l-\alpha_l}\omega_k^{\alpha_l}$, which vanishes as soon as
$\alpha_l>\gamma_l$ for some $l$, and this is the case if $|\gamma|\le|\alpha|$ and $\gamma\ne\alpha$.
\end{enumerate}

\begin{lemma}[Filtrations]\label{pf:onto1:filt}
Let $R$ be a commutative ring, $W$ an $R$-module, $a\in\N$ and $F_0=W\supseteq F_1\supseteq\dots$ submodules with $F_a=0$.
For $k\in\N$ let $\phi_k\colon F_k\to V_k$ be $R$-linear maps with $\ker\phi_k\subseteq F_{k+1}$. Let $N\subseteq M\subseteq W$ and
$P_k,Q_k\subseteq V_k$ be submodules with $\phi_k(F_k\cap M)\subseteq P_k$ and $Q_k\subseteq\phi_k(F_k\cap N)$ for $k<a$. Then
\[
\len_R(M/N)\le\sum_{k<a}\len_R\bigl(P_k/(P_k\cap Q_k)\bigr).
\]
\formallabel{filtration_length_bound}
\formallabel{PoleLengthBounds.filtLenGr_holds}
\end{lemma}

\begin{proof}
Put $M_k=M\cap F_k$ and $N_k=N\cap F_k$. The chain $M=M_0+N\supseteq M_1+N\supseteq\dots\supseteq M_a+N=N$ gives
$\len(M/N)=\sum_{k<a}\len\bigl((M_k+N)/(M_{k+1}+N)\bigr)$. The module $(M_k+N)/(M_{k+1}+N)$ is a quotient of
$M_k/(M_{k+1}+N_k)$. Consider the composite $M_k\to P_k\to P_k/(P_k\cap Q_k)$ of $\phi_k$ and the projection. If $x\in M_k$
lies in its kernel, then $\phi_k(x)\in Q_k$, so $\phi_k(x)=\phi_k(n)$ with $n\in N_k$, and $x-n\in\ker\phi_k\cap M\subseteq
M_{k+1}$; hence the kernel is contained in $M_{k+1}+N_k$, and $\len(M_k/(M_{k+1}+N_k))\le\len(P_k/(P_k\cap Q_k))$.
\end{proof}

\begin{lemma}[Counting monomials]\label{pf:onto1:monos}
For natural numbers $\delta<T$, $\sum_{k<T-\delta}|\mathrm{mon}(k)|=\binom{T-1}{\delta}$, where $\mathrm{mon}(k)$ is taken in
$\N^\delta$.
\formallabel{sum_antidiagonalTuple}
\formallabel{PoleLengthBounds.monoCount_holds}
\end{lemma}

\begin{proof}
$|\mathrm{mon}(k)|=\binom{k+\delta-1}{\delta-1}$ for $\delta\ge1$, and $\sum_{k<n}\binom{k+\delta-1}{\delta-1}=\binom{n+\delta-1}{\delta}$
(induction on $n$, Pascal's rule); take $n=T-\delta$. For $\delta=0$ both sides equal $1$.
\end{proof}

\begin{lemma}[Columns]\label{pf:onto1:columns}
Assume $(\mathrm C_{\mathfrak l})$, $(\mathrm T_{\mathfrak l})$ and $u_k\in W$ for all $k$. Then for every point $\zeta$,
$\ontooneord_\zeta(\omega_k)\ge0$ for all $k$ and $\ontooneord_\zeta(\omega_k)=0$ for some $k$. Consequently, for every $n\in\N$:
\begin{enumerate}[label=(\alph*)]
\item if $b\in\mathrm{SQ}$ and $\omega_k^nb=0$ for all $k$, then $b=0$;
\item if $b\in\mathrm{SQ}$, $m\colon\mathrm{Pts}\to\Z$ and $\omega_k^nb\in B(m)$ for all $k$, then $b\in B(m)$;
\item there are $a_0',\dots,a_{N-1}'\in\bar\OO$ with $\sum_ka_k'\omega_k^n=1$.
\end{enumerate}
\formallabel{PoleLengthBounds.column_holds}
\formallabel{PoleLengthBounds.uInj_of_column}
\formallabel{PoleLengthBounds.uBall_of_column}
\formallabel{PI.PoleLengthCount.unit_columns_holds}
\formallabel{PI.PoleLengthCount.exists_coef}
\end{lemma}

\begin{proof}
By $(\mathrm T_{\mathfrak l})$ applied to $g=u_k\in\mathcal F_{w_L}$, $\ontooneord_\zeta(\omega_k)=e_\zeta(v_{\mathfrak l,y_\zeta}(\iota u_k)-\nu(\zeta))$,
which is $\ge0$ because $u_k\in W$ and $\nu(\zeta)$ is the infimum of $v_{\mathfrak l,y_\zeta}$ on $\iota W$. If
$\ontooneord_\zeta(\omega_k)\ge1$ for all $k$, then by \eqref{pf:onto1:sumone}, since $X_{0,k}\in\OO'\subseteq\bar\OO$, the element $1$
would have positive order at $\zeta$. For (a) and (b) write $b=\mathrm{co}(b)\varepsilon$ (by $(\mathrm C_{\mathfrak l})$) and, at a
point $\zeta$, choose $k$ with $\ontooneord_\zeta(\omega_k)=0$; then $\ontooneord^Q_\zeta(\omega_k^nb)=\ontooneord^Q_\zeta(b)$. In (a) this
gives $\ontooneord_\zeta(\mathrm{co}(b))=\infty$ at all points, hence $\mathrm{co}(b)=0$ ($\OO_Z$ is reduced and finite over $K'$, so
an element of infinite order at every point of its normalisation is $0$); (b) is immediate. For (c): the $\omega_k$ lie in
$\bar\OO$, which consists of the elements of $\OO_Z$ of non-negative order at all points ($\bar\OO$ is a finite product of
Dedekind domains by \cref{pf:onto1:intclos}), and by \eqref{pf:onto1:sumone} they generate the unit ideal of $\bar\OO$; hence so do
their $n$-th powers.
\end{proof}

\begin{lemma}[Taylor coefficients]\label{pf:onto1:taylor}
Let $k\in\N$ and $x\in J^kH$, and write $x=\sum_{\mu\in\mathrm{mon}(k)}y^\mu c_\mu+r$ with $c_\mu\in H$ and $r\in J^{k+1}H$. Then
for every column $k'<N$ and every $\alpha\in\mathrm{mon}(k)$,
\[
\mathrm{res}\bigl(\partial^{\beta(k',\alpha)}x\bigr)=\omega_{k'}^{\,k}\,\mathrm{res}(c_\alpha).
\]
Moreover $\mathrm{res}(\partial^\beta x)=0$ for $|\beta|<k$, and $\mathrm{res}(\partial^\beta(ox))=o\,\mathrm{res}(\partial^\beta x)$ for $o\in\OO'$
and $|\beta|=k$.
\formallabel{PI.PoleLengthCount.ymDeriv}
\formallabel{PoleLengthBounds.taylorHasse_holds}
\formallabel{PoleLengthBounds.hasseVan_holds}
\formallabel{PoleLengthBounds.hasseSmul_holds}
\end{lemma}

\begin{proof}
By the Leibniz rule, $\partial^{\beta(k',\alpha)}(y^\mu c_\mu)$ is the sum of the products $\partial^{\beta(k',\gamma)}(y^\mu)\,
\partial^{\beta(k',\alpha-\gamma)}(c_\mu)$ over $\gamma\le\alpha$. For $\gamma\ne\alpha$ we have $|\gamma|<|\mu|$, and
$\partial^{\beta(k',\gamma)}(y^\mu)=\prod_l\binom{\mu_l}{\gamma_l}y_l^{\mu_l-\gamma_l}\omega_{k'}^{\gamma_l}$ lies in $J$; for $\gamma=\alpha$,
$\partial^{\beta(k',\alpha)}(y^\mu)=\delta_{\mu\alpha}\omega_{k'}^k$ by \ref{pf:onto1:hs4}. Since $\partial^{\beta''}(c_\mu)\in H$ by
\ref{pf:onto1:hs3}, the restriction of $\partial^{\beta(k',\alpha)}(\sum_\mu y^\mu c_\mu)$ is $\omega_{k'}^k\mathrm{res}(c_\alpha)$, and
$\partial^{\beta(k',\alpha)}r\in JH$ by \ref{pf:onto1:hs3}. The second assertion follows from $\partial^\beta(J^kH)\subseteq JH$
for $|\beta|<k$. For the third, $\partial^\beta(ox)=\sum_{\beta'+\beta''=\beta}\partial^{\beta'}(o)\,\partial^{\beta''}(x)$, where
$\partial^{\beta'}(o)\in\OO'$ and, for $\beta''\ne\beta$, $\mathrm{res}(\partial^{\beta''}x)=0$ by the second assertion.
\end{proof}

\begin{lemma}[Good elements]\label{pf:onto1:good}
Let $m\colon\mathrm{Pts}\to\Z$ and assume $(\mathrm C_{\mathfrak l})$.
\begin{enumerate}[label=(\alph*)]
\item The $m$-good elements form an $\OO'$-submodule of $H$ containing $J^{a_0}H$. Hence $x\in H$ is $m$-good if and only if
its image in $\mathrm{WK}$ lies in $\mathrm{AA}^+(m)$.
\item Assume moreover $(\mathrm T_{\mathfrak l})$ and $u_k\in W$ for all $k$. If $x$ is $m$-good, then so is $y^\alpha x$ for every
$\alpha\in\N^\delta$.
\item Assume moreover $\theta'_i\in\mathcal M_{w_F}$ for all $i$. Then $N_1'\subseteq\mathrm{AA}$.
\end{enumerate}
\formallabel{PoleLengthBounds.aaRep_holds}
\formallabel{PI.PoleLengthCount.liftMul_of}
\formallabel{Formalism.MainNotions.Gaps.EtaleWeightGaps.reduced_generators_le}
\end{lemma}

\begin{proof}
(a) For $o\in\OO'$ and $m$-good $x$, $\mathrm{res}(\partial^\beta(ox))=\sum_{\beta'+\beta''=\beta}\partial^{\beta'}(o)\,
\mathrm{res}(\partial^{\beta''}x)$ lies in $B(m)$, because $\partial^{\beta'}(o)\in\OO'\subseteq\bar\OO$ by \ref{pf:onto1:hs1}, $|\beta''|\le
a_0-1$ and $B(m)$ is an $\bar\OO$-module; sums of good elements are good. Elements of $J^{a_0}H$ are good by
\ref{pf:onto1:hs3}. As $\mathrm{WK}=H/J^{a_0}H$, the preimage of $\mathrm{AA}^+(m)$ in $H$ is therefore the set of $m$-good
elements.

(b) By \ref{pf:onto1:hs2}, every $\partial^{\beta'}(y^\alpha)$ lies in $\Z[y,\omega]$, and its image in $\OO_Z$ has non-negative
order at every point, since $\mathrm{res}(y_l)=0$ and $\ontooneord_\zeta(\omega_k)\ge0$ (\cref{pf:onto1:columns}); it therefore
lies in $\bar\OO$. By the Leibniz rule, $\mathrm{res}(\partial^\beta(y^\alpha x))=\sum\mathrm{res}(\partial^{\beta'}y^\alpha)\,\mathrm{res}
(\partial^{\beta''}x)\in B(m)$.

(c) $N_1'$ is the $\OO'$-span of the images of the elements $\rho_T(s^\alpha s_0^{T-|\alpha|}\mathrm{emb}\,\iota\theta'_i)$, which are
the classes of $y^\alpha t'_i$ ($|\alpha|<a_0$). As $t'_i$ does not involve the variables,
$\mathrm{res}(\partial^\beta(y^\alpha t'_i))=\mathrm{res}(\partial^\beta y^\alpha)\,t_i$, where the first factor lies in $\bar\OO$ (proof
of (b)) and $\ontooneord^Q_\zeta(t_i)=e_\zeta v_{\mathfrak l,y_\zeta}(\iota\theta'_i)\ge0$ by $(\mathrm T_{\mathfrak l})$ (with $k=0$) and
the integrality of $\theta'_i$. Hence $y^\alpha t'_i$ is $0$-good and its class lies in $\mathrm{AA}$; as $\mathrm{AA}$ is an
$\OO'$-module, $N_1'\subseteq\mathrm{AA}$.
\end{proof}

\begin{lemma}[Lengths of balls]\label{pf:onto1:balls}
Assume $(\mathrm C_{\mathfrak l})$ and let $m\colon\mathrm{Pts}\to\N$. Then $A_0\subseteq B(m)$ and
\[
\len_{\OO'}\bigl(B(m)/A_0\bigr)=\sum_{\zeta\in\mathrm{Pts}}m(\zeta)f_\zeta .
\]
\formallabel{ball_quotient_length}
\formallabel{PI.UnitaryGroupMassFormula.BallQuotLen.length_ball_quotient_eq_sum_mul_fdeg}
\end{lemma}

\begin{proof}
By $(\mathrm C_{\mathfrak l})$ and \cref{pf:onto1:intclos}, $\mathrm{Pts}$ is finite, $\bar\OO$ is a finite product of semi-local Dedekind
domains whose localisations $\bar\OO_\zeta$ are discrete valuation rings, and $v\mapsto\mathrm{co}(v)$ identifies $B(m)$ with
the $\bar\OO$-module of $a\in\OO_Z$ with $\ontooneord_\zeta(a)\ge-m(\zeta)-c(\zeta)$ for all $\zeta$. We argue by induction on
$\sum_\zeta m(\zeta)$, the case $m=0$ being trivial. If $m=m'+\mathbf 1_{\zeta_0}$ with $m'\ge0$, then $B(m')\subseteq B(m)$, and
$v\mapsto$ (the class of $\mathrm{co}(v)$ in $P^{-n}/P^{-n+1}$, where $P$ is the maximal ideal of $\bar\OO_{\zeta_0}$ and
$n=m(\zeta_0)+c(\zeta_0)$) maps $B(m)$ onto this one-dimensional vector space over the residue field $\kappa(\zeta_0)$, with
kernel $B(m')$; surjectivity follows from the approximation theorem in $\bar\OO$. Since $\kappa(\zeta_0)$ has length
$f_{\zeta_0}$ over $\OO'$, $\len(B(m)/A_0)=\len(B(m')/A_0)+f_{\zeta_0}$.
\end{proof}

\begin{theorem}[Base length]\label{pf:onto1:lenBase}
Let $G$, $\mathfrak l$, $Q$, $C_R$ and $c$ be as in \cref{sec:pf:onto1:local}, and assume $(\mathrm B_{\mathfrak l})$, $(\mathrm C_{\mathfrak
l})$, $(\mathrm T_{\mathfrak l})$, $(\mathrm Z_{\mathfrak l})$, $(\mathrm M_{\mathfrak l})$ (with $C_R,c$), $\theta'_i\in\mathcal M_{w_F}$ for all
$i$, $u_k\in W$ for all $k$, and $I_{\mathfrak l}<\infty$. Then
\[
\ontooneql\bigl(\mathrm{AA},N_1'\bigr)\le\ontoonebin\,I_{\mathfrak l}.
\]
\formallabel{Formalism.MainNotions.Gaps.EtaleWeightGaps.lenBase_all}
\formallabel{PoleLengthBounds.lenBase_holds}
\end{theorem}

\begin{proof}
We apply \cref{pf:onto1:filt} over $\OO'$ to $W=\mathrm{WK}$, the filtration $(F_k)$ (so $F_0=\mathrm{WK}$ and $F_{a_0}=0$), $a=a_0$,
$M=\mathrm{AA}$ and $N=\mathrm{AA}\cap N_1'$. For $k<a_0$ let $V_k$ be the module of families $(b_{k',\alpha})$ of elements of
$\mathrm{SQ}$ indexed by the columns $k'<N$ and $\alpha\in\mathrm{mon}(k)$, and define $\phi_k\colon F_k\to V_k$ by
$\phi_k(\bar x)_{k',\alpha}=\mathrm{res}(\partial^{\beta(k',\alpha)}x)$ for a lift $x\in J^kH$ of $\bar x$ (for $k\ge a_0$, $F_k=0$ and
$\phi_k=0$). Two lifts differ by an element of $J^{a_0}H\subseteq J^{k+1}H$, so $\phi_k$ is well defined by
\cref{pf:onto1:taylor}, which also shows that it is $\OO'$-linear and that
\[
\phi_k(\bar x)=U_k\bigl((\mathrm{res}\,c_\alpha)_\alpha\bigr),\qquad U_k(b)_{k',\alpha}=\omega_{k'}^{\,k}b_\alpha ,
\]
for any representation $x=\sum_\mu y^\mu c_\mu+r$ as there. By \cref{pf:onto1:columns}(a), $U_k$ is injective. Put
$P_k=U_k(A_0^{\mathrm{mon}(k)})$ and $Q_k=U_k\bigl((\sum_i\OO't_i)^{\mathrm{mon}(k)}\bigr)$.

\emph{Kernel.} If $\phi_k(\bar x)=0$, then $\mathrm{res}(c_\alpha)=0$, i.e.\ $c_\alpha\in JH$, for all $\alpha$, so $x\in J^{k+1}H$ and
$\bar x\in F_{k+1}$.

\emph{Image of $\mathrm{AA}$.} If $\bar x\in F_k\cap\mathrm{AA}$, its lift $x\in J^kH$ is $0$-good by \cref{pf:onto1:good}(a), so
$\omega_{k'}^k\mathrm{res}(c_\alpha)=\mathrm{res}(\partial^{\beta(k',\alpha)}x)\in A_0$ for all $k'$ (note $|\beta(k',\alpha)|=k\le a_0-1$),
hence $\mathrm{res}(c_\alpha)\in A_0$ by \cref{pf:onto1:columns}(b), and $\phi_k(\bar x)\in P_k$.

\emph{Realisation of $Q_k$.} Let $b_\alpha=\sum_io_{\alpha,i}t_i$ with $o_{\alpha,i}\in\OO'$, and put $x=\sum_\alpha\sum_io_{\alpha,i}
y^\alpha t'_i\in J^kH$. Its image $\bar x$ lies in $N_1'$, because $y^\alpha t'_i$ is the image under $\rho_T$ of the generator
$s^\alpha s_0^{T-k}\mathrm{emb}\,\iota\theta'_i$ of $N_1$ ($k<a_0$), and in $\mathrm{AA}$ by \cref{pf:onto1:good}(c). By
\cref{pf:onto1:taylor} with $c_\alpha=\sum_io_{\alpha,i}t'_i$, $\phi_k(\bar x)=U_k(b)$.

\emph{Graded lengths.} Since $U_k$ is injective, $P_k/(P_k\cap Q_k)\cong\bigl(A_0/(A_0\cap\sum_i\OO't_i)\bigr)^{\mathrm{mon}(k)}$, of
length $|\mathrm{mon}(k)|\,I_{\mathfrak l}$.

\Cref{pf:onto1:filt} gives $\ontooneql(\mathrm{AA},N_1')=\len(\mathrm{AA}/(\mathrm{AA}\cap N_1'))\le\sum_{k<a_0}|\mathrm{mon}(k)|\,I_{\mathfrak
l}$, which equals $\ontoonebin I_{\mathfrak l}$ by \cref{pf:onto1:monos} if $\delta<T$, and is $0$ if $\delta\ge T$ (then
$a_0=0$).
\end{proof}

\begin{theorem}[Pole length]\label{pf:onto1:lenPole}
Let $G$, $\mathfrak l$, $Q$, $C_R$ and $c$ be as in \cref{sec:pf:onto1:local}, and assume $(\mathrm B_{\mathfrak l})$, $(\mathrm C_{\mathfrak
l})$, $(\mathrm T_{\mathfrak l})$, $(\mathrm Z_{\mathfrak l})$, $(\mathrm M_{\mathfrak l})$ (with $C_R,c$), $(\mathrm E_{\mathfrak l})$ and $u_k\in W$
for all $k$. Then for every $e\in\N$,
\[
\ontooneql\bigl(\mathrm{AA}^+(m_e),\mathrm{AA}\bigr)\le\ontoonebin\sum_{\zeta}m_e(\zeta)f_\zeta .
\]
\formallabel{Formalism.MainNotions.Gaps.EtaleWeightGaps.lenPoleL_all}
\formallabel{PI.PoleLengthCount.pole_length_bound}
\end{theorem}

\begin{proof}
The argument holds for every $m\colon\mathrm{Pts}\to\N$ in place of $m_e$. Let $k<a_0$. By \cref{pf:onto1:columns}(c) choose
$a'_{k'}\in\bar\OO$ with $\sum_{k'}a'_{k'}\omega_{k'}^k=1$, and define $\psi_k\colon F_k\to\mathrm{SQ}^{\mathrm{mon}(k)}$ by
$\psi_k(\bar x)_\alpha=\sum_{k'}a'_{k'}\,\mathrm{res}(\partial^{\beta(k',\alpha)}x)$ for a lift $x\in J^kH$. As in the proof of
\cref{pf:onto1:lenBase}, $\psi_k$ is a well defined $\OO'$-linear map, and by \cref{pf:onto1:taylor}
$\psi_k(\bar x)=(\mathrm{res}\,c_\alpha)_\alpha$ for any representation $x=\sum_\mu y^\mu c_\mu+r$; in particular $\ker\psi_k\subseteq
F_{k+1}$. We apply \cref{pf:onto1:filt} with $M=\mathrm{AA}^+(m)$, $N=\mathrm{AA}$ (contained in $M$ since $m\ge0$),
$P_k=B(m)^{\mathrm{mon}(k)}$ and $Q_k=A_0^{\mathrm{mon}(k)}$.

\emph{Image of $\mathrm{AA}^+(m)$.} A lift $x\in J^kH$ of $\bar x\in F_k\cap\mathrm{AA}^+(m)$ is $m$-good by \cref{pf:onto1:good}(a);
since the $a'_{k'}$ lie in $\bar\OO$ and $B(m)$ is an $\bar\OO$-module, $\psi_k(\bar x)\in B(m)^{\mathrm{mon}(k)}$.

\emph{Realisation of $Q_k$.} Let $b\in A_0^{\mathrm{mon}(k)}$. By $(\mathrm E_{\mathfrak l})$ choose $x_\alpha\in H$ with
$\mathrm{res}(x_\alpha)=b_\alpha$ and image in $\mathrm{AA}$; then $x_\alpha$ is $0$-good, $x=\sum_\alpha y^\alpha x_\alpha\in J^kH$ is
$0$-good by \cref{pf:onto1:good}(b), so $\bar x\in F_k\cap\mathrm{AA}$, and $\psi_k(\bar x)=b$.

\emph{Graded lengths.} $\len\bigl(P_k/(P_k\cap Q_k)\bigr)=|\mathrm{mon}(k)|\,\len(B(m)/A_0)=|\mathrm{mon}(k)|\sum_\zeta m(\zeta)f_\zeta$ by
\cref{pf:onto1:balls}.

\Cref{pf:onto1:filt} and \cref{pf:onto1:monos} conclude as in the proof of \cref{pf:onto1:lenBase}.
\end{proof}

\begin{remark}\label{pf:onto1:rem-lengths}
In \cref{pf:onto1:lenBase,pf:onto1:lenPole} the conditions $(\mathrm B_{\mathfrak l})$, $(\mathrm Z_{\mathfrak l})$ and $(\mathrm
M_{\mathfrak l})$ are part of the formal hypotheses, which bundle the local conditions; the proofs above use only
$(\mathrm C_{\mathfrak l})$, $(\mathrm T_{\mathfrak l})$, the integrality of the $\theta'_i$ (base length), $(\mathrm E_{\mathfrak l})$ (pole
length) and $u_k\in W$.
\end{remark}

%----------------------------------------------------------------------------------------------------
\subsection{Theorem Zero}\label{sec:pf:onto1:thmzero}

Theorem Zero takes as input a long list of conditions. We state precisely those that enter the arguments
given here, and describe the purely numerical ones by their role, the explicit inequalities being those of the
formal definitions tagged at the end of each definition.

\begin{definition}[Hypotheses on the data]\label{pf:onto1:hyps}
Let $L$ be a size datum, $X=(X_j)_{1\le j\le m}$ a family of valued form algebras over a form provider $B$, let
\[
\mathbf K=(C_G,A_G,C_f,A_f,C_W,A_W,C_L,A_L,C_{Rm}),\qquad\mathbf H=(C_K,C_A,C_1,A_1,C_c,A_c,C_{Rc})
\]
be two records of real constants, and let $C_h,A_h,C_q,A_q,C_2,A_2,C_R,C_T\in\R$. Let $j$ be in range, $S$ a scalar
quadruple, $\kappa,w_L,w_F$ weights, $T,n_P,n_W\in\N$, $h_0\in\R$ and $\nu$ a real function on the primes of $K$. Put
\[
c=\frac{R(C_R)\,\ontoonePP(C_2,A_2)\,\ontooneDF(w_L)^2}{\min(g_0,g_1)} .
\]
\begin{itemize}
\item[$(\mathrm H)$] The \emph{size hypotheses} are the conjunction of: (H1) parameter bounds such as $x\le fr\le
x^{C_A}$, $x^{-\eps}\le\delta_2$ and $\log g_i\le9\mathrm L+1$; (H2) conditions on the weights and the parameters
$\kappa,w_L,w_F,T,n_P,n_W,h_0$, which include $\kappa=w_F+Tw_L$, $1\le T$, $\delta\le T$, $\log[K:\Q]\le\mathrm{PO}$ and
$\log|d_K|/[K:\Q]\le\mathrm{PO}$; (H3) scalar conditions:
the entries of $S$ are non-zero, $\log|\sigma(s)|$ is bounded above and below at every embedding $\sigma$ for each
entry $s$, and each of $W$, $\mathrm{Th}_F$, $\mathrm{Th}_0$, $\mathrm{Th}_K$ contains a finite set spanning it (resp.\ the whole
space over $K$) consisting of forms of bounded mean log sup, in particular $K\mathrm{Th}_K=\mathcal F_\kappa$; (H4) the
\emph{prime sum bound}: $\nu\ge0$, $\nu$ has finite support and $\sum_{\mathfrak l}\nu(\mathfrak l)\ell(\mathfrak l)\le
\mathrm{PN}\,\ontoonesize(w_L)$; (H5) integrality of the theta lattices: $W\subseteq\mathcal M_{w_L}$,
$\mathrm{Th}_F\subseteq\mathcal M_{w_F}$, $\mathrm{Th}_0\subseteq\mathcal M_{w_0}$ and $\mathrm{Th}_K\subseteq\mathcal M_\kappa$; (H6) numerical
conditions, and the existence of integral forms which are not zero divisors, of weight $w_1$ with
$a^\pm(w_1)=2a^\pm(w_L)$, $n_{w_1}=2n_W-2$ if $j<m$, and of weight $w_2$ with $w_2+w_F=(\delta+4)w_L$, with bounded
sup norms; and (H7) $T\le x^{C_T}$. (Here the parameter $B_x$ of the formal definition is $0$.)
\item[$(\mathrm L)$] (lattice) $\mathcal M_\kappa$ is a lattice in $\mathcal F_\kappa$.
\item[$(\mathrm S)$] (sup norms by $L^2$-norms) $\|F\|_\sigma^2\le\exp\bigl(\ontoonePP(C_G,A_G)(1+\log\ontoonesize\kappa)\bigr)
\|F\|_{2,\sigma}^2$ for all $F\in\mathcal F_\kappa$ and all $\sigma$.
\item[$(\mathrm A)$] (lower bound for integral forms) $\ontoonemls(F)\ge-\ontoonePPz(C_L,A_L)(1+\log\ontoonesize\kappa)
\ontoonesize\kappa$ for $0\ne F\in\mathcal M_\kappa$.
\item[$(\mathrm Q)$] (quotient bound) If $C_{Rm}\ge0$, $\log[K:\Q]\le\ontoonePP(C_h,A_h)$, $\log|d_K|^{1/[K:\Q]}\le
\ontoonePP(C_h,A_h)$, $n_\kappa=n_V$, $a^+_\tau(p_j-1)\ge1$ for all $\tau$, $\ontoonesize\kappa\le C_Kkfb$, $T>0$, $c>0$ and
$Tc\le1$, then there is a finitely supported $b_e\ge0$ with $\sum_{\mathfrak l}b_e(\mathfrak l)\ell(\mathfrak l)\le
\ontoonePP(C_q,A_q)\ontoonesize\kappa$ such that for every family $\theta_1,\dots,\theta_{\rho_P}\in\mathcal M_\kappa$ with
$\sum K\theta_i+\mathcal M^{\mathrm{gg}}_\kappa(C_{Rm},T,c)=\mathcal F_\kappa$ and $\ontoonemls(\theta_i)\le\ontoonePP(C_h,A_h)
\ontoonesize\kappa$ there is a finitely supported $n_n$ such that $\ontoonelg_{\mathfrak l}(Q_\kappa/\sum\OO_K\bar\theta_i)$ is
finite and at most $n_n(\mathfrak l)+\rho_Pb_e(\mathfrak l)$ for all $\mathfrak l$, and $\sum_{\mathfrak l}n_n(\mathfrak l)\ell(\mathfrak
l)\le\rho_P\ontoonePP(C_q,A_q)\ontoonesize\kappa$.
\item[$(\mathrm K)$] (ties of constants) $C_1=C_c=C_2$, $A_1=A_c=A_2$, $C_{Rc}=C_R=C_{Rm}=7$, $12\le A_2$, $A_G,A_f,A_W\le2$,
$A_L\le4$, $1\le C_h$, $4\le A_h$ and $A_q\le11$.
\item[$(\mathrm{PA})$] (page parameters) the size datum $L$ admits page parameters with constants $(C_2,A_2)$.
\item[$(\mathrm{PS})$] (page support) an explicit polynomial in $n_W,m,f,r$ is at most $\ontoonePP(C_2,A_2)$;
$\frac12\ontooneDF(w_L)\log(g_0g_1)\le\log|\sigma(s_W)|$ for all $\sigma$; $f^2\le\mathrm{rd}_E$; $8\le r$; and
$3\le\ontooneDF(w_L)$.
\end{itemize}
\formal{Formalism.MainNotions.Gaps.EtaleWeightGaps.HypC}
\formal{Formalism.MainNotions.WeightAlgebra.cc}
\formal{Formalism.MainNotions.Gaps.EtaleWeightGaps.LatAt}
\formal{Formalism.MainNotions.Gaps.EtaleWeightGaps.SupBoundByL2AtWeight}
\formal{Formalism.MainNotions.Gaps.EtaleWeightGaps.LAMwAt}
\formal{Formalism.MainNotions.Gaps.EtaleWeightGaps.QAt}
\formal{Formalism.MainNotions.WeightAlgebra.TieC}
\formal{Formalism.MainNotions.WeightAlgebra.PageAt}
\formal{Formalism.MainNotions.Gaps.EtaleWeightGaps.PageSupp}
\end{definition}

\begin{definition}[Conditions on a global family datum]\label{pf:onto1:famconds}
In the situation of \cref{pf:onto1:hyps}, let $G$ be a global family datum for $(X_j,\kappa,w_L,w_F)$ and let
$G_c,P_H,P_{rw},P_{et},P_d\in\R$.
\begin{itemize}
\item[$(\mathrm G1)$] (generators) the $u_k$ span $W$ over $\OO_K$ and lie in $\mathcal M_{w_L}$, $\log N\le
P_{rw}(\ontoonesize w_L-1)$, $\log\|u_k\|_\sigma\le m_x(\sigma)$ for all $\sigma,k$, and $[K:\Q]^{-1}\sum_\sigma m_x(\sigma)\le
G_c\ontoonesize w_L$.
\item[$(\mathrm G2)$] $\theta'_i\in\mathrm{Th}_F$ and $\ontoonemls(\theta'_i)\le G_c\ontoonesize w_F$ for all $i$.
\item[$(\mathrm G3)$] (zero family) $\theta^0_k\in\mathrm{Th}_0\cap\mathcal M_{w_0}$, $\ontoonemls(\theta^0_k)\le G_c\ontoonesize w_0$,
$|\alpha_k|=a_0$, $\iota z_k(\beta)$ is the coefficient of $X^\beta$ in $z(\alpha_k,\theta^0_k)$, and $\det(\phi,z)\ne0$ for
some $\phi$.
\item[$(\mathrm G4)$] (quotient family) $r_q=\rho_P(C_{Rm},T,c)$, $\theta_i\in\mathrm{Th}_K$, $\sum K\theta_i+\mathcal
M^{\mathrm{gg}}_\kappa(C_{Rm},T,c)=\mathcal F_\kappa$ and $\ontoonemls(\theta_i)\le\ontoonePP(C_h,A_h)\ontoonesize\kappa$.
\item[$(\mathrm G5)$] (auxiliary form) if $j<m$: $i_\eta>0$, $\ontooneet\ne0$, $\iota\ontooneet$ is not a zero divisor on the
forms, and $\log\|\ontooneet\|_\sigma/i_\eta\le\ontoonePP(C_f,A_f)$ for all $\sigma$.
\item[$(\mathrm D1)$] (hermitian structure) for some $d_H$ there are $\sigma$-semilinear maps $e_\sigma\colon\mathcal F_\kappa\to
\C^{d_H}$ with $|e_\sigma F|=\|F\|_{2,\sigma}$ and $\langle e_{\bar\sigma}F,e_{\bar\sigma}F'\rangle=\overline{\langle e_\sigma
F,e_\sigma F'\rangle}$.
\item[$(\mathrm D2)$] (submultiplicativity) $\|H\|_\sigma\le\|F\|_\sigma\|F'\|_\sigma$ whenever $\iota H=\iota F\cdot\iota F'$.
\item[$(\mathrm D3)$] $\log\dim_K\mathcal F_\kappa\le P_d$, and $(\mathrm D4)$: $\ontoonePP(C_G,A_G)(1+\log\ontoonesize\kappa)\le P_d$.
\item[$(\mathrm{Sh})$] (shares) $(\mathrm{Sh}_\Delta)$: $\binom{T-1}\delta\Delta\le\ontoonerk_\kappa$; $(\mathrm{Sh}_z)$:
$r_z\le22\delta^2\ontoonerk_\kappa/T$; $(\mathrm{Sh}'_z)$: $r_z\le r_q+22\delta^2\ontoonerk_\kappa/T$; $(\mathrm{Sh}_q)$:
$r_q\le22\delta^2\ontoonerk_\kappa/T$.
\item[$(\mathrm P)$] (domination) explicit inequalities saying that the auxiliary constants are dominated by powers of
$\ontoonePP(C_2,A_2)$ and $R(C_R)$, e.g.\ $25(2G_c+6P_H+P_{rw})\le\ontoonePP(C_2,A_2)$, $\tfrac12\le P_{rw}\le P_H$ and
$C_{Rm}\ge2$.
\item[$(\mathrm{DB})$] (depth bound) for all $\mathfrak l$ and $y\in\ontoonePt(\mathfrak l)$: if $j<m$ then
$v_{\mathfrak l,y}(\iota\ontooneet)\ge0$ and, when this value is finite, $\ontoonedepth_{\mathfrak l,y}\le\nu(\mathfrak l)+\frac
c{2i_\eta}v_{\mathfrak l,y}(\iota\ontooneet)$; if $j\ge m$ then $\ontoonedepth_{\mathfrak l,y}\le\nu(\mathfrak l)$.
\item[$(\mathrm{Ht})$] (heights) for a family $Q=(Q_{\mathfrak l})$ of local point data: $I_{\mathfrak l}<\infty$ for all
$\mathfrak l$, and
\begin{equation}\label{pf:onto1:heights}
2\sum_{\mathfrak l}I_{\mathfrak l}\ell(\mathfrak l)\le\Delta\log\Delta+2\sum_i\ontoonemls(\theta'_i)
+\Bigl(\tfrac83\delta(\delta+4)P_{rw}+2(\delta+4)(\delta+1)(2G_c+6P_H)\Bigr)\Delta\,\ontoonesize w_L,
\end{equation}
and if $j<m$ then $i_\eta^{-1}\sum_{\mathfrak l}\mathrm{Net}_{\mathfrak l}\ell(\mathfrak l)\le c_0\bigl((\delta+1)(2G_c+6P_H)
+\tfrac43\delta P_{rw}\bigr)\Delta\,\ontoonesize w_L+\Delta P_{et}$, the sums being convergent.
\end{itemize}
We write $(\mathrm E'_{\mathfrak l})$ for the implication $(\mathrm G1)\wedge(\mathrm C_{\mathfrak l})\wedge(\mathrm
T_{\mathfrak l})\Rightarrow(\mathrm E_{\mathfrak l})$.
\formal{Formalism.MainNotions.Gaps.EtaleWeightGaps.GlobalFamilyData.TiedW}
\formal{Formalism.MainNotions.Gaps.EtaleWeightGaps.GlobalFamilyData.TiedTh}
\formal{Formalism.MainNotions.Gaps.EtaleWeightGaps.GlobalFamilyData.TiedZ}
\formal{Formalism.MainNotions.Gaps.EtaleWeightGaps.GlobalFamilyData.TiedQ}
\formal{Formalism.MainNotions.Gaps.EtaleWeightGaps.GlobalFamilyData.TiedEt}
\formal{Formalism.MainNotions.Gaps.EtaleWeightGaps.DeepHerm}
\formal{Formalism.MainNotions.Gaps.EtaleWeightGaps.DeepSubMul}
\formal{Formalism.MainNotions.Gaps.EtaleWeightGaps.DeepLogD}
\formal{Formalism.MainNotions.Gaps.EtaleWeightGaps.PageGd}
\formal{Formalism.MainNotions.Gaps.EtaleWeightGaps.ShareRd}
\formal{Formalism.MainNotions.Gaps.EtaleWeightGaps.ShareRho}
\formal{Formalism.MainNotions.Gaps.EtaleWeightGaps.ShareRhoT}
\formal{Formalism.MainNotions.Gaps.EtaleWeightGaps.ShareRhoP}
\formal{Formalism.MainNotions.Gaps.EtaleWeightGaps.PolyDom}
\formal{Formalism.MainNotions.Gaps.EtaleWeightGaps.CombBound}
\formal{Formalism.MainNotions.Gaps.EtaleWeightGaps.Heights}
\formal{Formalism.MainNotions.Gaps.EtaleWeightGaps.LocalPointData.InLiftC}
\end{definition}

\begin{proposition}[The global sum]\label{pf:onto1:rLa_tsum}
In the situation of \cref{pf:onto1:famconds}, assume $(\mathrm H)$, $(\mathrm L)$, $(\mathrm S)$, $(\mathrm A)$,
($j<m\Rightarrow(\mathrm Q)$), $(\mathrm G1)$, $(\mathrm G3)$, $(\mathrm G4)$, $(\mathrm D1)$--$(\mathrm D4)$, $(\mathrm P)$, and
$(\mathrm{Ht})$ for a family $Q$ of local point data. Then $r(\mathfrak l)\ell(\mathfrak l)$ is summable and
\begin{multline*}
\sum_{\mathfrak l}r(\mathfrak l)\ell(\mathfrak l)\le\ontoonebin\Bigl(\sum_{\mathfrak l}I_{\mathfrak l}\ell(\mathfrak l)+[j<m]
T\frac{\max(0,c)}{2i_\eta}\sum_{\mathfrak l}\mathrm{Net}_{\mathfrak l}\ell(\mathfrak l)\Bigr)\\
+r_z\,\Xi\,\ontoonesize\kappa+2r_q\ontoonePP(C_q,A_q)\ontoonesize\kappa,
\end{multline*}
where $\Xi=2G_c+P_{rw}+\ontoonePPz(C_L,A_L)(1+\log\ontoonesize\kappa)+\tfrac32(P_d+\mathrm{PO})$.
\formallabel{Formalism.MainNotions.Gaps.EtaleWeightGaps.rLa_tsum_le_comb}
\end{proposition}

\begin{proof}
Three bounds are combined.
\begin{enumerate}
\item (Gauss contents.) By \cref{pf:onto1:deepC} (proved in \cref{sec:pf:onto1:gausssum}), applied with
$p=\ontoonePPz(C_L,A_L)(1+\log\ontoonesize\kappa)$, which is possible by (H2), $(\mathrm G1)$, $(\mathrm G3)$, $(\mathrm D1)$--$(\mathrm
D4)$, $(\mathrm L)$, $(\mathrm A)$ and $(\mathrm S)$, and since $G_c,P_{rw}\ge0$ by $(\mathrm P)$, the family $\gamma_{\mathfrak
l}\ell(\mathfrak l)$ is summable with sum at most $r_z\Xi\ontoonesize\kappa$.
\item (Quotient lengths.) By \lref{Formalism.MainNotions.Gaps.EtaleWeightGaps.quotC}, which states that under
($j<m\Rightarrow(\mathrm Q)$), $(\mathrm H)$, $(\mathrm G4)$ and $(\mathrm P)$ the family $\lambda_{\mathfrak l}\ell(\mathfrak l)$ is
summable with sum at most $2r_q\ontoonePP(C_q,A_q)\ontoonesize\kappa$, the second bound holds. For $j<m$ this follows from
$(\mathrm Q)$ applied to the $\theta_i$ (its premises follow from $(\mathrm H)$ and $(\mathrm P)$), since
$\sum_{\mathfrak l}(n_n(\mathfrak l)+\rho_Pb_e(\mathfrak l))\ell(\mathfrak l)\le2\rho_P\ontoonePP(C_q,A_q)\ontoonesize\kappa$ and
$r_q=\rho_P$; for $j=m$ one has $q_m=0$, so $Q_\kappa$ has rank $0$, $r_q=0$ and all $\lambda_{\mathfrak l}$ vanish.
\item (Splitting.) Given the summability of $I_{\mathfrak l}\ell(\mathfrak l)$ and, for $j<m$, of
$\mathrm{Net}_{\mathfrak l}\ell(\mathfrak l)$ from $(\mathrm{Ht})$, and of the two families above, the family $r(\mathfrak
l)\ell(\mathfrak l)$ is summable and its sum is the corresponding linear combination of the four sums, by the
definition of $r(\mathfrak l)$.
\end{enumerate}
Adding the bounds of (1) and (2) to the identity of (3) gives the assertion.
\end{proof}

\begin{lemma}[Assembly of the index bound]\label{pf:onto1:concl_of_nat}
Let $X,S,\kappa,w_L,w_F,T$ be given, $\nu\ge0$, $\beta\in\R$, and $r\ge0$ a function on the primes such that
$r(\mathfrak l)\ell(\mathfrak l)$ is summable with sum $\le\beta$. Assume that for every $\mathfrak l$, every $e\in\N$ with
$T\nu(\mathfrak l)\le e$ and every $V'\subseteq K_{\mathfrak l}\ontooneMsup_\kappa$, the length
$\ontooneql_{\mathfrak l}(\ontooneMsup_{\kappa,\mathfrak l}\cap V',\Theta_{\mathfrak l}\cap V')$ is finite and $\le e\dim V'+r(\mathfrak
l)$. Then the index bound $\ontooneIB(\beta)$ holds.
\formallabel{Formalism.MainNotions.Gaps.EtaleWeightGaps.concl_of_nat}
\end{lemma}

\begin{proof}
Take the same $r$. In $\ontooneIB(\beta)$ the integer $e$ satisfies $e\ge T\nu(\mathfrak l)\ge0$, so $e\in\N$ and the
hypothesis applies.
\end{proof}

\begin{theorem}[Theorem Zero]\label{pf:onto1:thmzero}
In the situation of \cref{pf:onto1:famconds}, let $Q=(Q_{\mathfrak l})$ be a family of local point data for $G$. Assume
$(\mathrm H)$, $(\mathrm L)$, $(\mathrm S)$, $(\mathrm A)$, ($j<m\Rightarrow(\mathrm Q)$), $(\mathrm G1)$--$(\mathrm G5)$, $(\mathrm
D1)$--$(\mathrm D4)$, $(\mathrm{Sh}_\Delta)$, $(\mathrm{Sh}_z)$, $(\mathrm{Sh}_q)$, $(\mathrm P)$, $(\mathrm{DB})$, $(\mathrm{Ht})$, and
for every prime $\mathfrak l$: $(\mathrm B_{\mathfrak l})$, $(\mathrm C_{\mathfrak l})$, $(\mathrm T_{\mathfrak l})$, $(\mathrm E'_{\mathfrak
l})$, $(\mathrm Z_{\mathfrak l})$ and $(\mathrm M_{\mathfrak l})$ with $(C_{Rm},c)$. Then the index bound $\ontooneIB(\beta)$ holds
for $(X_j,S,\kappa,w_L,w_F,T,\nu)$ with
\[
\beta=\ontoonerk_\kappa\cdot f\cdot m\cdot b\cdot\delta_2 .
\]
\formallabel{Formalism.MainNotions.Gaps.EtaleWeightGaps.theorem_zero_of_named_inputs}
\end{theorem}

\begin{proof}
\emph{Global part.} By \cref{pf:onto1:rLa_tsum}, $r(\mathfrak l)\ell(\mathfrak l)$ is summable with sum bounded by the
expression displayed there. By \lref{Formalism.MainNotions.Gaps.EtaleWeightGaps.sumGC}, which states that under
$(\mathrm H)$, $(\mathrm G2)$, $(\mathrm D3)$, $(\mathrm{Sh}_\Delta)$, $(\mathrm{Sh}_z)$, $(\mathrm{Sh}_q)$, $(\mathrm P)$ and
$(\mathrm{Ht})$ that expression is at most $\ontoonerk_\kappa fmb\delta_2$ (a numerical estimate), the sum is at most
$\beta$.

\emph{Local part.} Fix $\mathfrak l$, $e\in\N$ with $T\nu(\mathfrak l)\le e$ and $V'$. We apply
\cref{pf:onto1:localBound} with $\nu_s=\nu(\mathfrak l)$, $C_R=C_{Rm}$ and the constant $c$, and check its hypotheses.
(a) $\mathcal M_\kappa$ is a lattice by $(\mathrm L)$. $\Theta\subseteq\ontooneMsup_\kappa$: by (H5) the modules $W$,
$\mathrm{Th}_F$, $\mathrm{Th}_0$, $\mathrm{Th}_K$ consist of integral forms; a product $\iota F\cdot\iota F'=\iota H$ of integral forms
has $v_{\mathfrak l,y}(\iota H)=v_{\mathfrak l,y}(\iota F)+v_{\mathfrak l,y}(\iota F')\ge0$, so $H$ is integral by the integrality
condition; and $Tw_L+w_F=(T-\delta)w_L+w_0=\kappa$ because $\delta\le T$ and $\kappa=w_F+Tw_L$ by (H2). Finally
$K\mathrm{Th}_K=\mathcal F_\kappa$ by (H3). (b) The $u_k$ lie in $W$ (they span it) and in $\mathcal M_{w_L}$; $\theta'_i\in\mathrm{Th}_F
\subseteq\mathcal M_{w_F}$ by $(\mathrm G2)$ and (H5); the conditions on $\theta^0_k$, $\alpha_k$, $z_k$ are $(\mathrm G3)$. (c) is
$(\mathrm G4)$. (d) is assumed. (e) By $(\mathrm{DB})$ at the points $y_\zeta$: if $j<m$, $v_{\mathfrak l,y_\zeta}(\iota\ontooneet)\ge0$,
it is finite by $(\mathrm T_{\mathfrak l})$, and $\nu(\zeta)\le\nu(\mathfrak l)+\frac c{2i_\eta}v_{\mathfrak l,y_\zeta}(\iota
\ontooneet)\le\nu(\mathfrak l)+\max(0,c)v_\eta(\zeta)$; if $j\ge m$, $v_\eta=0$ and $\nu(\zeta)\le\nu(\mathfrak l)$. (f) and (g)
are given by \cref{pf:onto1:lenBase} and \cref{pf:onto1:lenPole}: $(\mathrm B_{\mathfrak
l})$, $(\mathrm C_{\mathfrak l})$, $(\mathrm T_{\mathfrak l})$, $(\mathrm Z_{\mathfrak l})$, $(\mathrm M_{\mathfrak l})$ hold, $I_{\mathfrak
l}<\infty$ by $(\mathrm{Ht})$, and $(\mathrm E_{\mathfrak l})$ follows from $(\mathrm E'_{\mathfrak l})$, $(\mathrm G1)$, $(\mathrm
C_{\mathfrak l})$ and $(\mathrm T_{\mathfrak l})$. Hence the length in question is finite and at most $e\dim V'+r(\mathfrak l)$.

Finally $\nu\ge0$ by (H4), and $r\ge0$ since all its terms are non-negative ($\mathrm{Net}_{\mathfrak l}\ge0$ by
$(\mathrm{DB})$). \Cref{pf:onto1:concl_of_nat} concludes.
\end{proof}

The development also uses a variant in which the depth bound is only required at the points $y_\zeta$ of the
local point data, $(\mathrm{Sh}_z)$ is replaced by $(\mathrm{Sh}'_z)$, and the case $j\ge m$ is supplied as a hypothesis.

\begin{corollary}[Theorem Zero, refined form]\label{pf:onto1:thmzero-refined}
In the situation of \cref{pf:onto1:famconds}, assume $(\mathrm H)$, $(\mathrm L)$, $(\mathrm S)$, $(\mathrm A)$, ($j<m\Rightarrow
(\mathrm Q)$), $(\mathrm K)$, $(\mathrm{PA})$, $(\mathrm{PS})$, $(\mathrm G1)$--$(\mathrm G5)$, $(\mathrm D1)$--$(\mathrm D4)$,
$(\mathrm{Sh}_\Delta)$, $(\mathrm{Sh}'_z)$, $(\mathrm{Sh}_q)$, $(\mathrm P)$, and $(\mathrm Z_{\mathfrak l})$, $(\mathrm M_{\mathfrak l})$ for
all $\mathfrak l$; and:
\begin{enumerate}[label=(\roman*)]
\item if $j<m$: there is a family $Q$ of local point data satisfying $(\mathrm B_{\mathfrak l})$, $(\mathrm C_{\mathfrak l})$,
$(\mathrm T_{\mathfrak l})$, $(\mathrm E'_{\mathfrak l})$ for all $\mathfrak l$, $(\mathrm{Ht})$, and the depth bound at the points:
for all $\mathfrak l,\zeta$, $v_{\mathfrak l,y_\zeta}(\iota\ontooneet)\ge0$ and $\ontoonedepth_{\mathfrak l,y_\zeta}\le\nu(\mathfrak
l)+\frac c{2i_\eta}v_{\mathfrak l,y_\zeta}(\iota\ontooneet)$ when $v_{\mathfrak l,y_\zeta}(\iota\ontooneet)<\infty$;
\item if $j\ge m$: $\ontoonedepth_{\mathfrak l,y}\le\nu(\mathfrak l)$ for all $\mathfrak l$ and $y$;
\item if $j\ge m$ and the depth bound of (ii) holds, then $\ontooneIB(\ontoonerk_\kappa fmb\delta_2)$ holds.
\end{enumerate}
Then $\ontooneIB(\ontoonerk_\kappa fmb\delta_2)$ holds for $(X_j,S,\kappa,w_L,w_F,T,\nu)$.
\formallabel{Formalism.MainNotions.Gaps.EtaleWeightGaps.theorem_zero_of_refined_inputs}
\end{corollary}

\begin{proof}
If $j\ge m$, apply (iii) with (ii). If $j<m$, take $Q$ from (i). By \cref{pf:onto1:rLa_tsum}, $r(\mathfrak l)\ell(\mathfrak l)$
is summable with sum bounded as there, and by \lref{Formalism.MainNotions.Gaps.EtaleWeightGaps.sumGC4}, the
analogue of the numerical estimate used in the proof of \cref{pf:onto1:thmzero} with $(\mathrm{Sh}'_z)$ in place of
$(\mathrm{Sh}_z)$, this bound is at most $\ontoonerk_\kappa fmb\delta_2$. The local part is the same as in the proof of
\cref{pf:onto1:thmzero}: only the depth bound at the points $y_\zeta$ is used there, and (f), (g) of
\cref{pf:onto1:localBound} come from \cref{pf:onto1:lenBase} and \cref{pf:onto1:lenPole}. \Cref{pf:onto1:concl_of_nat} concludes.
\end{proof}

\begin{remark}\label{pf:onto1:rem-refined}
In \cref{pf:onto1:thmzero-refined} the hypotheses $(\mathrm K)$, $(\mathrm{PA})$, $(\mathrm{PS})$ and $(\mathrm G5)$ are not used, and the
case $j=m$ is not proved but assumed in (iii); in the application (iii) is supplied by \cref{pf:onto1:thmzero}
(see the proof of \cref{pf:onto1:inputs4_of_inputsR}).
\end{remark}

%----------------------------------------------------------------------------------------------------
\subsection{The sum of the Gauss contents}\label{sec:pf:onto1:gausssum}

In this subsection we prove the bound on the Gauss contents used in \cref{pf:onto1:rLa_tsum}. Let $K$ be a number
field, $n_K=[K:\Q]$, and $\mathfrak d_K^{-1}=\{x\in K:\Tr_{K/\Q}(x\OO_K)\subseteq\Z\}$ the codifferent. For a finite
dimensional $K$-space $V$, a \emph{hermitian structure} on $V$ is a family of $\sigma$-semilinear maps $e_\sigma\colon
V\to\C^{d_H}$ ($\sigma\colon K\to\C$); we write $N_\sigma(f)=|e_\sigma f|$ and say that the structure is \emph{conjugation
symmetric} if $\langle e_{\bar\sigma}f,e_{\bar\sigma}g\rangle=\overline{\langle e_\sigma f,e_\sigma g\rangle}$. A $K$-linear form $\phi$
on $V$ has \emph{$\sigma$-norm at most $C$} if there is $a\in\C^{d_H}$ with $|a|\le C$ and $\sigma(\phi(f))=\langle a,e_\sigma
f\rangle$ for all $f$.

\begin{lemma}[Dual forms from transference]\label{pf:onto1:transf}
Let $V$ be a $K$-space of dimension $d$ with a conjugation symmetric hermitian structure, and $M\subseteq V$ an
$\OO_K$-submodule spanning $V$. Let $\Lambda>0$ with $\sum_\sigma N_\sigma(f)^2\ge\Lambda^2$ for all $0\ne f\in M$, and let
$C>dn_K$. Then there are $K$-linear forms $\phi_1,\dots,\phi_n$ on $V$ spanning the dual space, with $\phi_j(M)\subseteq
\mathfrak d_K^{-1}$ and of $\sigma$-norm at most $C/\Lambda$ for all $j$ and $\sigma$.
\formallabel{Formalism.MainNotions.Gaps.SizeBoundsAlt.trans_of_minima}
\formallabel{Formalism.MainNotions.Gaps.SizeBoundsAlt.trans_of_stages}
\end{lemma}

\begin{proof}
Let $E$ be the real span of the image of $\iota\colon f\mapsto(e_\sigma f)_\sigma$ in $\bigoplus_\sigma\C^{d_H}$, with the real
inner product $\mathrm{Re}\sum_\sigma\langle\cdot,\cdot\rangle$, so that $|\iota f|^2=\sum_\sigma N_\sigma(f)^2$. Since every $f\in V$
has a non-zero integer multiple in $M$, $E$ is spanned by $L=\iota(M)$ and $\dim_\R E\le dn_K$; by hypothesis every non-zero
vector of $L$ has length $\ge\Lambda$ (if $\iota f=0$ for some $f\ne0$, a multiple of $f$ in $M$ would contradict the
hypothesis), so $L$ is a lattice in $E$ of first minimum $\ge\Lambda$. By the transference theorem for successive minima
(\lref{transference_successive_minima}, which states $\lambda_i(L)\lambda_{n+1-i}(L^*)\le n$ for a lattice $L$ of rank
$n$ and its dual lattice $L^*$; see \cite{Banaszczyk1993}), the dual lattice has $\dim E$ linearly independent vectors
$y_j$ with $|y_j|\le C/\Lambda$, and $\langle\iota f,y_j\rangle\in\Z$ for $f\in M$. The map $\ell_j(f)=\langle\iota f,y_j\rangle$ is
additive and $\Q$-valued on $V$ (again because multiples of $f$ lie in $M$), so by the non-degeneracy of the trace form
there is a $K$-linear form $\phi_j$ with $\Tr_{K/\Q}(\phi_j(f))=\ell_j(f)$ for all $f$. For $f\in M$ and $a\in\OO_K$,
$\Tr(a\phi_j(f))=\ell_j(af)\in\Z$, so $\phi_j(M)\subseteq\mathfrak d_K^{-1}$. The $\phi_j$ span the dual: if $\phi_j(f)=0$ for all
$j$, then $\iota f$ is orthogonal to the basis $(y_j)$ of $E$, hence $\iota f=0$ and $f=0$. Finally, writing $y_j=(y_{j,\sigma})_\sigma$,
the identity $\sum_\sigma\sigma(a\phi_j(f))=\Tr(\phi_j(af))=\mathrm{Re}\sum_\sigma\langle e_\sigma(af),y_{j,\sigma}\rangle$ for all $a\in K$,
together with the conjugation symmetry and the linear independence of the embeddings, gives $\sigma(\phi_j(f))=\langle
y_{j,\sigma},e_\sigma f\rangle$ with $|y_{j,\sigma}|\le|y_j|\le C/\Lambda$.
\end{proof}

\begin{lemma}[Coefficients of determinants]\label{pf:onto1:detcoef}
Let $V,M$ and the hermitian structure be as in \cref{pf:onto1:transf}, let $z_1,\dots,z_\rho$ be finitely supported maps from
the monomials in a set of variables to $V$, and $c\ge0$ a function on the primes of $K$ such that: whenever
$\phi_1,\dots,\phi_\rho$ are $K$-linear forms with $\ontooneord_{\mathfrak l}\phi_i(M)\ge0$ and a coefficient $a$ of
$\det(\phi,z)$ is non-zero, $c(\mathfrak l)\le\ontooneord_{\mathfrak l}(a)$. Assume $\det(\phi,z)\ne0$ for some $\phi$, and let
$\phi_1,\dots,\phi_n$ be forms spanning the dual, with $\phi_j(M)\subseteq\mathfrak d_K^{-1}$ and of $\sigma$-norm at most $B>0$ for all
$j,\sigma$. If $\sum_\gamma N_\sigma(z_k(\gamma))\le e^{E(\sigma,k)}$ for all $\sigma,k$, then $c(\mathfrak l)\log\Nm\mathfrak l$ has finite
support and
\[
\sum_{\mathfrak l}c(\mathfrak l)\log\Nm\mathfrak l\le\rho\log|d_K|+\sum_\sigma\Bigl(\rho\log B+\sum_kE(\sigma,k)\Bigr).
\]
\formallabel{Formalism.MainNotions.Gaps.SizeBoundsAlt.deepCore}
\formallabel{Formalism.MainNotions.Gaps.SizeBoundsAlt.core_arch}
\formallabel{Formalism.MainNotions.Gaps.SizeBoundsAlt.core_val}
\formallabel{Formalism.MainNotions.Gaps.SizeBoundsAlt.exists_minor_ne_zero}
\end{lemma}

\begin{proof}
The polynomial $\det(\phi,z)$ is multilinear in $(\phi_1,\dots,\phi_\rho)$; expanding a family with non-zero determinant in
terms of the spanning family gives indices $J(1),\dots,J(\rho)$ with $\det((\phi_{J(i)})_i,z)\ne0$. Let $a\ne0$ be one of its
coefficients, say that of $X^\beta$.

\emph{Finite places.} Let $\mathfrak l$ be a prime and choose $t\in K$ with $\ontooneord_{\mathfrak l}(t)=\ontooneord_{\mathfrak l}(\mathfrak
d_K)$; then the forms $t\phi_{J(i)}$ are $\mathfrak l$-integral on $M$, and the corresponding coefficient is $t^\rho a$, so
$c(\mathfrak l)\le\ontooneord_{\mathfrak l}(a)+\rho\,\ontooneord_{\mathfrak l}(\mathfrak d_K)$. Both terms vanish for almost all $\mathfrak
l$; by the product formula $\sum_{\mathfrak l}\ontooneord_{\mathfrak l}(a)\log\Nm\mathfrak l=\log|\Nm_{K/\Q}(a)|=\sum_\sigma\log|\sigma(a)|$,
and $\sum_{\mathfrak l}\ontooneord_{\mathfrak l}(\mathfrak d_K)\log\Nm\mathfrak l=\log\Nm\mathfrak d_K=\log|d_K|$.

\emph{Archimedean places.} Fix $\sigma$ and vectors $b_i$ with $|b_i|\le B$ and $\sigma\phi_{J(i)}=\langle b_i,e_\sigma\cdot\rangle$. The
coefficient $\sigma(a)$ is the sum, over the families $(\gamma_k)_k$ of monomials with $\sum_k\gamma_k=\beta$, of the determinants
$\det\bigl(\langle b_i,e_\sigma z_k(\gamma_k)\rangle\bigr)_{i,k}$. By Hadamard's inequality (applied to the Gram-type product of the
matrices with columns $b_i$ and $e_\sigma z_k(\gamma_k)$), each such determinant has absolute value at most
$\prod_i|b_i|\prod_kN_\sigma(z_k(\gamma_k))$. Summing over all families gives $|\sigma(a)|\le B^\rho\prod_k\sum_\gamma
N_\sigma(z_k(\gamma))\le B^\rho\exp(\sum_kE(\sigma,k))$. Combining the two parts gives the assertion.
\end{proof}

\begin{lemma}[Coefficients of products]\label{pf:onto1:l1coef}
Let $X$ be a valued form algebra over $B$ satisfying $(\mathrm D2)$, let $u_0,\dots,u_{N-1}\in\mathcal F_{w_L}$, $\theta\in\mathcal
F_{w_0}$ and $\alpha\in\N^\delta$, and let $z(\beta)\in\mathcal F_{|\alpha|w_L+w_0}$ be the forms with $\iota z(\beta)$ the coefficient of
$X^\beta$ in $z(\alpha,\theta)$. Then for every $\sigma$,
$\sum_\beta\|z(\beta)\|_\sigma\le\bigl(\sum_k\|u_k\|_\sigma\bigr)^{|\alpha|}\|\theta\|_\sigma$.
\formallabel{Formalism.MainNotions.Gaps.SizeBoundsAlt.polynomial_coefficient_l1_bound}
\end{lemma}

\begin{proof}
For a polynomial $p$ with coefficients $\iota F_\gamma$ ($F_\gamma$ of a fixed weight) put $\|p\|_1=\sum_\gamma\|F_\gamma\|_\sigma$. By
the triangle inequality and $(\mathrm D2)$, $\|pq\|_1\le\|p\|_1\|q\|_1$. The linear form $\sum_kX_{i,k}\iota(u_k)$ has
$\|\cdot\|_1=\sum_k\|u_k\|_\sigma$, and $z(\alpha,\theta)$ is the product of $|\alpha|$ such forms and of $\iota\theta$.
\end{proof}

\begin{theorem}[Sum of the Gauss contents]\label{pf:onto1:deepC}
In the situation of \cref{pf:onto1:famconds}, let $p,G_c,P_{rw},P_d\in\R$ with $G_c,P_{rw}\ge0$, and assume the weight conditions
of (H2), $(\mathrm G1)$, $(\mathrm G3)$, $(\mathrm D1)$--$(\mathrm D4)$, $(\mathrm L)$, $(\mathrm S)$, and
$\ontoonemls(F)\ge-p\,\ontoonesize\kappa$ for all $0\ne F\in\mathcal M_\kappa$. Then $\gamma_{\mathfrak l}\ell(\mathfrak l)$ is summable and
\[
\sum_{\mathfrak l}\gamma_{\mathfrak l}\,\ell(\mathfrak l)\le r_z\Bigl(2G_c+P_{rw}+p+\tfrac32\bigl(P_d+\mathrm{PO}\bigr)\Bigr)\ontoonesize\kappa .
\]
\formallabel{Formalism.MainNotions.Gaps.EtaleWeightGaps.deepC}
\formallabel{Formalism.MainNotions.Gaps.EtaleWeightGaps.deepPacketP}
\end{theorem}

\begin{proof}
Let $d=\dim_K\mathcal F_\kappa$ and $n_K=[K:\Q]$. By (H2), $\kappa=a_0w_L+w_0$, $a_0\le\ontoonesize\kappa$, $\log n_K\le\mathrm{PO}$ and
$\log|d_K|/n_K\le\mathrm{PO}$; moreover $\ontoonesize$ is affine, so $\ontoonesize\kappa=a_0(\ontoonesize w_L-1)+\ontoonesize w_0$.

\emph{Step 1: integrality.} As in the proof of \cref{pf:onto1:gaussC}, $(\mathrm G1)$ and $(\mathrm G3)$ imply $z_k(\beta)\in\mathcal M_\kappa$
for all $k,\beta$. By the definition of $\gamma_{\mathfrak l}$ as a least valuation, the function $c(\mathfrak l)=\gamma_{\mathfrak l}$
satisfies the hypothesis of \cref{pf:onto1:detcoef} with $M=\mathcal M_\kappa$, and $\det(\phi,z)\ne0$ for some $\phi$ by $(\mathrm G3)$.

\emph{Step 2: sizes of coefficients.} By \cref{pf:onto1:l1coef}, $N_\sigma\le\|\cdot\|_\sigma$ and $(\mathrm G1)$,
$\sum_\beta\|z_k(\beta)\|_{2,\sigma}\le(Ne^{m_x(\sigma)})^{a_0}\|\theta^0_k\|_\sigma\le e^{E(\sigma,k)}$ with $E(\sigma,k)=a_0(\log N+m_x(\sigma))+\log
\|\theta^0_k\|_\sigma$.

\emph{Step 3: a lower bound on the lattice.} Let $0\ne F\in\mathcal M_\kappa$ and $s_\sigma=\|F\|_\sigma$. By the arithmetic-geometric
mean inequality and the hypothesis on $\ontoonemls$, $\sum_\sigma s_\sigma^2\ge n_K\exp\bigl(\tfrac2{n_K}\sum_\sigma\log s_\sigma\bigr)\ge
n_Ke^{-2p\,\ontoonesize\kappa}$, and by $(\mathrm S)$ and $(\mathrm D4)$, $s_\sigma^2\le e^{P_d}\|F\|_{2,\sigma}^2$. Hence $\sum_\sigma\|F\|_{2,\sigma}^2\ge
\Lambda^2$ with $\Lambda=\sqrt{n_K}\exp\bigl(-(p\,\ontoonesize\kappa+P_d/2)\bigr)$.

\emph{Step 4: transference and product formula.} Let $C_0=\max(dn_K,1)$ and $C>C_0$. The hermitian structure of $(\mathrm D1)$
is conjugation symmetric, and $\mathcal M_\kappa$ spans $\mathcal F_\kappa$ by $(\mathrm L)$, so \cref{pf:onto1:transf} gives forms of
$\sigma$-norm at most $C/\Lambda$, and \cref{pf:onto1:detcoef} with $B=C/\Lambda$ gives
\[
\sum_{\mathfrak l}\gamma_{\mathfrak l}\log\Nm\mathfrak l\le r_z\log|d_K|+r_zn_K\bigl(\log C-\log\Lambda\bigr)+\sum_k\sum_\sigma E(\sigma,k).
\]
Since $C>C_0$ is arbitrary, the same holds with $C_0$ in place of $C$. Dividing by $n_K$ and inserting $\log\Lambda$,
\begin{multline*}
\sum_{\mathfrak l}\gamma_{\mathfrak l}\ell(\mathfrak l)\le r_z\Bigl(\tfrac{\log|d_K|}{n_K}+\log C_0-\tfrac12\log n_K+p\,\ontoonesize\kappa
+\tfrac{P_d}2\Bigr)\\+\sum_k\Bigl(a_0\log N+\tfrac{a_0}{n_K}\sum_\sigma m_x(\sigma)+\ontoonemls(\theta^0_k)\Bigr).
\end{multline*}

\emph{Step 5: numerical estimate.} We have $\log C_0-\frac12\log n_K\le\log d+\frac12\log n_K\le P_d+\frac12\mathrm{PO}$ by
$(\mathrm D3)$ (if $d=0$ the left side is $\le0$), so the first bracket is at most $\tfrac32(P_d+\mathrm{PO})+p\,\ontoonesize\kappa$. By $(\mathrm G1)$ and $(\mathrm G3)$ each term of the second sum is at most
$a_0P_{rw}(\ontoonesize w_L-1)+G_c(a_0\ontoonesize w_L+\ontoonesize w_0)\le P_{rw}\ontoonesize\kappa+2G_c\ontoonesize\kappa$, using the
affine formula for $\ontoonesize\kappa$ and $a_0\le\ontoonesize\kappa$. Since $\ontoonesize\kappa\ge1$ and $P_d,\mathrm{PO}\ge0$, the
assertion follows; summability holds because $\gamma_{\mathfrak l}$ has finite support (\cref{pf:onto1:detcoef}).
\end{proof}

%----------------------------------------------------------------------------------------------------
\subsection{Second norms from finite nets, and the conversion of the inputs}\label{sec:pf:onto1:net}

\begin{definition}[Second norms]\label{pf:onto1:secnorm}
A \emph{second norm} $n$ for a form provider $B$ is a family of functions $N_\sigma\colon\mathcal F_w\to\R$ (for all
$j$, all weights $w$ and all embeddings $\sigma$) satisfying the axioms of \cref{pf:onto1:prov}(iii) for the
$L^2$-norm: $N_\sigma\ge0$, $N_\sigma(F)=0$ iff $F=0$, $N_\sigma(aF)=|\sigma a|N_\sigma(F)$, $N_\sigma(F+G)\le
N_\sigma(F)+N_\sigma(G)$ and $N_\sigma\le\|\cdot\|_\sigma$. We write $B^n$ for the form provider obtained from $B$ by
replacing $\|\cdot\|_{2,\sigma}$ by $N_\sigma$, and $X^n$, $S^n$ for a valued form algebra over $B$ and a scalar quadruple
regarded over $B^n$; all other data are unchanged.
\formal{GlobalTheory.RunIn.SecNorm}
\formal{Formalism.L2.Pj.PJ.tol2NormAlt}
\formal{Formalism.L2.Pj.scalarsWithNorm}
\end{definition}

\begin{definition}[Nets]\label{pf:onto1:net}
A \emph{net} for $(B,j,\kappa)$ consists of, for each embedding $\sigma$, a finite set $\mathcal N_\sigma$ and
$\sigma$-semilinear functionals $\mathrm{ev}_{\sigma,y}\colon\mathcal F_\kappa\to\C$ ($y\in\mathcal N_\sigma$), together with
bijections $c_\sigma\colon\mathcal N_\sigma\to\mathcal N_{\bar\sigma}$ such that $\mathrm{ev}_{\bar\sigma,c_\sigma(y)}(F)=
\overline{\mathrm{ev}_{\sigma,y}(F)}$. Its \emph{net norm} is $N^{\mathrm{net}}_\sigma(F)=\bigl(|\mathcal N_\sigma|^{-1}\sum_{y}
|\mathrm{ev}_{\sigma,y}(F)|^2\bigr)^{1/2}$. The \emph{net condition} $(\mathrm N)$ with constant $P_d$ asks for a net and an
embedding $\sigma_0$ such that:
\begin{itemize}
\item for all $\sigma,F$ some $y$ has $\|F\|_\sigma/4\le|\mathrm{ev}_{\sigma,y}(F)|$, and $|\mathrm{ev}_{\sigma,y}(F)|\le\|F\|_\sigma$ for
all $y$;
\item $\log(16|\mathcal N_\sigma|)\le\ontoonePP(C_G,A_G)(1+\log\ontoonesize\kappa)$ for all $\sigma$;
\item every $K$-linearly independent family $F_1,\dots,F_n$ in $\mathcal F_\kappa$ gives $\C$-linearly independent vectors
$(\mathrm{ev}_{\sigma_0,y}(F_i))_y$, and $\log|\mathcal N_{\sigma_0}|\le P_d$.
\end{itemize}
\formal{Formalism.MainNotions.Gaps.NetNorms.NetData}
\formal{Formalism.MainNotions.Gaps.NetNorms.NetData.netN}
\formal{GlobalTheory.FinalJoin.NetIn}
\end{definition}

\begin{lemma}[Second norms from nets]\label{pf:onto1:secNorm_of_netIn}
Assume the net condition $(\mathrm N)$ with constant $P_d$ for $(B,j,\kappa)$, and let $X$ be a valued form algebra over
$B$ at $j$. Then there is a second norm $n$ such that $B^n$ satisfies $(\mathrm S)$ at $\kappa$ with the constants $C_G,A_G$,
and $X^n$ satisfies $(\mathrm D1)$ and $(\mathrm D3)$ at $\kappa$.
\formallabel{GlobalTheory.FinalJoin.secNorm_of_netIn}
\end{lemma}

\begin{proof}
Let $N_\sigma$ be the net norm on $\mathcal F_\kappa$ (at the index $j$) and the sup norm on all other spaces. The axioms are
clear except definiteness, which follows from the first item of $(\mathrm N)$, and $N_\sigma\le\|\cdot\|_\sigma$, which
follows from $|\mathrm{ev}_{\sigma,y}(F)|\le\|F\|_\sigma$.

$(\mathrm S)$: given $F,\sigma$ choose $y$ with $\|F\|_\sigma\le4|\mathrm{ev}_{\sigma,y}(F)|$; then
$\|F\|_\sigma^2\le16|\mathrm{ev}_{\sigma,y}(F)|^2\le16\sum_{y'}|\mathrm{ev}_{\sigma,y'}(F)|^2=16|\mathcal N_\sigma|\,N_\sigma(F)^2$, and
$16|\mathcal N_\sigma|\le\exp(\ontoonePP(C_G,A_G)(1+\log\ontoonesize\kappa))$ by the second item.

$(\mathrm D1)$: let $d_H=\sum_\tau|\mathcal N_\tau|$ and $e_\sigma(F)\in\C^{d_H}$ the vector whose $(\tau,y)$-coordinate is
$|\mathcal N_\sigma|^{-1/2}\mathrm{ev}_{\sigma,y}(F)$ if $\tau=\sigma$ and $0$ otherwise. Then $e_\sigma$ is $\sigma$-semilinear,
$|e_\sigma F|=N_\sigma(F)$, and the conjugation relation of the net gives
$\langle e_{\bar\sigma}F,e_{\bar\sigma}F'\rangle=\overline{\langle e_\sigma F,e_\sigma F'\rangle}$.

$(\mathrm D3)$: applying the last item to a $K$-basis of $\mathcal F_\kappa$ gives $\dim_K\mathcal F_\kappa\le|\mathcal N_{\sigma_0}|$, so
$\log\dim_K\mathcal F_\kappa\le\log|\mathcal N_{\sigma_0}|\le P_d$ (if $\dim\mathcal F_\kappa=0$ the left side is $0\le
\log|\mathcal N_{\sigma_0}|$).
\end{proof}

\begin{definition}[Input packages]\label{pf:onto1:packages}
In the situation of \cref{pf:onto1:hyps}, define:
\begin{itemize}
\item the \emph{base package} $\Pi_{\mathrm{base}}$: $(\mathrm H)$, $(\mathrm L)$, $(\mathrm A)$, ($j<m\Rightarrow(\mathrm Q)$) and
$(\mathrm D2)$;
\item the \emph{family package} $\Pi_{\mathrm{fam}}(G;G_c,P_H,P_{rw},P_{et},P_d)$: $(\mathrm G1)$, $(\mathrm G2)$, $(\mathrm G5)$,
$(\mathrm D4)$, the growth bound $G_c\ge\mathrm{Gc}(p)$ for $p\in\{p_{\mathrm{de}},p'_{\mathrm{de}}\}$, where $\mathrm{Gc}(p)$ is an
explicit expression in $L$, $K$, $E$ and $p$, the bound $G_c\le\ontoonePP(C_h,A_h)$, $(\mathrm{Sh}_\Delta)$, the share
condition $(\mathrm{Sh}_R)$ (for every $\mathfrak l$ the $K'$-span of $\{s^\alpha\mathrm{emb}\,\iota\theta:|\alpha|=a_0,\theta\in\mathrm{Th}_0\}$
has dimension $\le22\delta^2\ontoonerk_\kappa/T$), $(\mathrm{Sh}_q)$ for $\rho_P$ in place of $r_q$, $(\mathrm P)$, and
$(\mathrm B_{\mathfrak l})$, $(\mathrm M'_{\mathfrak l})$ for all $\mathfrak l$;
\item the \emph{point package} $\Pi_{\mathrm{pt}}(G;G_c,P_H,P_{rw},P_{et})$: $(\mathrm{DB})$, and there is a family $Q$ of local
point data with $(\mathrm C_{\mathfrak l})$, $(\mathrm T_{\mathfrak l})$, $(\mathrm E'_{\mathfrak l})$ for all $\mathfrak l$ and
$(\mathrm{Ht})$.
\end{itemize}
The three input packages are:
\begin{itemize}
\item $\Pi_{\mathrm{net}}$: $\Pi_{\mathrm{base}}$, and there are $G,G_c,P_H,P_{rw},P_{et},P_d$ with $(\mathrm N)$ (constant
$P_d$), $\Pi_{\mathrm{fam}}$ and $\Pi_{\mathrm{pt}}$;
\item $\Pi_{\mathrm{red}}$: $\Pi_{\mathrm{base}}$, $(\mathrm S)$, $(\mathrm D1)$, and there are $G,G_c,P_H,P_{rw},P_{et},P_d$ with
$(\mathrm D3)$, $\Pi_{\mathrm{fam}}$ and $\Pi_{\mathrm{pt}}$;
\item $\Pi_{\mathrm{full}}$: the hypotheses of \cref{pf:onto1:thmzero-refined}, with existential quantifiers over the
global family datum, the auxiliary constants and the families of local point data, together with the existence, when
$j\ge m$, of a family $Q$ satisfying $(\mathrm B_{\mathfrak l})$, $(\mathrm C_{\mathfrak l})$, $(\mathrm T_{\mathfrak l})$ and
$(\mathrm E'_{\mathfrak l})$ for all $\mathfrak l$, and $(\mathrm{Ht})$.
\end{itemize}
\formal{GlobalTheory.FinalJoin.BaseIn}
\formal{GlobalTheory.FinalJoin.FamIn}
\formal{GlobalTheory.FinalJoin.PtIn}
\formal{GlobalTheory.FinalJoin.InputsNet}
\formal{GlobalTheory.FinalJoin.InputsR}
\formal{GlobalTheory.FinalJoin.Inputs4}
\formal{Formalism.MainNotions.Gaps.EtaleWeightGaps.ShareRhoR}
\formal{Formalism.MainNotions.Gaps.EtaleWeightGaps.GrowthConstantBound}
\end{definition}

\begin{lemma}[Index bound from the full package]\label{pf:onto1:concl_of_inputs4}
$\Pi_{\mathrm{full}}$ implies the index bound $\ontooneIB(\ontoonerk_\kappa fmb\delta_2)$ for $(X_j,S,\kappa,w_L,w_F,T,\nu)$.
\formallabel{GlobalTheory.FinalJoin.concl_of_inputs4}
\end{lemma}

\begin{proof}
This is \cref{pf:onto1:thmzero-refined}.
\end{proof}

\begin{proposition}[From the reduced to the full package]\label{pf:onto1:inputs4_of_inputsR}
$\Pi_{\mathrm{red}}$, together with $(\mathrm K)$, $(\mathrm{PA})$ and $(\mathrm{PS})$, implies $\Pi_{\mathrm{full}}$.
\formallabel{GlobalTheory.FinalJoin.inputs4_of_inputsR}
\formallabel{GlobalTheory.T.indexBoundChosenFamily_of_intrinsicHypotheses}
\end{proposition}

\begin{proof}
Take $G,G_c,P_H,P_{rw},P_{et},P_d$ from $\Pi_{\mathrm{red}}$. The global family datum $G$ need not satisfy $(\mathrm G3)$ and
$(\mathrm G4)$; we replace its components $(r_z,\alpha,\theta^0,z)$ and $(r_q,\theta)$.
\begin{enumerate}
\item By (H3) and the growth bound, $\mathrm{Th}_0$ and $\mathrm{Th}_K$ contain finite sets spanning $\mathcal F_{w_0}$, resp.\
$\mathcal F_\kappa$, over $K$ and consisting of forms of mean log sup at most $G_c\ontoonesize w_0$, resp.\
$G_c\ontoonesize\kappa$. Moreover $a_0w_L+w_0=\kappa$ since $\delta\le T$.
\item By \lref{Formalism.MainNotions.Gaps.EtaleWeightGaps.exists_global_family_zero_choice}, which states that
from (1) and $\mathrm{Th}_0\subseteq\mathcal M_{w_0}$ one obtains $r_z,\alpha,\theta^0,z$ such that the modified datum satisfies
$(\mathrm G3)$ and $(\mathrm Z_{\mathfrak l})$ for all $\mathfrak l$, we make this replacement.
\item By \lref{Formalism.MainNotions.Gaps.EtaleWeightGaps.exists_global_family_quotient_bounded}, which states
that from (1), $\mathrm{Th}_K\subseteq\mathcal M_\kappa$ and $G_c\ontoonesize\kappa\le\ontoonePP(C_h,A_h)\ontoonesize\kappa$ one obtains
$\theta_1,\dots,\theta_{\rho_P}$ satisfying $(\mathrm G4)$, we obtain a datum $G''$ satisfying $(\mathrm G3)$ and $(\mathrm G4)$.
\item By \lref{GlobalTheory.shareRho_and_shareRhoT_of_shareRhoR}, which states that $(\mathrm{Sh}_R)$, $(\mathrm G3)$ and
$(\mathrm Z_{\mathfrak l})$ for all $\mathfrak l$ imply $(\mathrm{Sh}_z)$ and $(\mathrm{Sh}'_z)$ (the elements
$s^{\alpha_k}\mathrm{emb}\,\iota\theta^0_k$ then form a basis of the span in $(\mathrm{Sh}_R)$, so $r_z$ is its dimension), $G''$
satisfies $(\mathrm{Sh}_z)$ and $(\mathrm{Sh}'_z)$. Moreover $(\mathrm M'_{\mathfrak l})$ and $(\mathrm Z_{\mathfrak l})$ give $(\mathrm
M_{\mathfrak l})$.
\item The other conditions on $G$ ($(\mathrm G1)$, $(\mathrm G2)$, $(\mathrm G5)$, $(\mathrm B_{\mathfrak l})$, $(\mathrm{DB})$), the
local point data and the conditions on them do not involve the replaced components, and are transported to $G''$.
\item $(\mathrm{DB})$ gives the depth bound at the points required in (i) of \cref{pf:onto1:thmzero-refined} when $j<m$,
and the depth bound (ii) when $j\ge m$. Finally (iii) holds because \cref{pf:onto1:thmzero} applies to $G''$: all its
hypotheses have been established.
\end{enumerate}
Collecting these facts gives $\Pi_{\mathrm{full}}$.
\end{proof}

\begin{proposition}[From net inputs to the index bound]\label{pf:onto1:inputs4_of_net}
Let $X=(X_j)$ be a family of valued form algebras over $B$, $j$ in range, $C_L\in\R$, and put
\[
\mathbf K_R(C_L)=(1041,3,200000,2,1,2,C_L,2,7),\qquad \mathbf K_4(C_L)=(1041z_B,2,200000,2,1,2,C_L,2,7).
\]
Assume $\Pi_{\mathrm{net}}$ for $X_j$ with $\mathbf K=\mathbf K_R(C_L)$, $(C_h,A_h,A_q,C_R)=(1,4,4,7)$ and constants
$C_q,C_2,A_2,C_T$, and assume $C_1=C_c=C_2$, $A_1=A_c=A_2$, $C_{Rc}=7$, $12\le A_2$, $(\mathrm{PA})$ and $(\mathrm{PS})$. Then
there is a second norm $n$ such that $\Pi_{\mathrm{full}}$ holds for $X_j^n$ and $S^n$ with $\mathbf K=\mathbf K_4(C_L)$ and
the same other constants, and the index bound $\ontooneIB(\ontoonerk_\kappa fmb\delta_2)$ holds for
$(X_j^n,S^n,\kappa,w_L,w_F,T,\nu)$.
\formallabel{GlobalTheory.FinalJoin.inputs4_of_net}
\end{proposition}

\begin{proof}
\emph{Step 1.} From $\Pi_{\mathrm{net}}$ we obtain $(\mathrm N)$, and \cref{pf:onto1:secNorm_of_netIn} gives a second norm $n$
with $(\mathrm S)$, $(\mathrm D1)$ and $(\mathrm D3)$ for $X_j^n$. The packages $\Pi_{\mathrm{base}}$, $\Pi_{\mathrm{fam}}$ and
$\Pi_{\mathrm{pt}}$ do not involve the $L^2$-norm; they are transported from $(X_j,S)$ to $(X_j^n,S^n)$, the global family
datum and the local point data being transported along
(\lref{GlobalTheory.FinalJoin.baseIn_secNorm}, which states the transport of $\Pi_{\mathrm{base}}$, and
\lref{GlobalTheory.T.exists_formFamilyHypotheses_changedSecondNorm}, which states that of $\Pi_{\mathrm{fam}}$ and
$\Pi_{\mathrm{pt}}$). Hence $\Pi_{\mathrm{red}}$ holds for $X_j^n$ with $\mathbf K_R(C_L)$.

\emph{Step 2.} $\mathbf K_R(C_L)$ and $\mathbf K_4(C_L)$ differ only in $(C_G,A_G)$, and
$\ontoonePP(1041,3)=1041z_B^3=\ontoonePP(1041z_B,2)$; hence $\Pi_{\mathrm{red}}$ holds with $\mathbf K_4(C_L)$.

\emph{Step 3.} With the given equalities, $(\mathrm K)$ holds for $\mathbf K_4(C_L)$, $(C_h,A_h,A_q)=(1,4,4)$ and $C_R=7$:
indeed $C_{Rm}=7$, $A_G=A_f=A_W=2$, $A_L=2\le4$, $1\le1$, $4\le4$ and $4\le11$. \Cref{pf:onto1:inputs4_of_inputsR} gives
$\Pi_{\mathrm{full}}$, and \cref{pf:onto1:concl_of_inputs4} the index bound.
\end{proof}

%----------------------------------------------------------------------------------------------------
\subsection{The record of a run and the index bound along runs}\label{sec:pf:onto1:record}

In this subsection and the next, $T$ denotes a top structure; the integer exponent of the previous subsections is
denoted $T_n$.

\begin{definition}[The record of a top structure]\label{pf:onto1:record-def}
Let $T$ be a top structure over $D$ with laws $h$: a family of moduli data of unitary type with a common level and
coefficient field, together with theta rules, short theta rules and auxiliary shortening data, subject to the laws of
the development. To $(T,h)$ the development attaches a form provider $B_T$ whose forms are the cusp forms of the
moduli data, and for $j$ in range a valued form algebra $X_{T,j}$ over $B_T$ at $j$ whose points at $\mathfrak l$ are
integral points of the moduli spaces and whose valuations are the corresponding orders of forms. Further:
\begin{enumerate}[label=(\roman*)]
\item An \emph{enlargement} $E$ of a valued form algebra $X$ is a family of additional points with valuations
satisfying the axioms of \cref{pf:onto1:PJ}(ii), at which integral forms have non-negative valuation; $X\oplus E$ is
the valued form algebra with points $\ontoonePt(\mathfrak l)\sqcup E(\mathfrak l)$ (the integrality condition still holds).
We put $X^+_{T,j}=X_{T,j}\oplus E_{\mathrm{all}}$, where $E_{\mathrm{all}}$ consists of the valued points of the moduli spaces
satisfying the trivial class condition, and for a second norm $n$ of $B_T$ we write $X^{+,n}_{T,j}$ for $(X^+_{T,j})^n$. A
global family datum for $X$ is regarded as one for $X\oplus E$.
\item For an error datum $\mathrm{EX}$, a prime $p$, weights $\kappa,w_L,w_F$ and $T_n\in\N$, the scalar quadruple
$S_{\mathrm{rec}}=\mathrm{scal}(B_T,\mathrm{EX},p,\kappa,w_L,w_F,T_n)$ consists of explicit products of norms of elements of
$\mathrm{EX}$ and of powers of the auxiliary primes; its first entry $s_W=s(w_L)$ depends only on $w_L$. Its definition uses
the injectivity of the map $\OO\to\mathrm{LL}$ of the datum.
\item The function $\nu_T$ on the primes of $K$: at a prime $\mathfrak l$ dividing $g_i$ it is
$e_{\mathfrak l}\ontooneDF(w_L)/(g_i-1)$, and at other primes it is $e_{\mathfrak l}$ times the sum of an explicit local weight
of $(T,w_L)$ and a correction at the primes dividing $Q_S$, $q$ or $d$.
\end{enumerate}
\formal{Formalism.L2.Pj.Top}
\formal{Formalism.L2.Pj.LawsP}
\formal{Formalism.L2.Pj.PtExt}
\formal{Formalism.L2.Pj.IntAtExt}
\formal{Formalism.MainNotions.DIV.Cons.scalOfP}
\formal{GlobalTheory.RunJoin.nusR}
\end{definition}

The objects of \cref{pf:onto1:record-def} are constructed in other parts of the development:
$B_T$ in \lref{Formalism.L2.Pj.recB}, $X_{T,j}$ in \lref{Formalism.L2.Pj.recPJ}, the enlargement $X\oplus E$ in
\lref{Formalism.L2.Pj.xplus}, $E_{\mathrm{all}}$ and $X^+_{T,j}$ in \lref{Formalism.L2.Pj.extAll} and
\lref{Formalism.L2.Pj.recPlus}, and $X^{+,n}_{T,j}$ in \lref{GlobalTheory.RunIn.recPlusN}. For a second norm
$n$ one has $\mathrm{scal}(B_T^n,\mathrm{EX},p,\dots)=(S_{\mathrm{rec}})^n$ by definition.

\begin{definition}[Parameter and page conditions]\label{pf:onto1:parconds}
Let $T,h,\mathrm{EX},p$ be as above and $j$ in range.
\begin{enumerate}[label=(\roman*)]
\item The \emph{parameter conditions} $\mathcal C_{\mathrm{par}}(T,\mathrm{EX},p,j)$ are: the laws of the shortening data hold;
their levels are trivial; the components of the unfolding at $j$ share their group and character data; the bad integer
$N_{\mathrm{bad}}(D)=Q_S\,q\,\Nm(\mathfrak c_K)$ of the shortening datum $D$ (a product of two positive integers attached to its
host and of the norm of an ideal of the host field) is positive; all shortening data at $j$ have prime $p$; $T$ is
\emph{oriented} for $\mathrm{EX}$,
i.e.\ the level ideal $M_b$ is contained in the principal ideals generated by the conjugates $\bar\pi_i$ of the elements
$\pi_i$ of $\mathrm{EX}$; and the coefficient field is algebraically closed in the function fields of the moduli spaces at
$j$.
\item For $L$, a valued form algebra $X$, $\mathbf H$, $C_2,A_2,C_T$, weights $\kappa,w_L,w_F$, $T_n,n_P,n_W\in\N$ and
$h_0\in\R$: the \emph{page hypotheses} $\mathcal H_{\mathrm{page}}$ are (H1), (H2), the numerical parts of (H6) and (H7) of
$(\mathrm H)$ with $C_R=7$; the \emph{page conditions} $\mathcal C_{\mathrm{page}}$ are $\mathcal H_{\mathrm{page}}$, $(\mathrm{PA})$,
$C_1=C_c=C_2$, $A_1=A_c=A_2$, $C_{Rc}=7$ and $12\le A_2$; and the \emph{strengthened page conditions} $\mathcal
C^+_{\mathrm{page}}$ are $\mathcal C_{\mathrm{page}}$ together with the first, third, fourth and fifth items of $(\mathrm{PS})$ and
$4C_K\le\ontoonePP(C_2,A_2)^2f$.
\end{enumerate}
\formal{GlobalTheory.RunJoin.ParOK}
\formal{GlobalTheory.PtTie.Oriented}
\formal{Formalism.L2.Pj.FnAlgClosed}
\formal{Formalism.MainNotions.Gaps.FamilyHypotheses.HypPage}
\formal{GlobalTheory.FinalJoin.PageRun}
\formal{GlobalTheory.RunJoin.PageOK}
\end{definition}

\begin{definition}[Admissible generating families]\label{pf:onto1:genok-def}
Let $T,h$, $j$ in range, $L$ and weights be as above, $S$ a scalar quadruple for $B_T$ and $c\in\R$. Put
\[
G^{\max}=\max\bigl(1,\mathrm{Gc}(p_{\mathrm{de}}),\mathrm{Gc}(p'_{\mathrm{de}})\bigr),\qquad
\mathrm{nb}(w)=\log2+frn_w\bigl(\log(g_0g_1)+\log(1+\ontooneDF(w))\bigr)+\mathrm{PO}
\]
with $\mathrm{Gc}$ as in \cref{pf:onto1:packages}. A global family datum $G$ for $(X_{T,j},\kappa,w_L,w_F)$ is an
\emph{admissible generating family} for $(S,c)$ if: (i) $(\mathrm G1)$ holds with $(G_c,P_{rw})=(G^{\max},38\mathrm
L+\mathrm{PO})$; (ii) $N>0$ and $\log N\le\mathrm{nb}(w_L)$; (iii) the $u_k$ have no common zero: for every algebraically
closed field $k$ over the coefficient field and every $k$-point of the moduli space at $j$, some $u_k$ does not vanish
there; and (iv) the \emph{lower bound at frames}: for every good embedding $\sigma$ of the coefficient field, every
component of the unfolding at $\sigma$ and every frame $x$ in its domain,
\[
-\log\Bigl(\sum_{k}\|u_k\|_{x}^2\Bigr)^{1/2}\le p_{14}(w_L)\,\ontooneDF(w_L)+c\,\phi_{FD}(x),
\]
where $\|\cdot\|_x$ is the normalised frame norm at $x$, $\phi_{FD}$ an explicit function of frames and $p_{14}(w_L)\ge1$ an
explicit constant. The \emph{auxiliary form condition} for $(i,\ontooneet)$ with $\ontooneet\in\mathcal F_{\epsilon_i}$ is: $i>0$;
$\ontooneet\ne0$ and $\iota\ontooneet$ is not a zero divisor on the forms; $v_{\mathfrak l,y}(\iota\ontooneet)\ge2i\phi_C(\mathfrak l,y)$
for all $\mathfrak l,y$; $\log\|\ontooneet\|_\sigma/i\le\ontoonePP(200000,2)$ for all $\sigma$; at every frame $x$ as in (iv)
the normalised frame norm of $\ontooneet$ is at most $\exp\bigl(-i(\phi_{FD}(x)-200000z_B^2)\bigr)$; and $\ontooneet\in\mathcal
M_{\epsilon_i}$.
\formal{GlobalTheory.RunJoin.GenOK}
\formal{GlobalTheory.RunJoin.gcR}
\formal{GlobalTheory.RunJoin.nbR}
\formal{GlobalTheory.RunJoin.LowW}
\formal{GlobalTheory.RunJoin.p14R}
\formal{Formalism.L2.Pj.NoCommonZero}
\formal{GlobalTheory.RunJoin.EtaOK}
\end{definition}

\begin{lemma}[Enlarged points and second norms]\label{pf:onto1:transport}
Let $X$ be a valued form algebra, $E$ an enlargement and $n$ a second norm. Then the index bound $\ontooneIB(\beta)$
holds for $X\oplus E$ if and only if it holds for $X$, and it holds for $X^n$ if and only if it holds for $X$; likewise
$(\mathrm H)$ holds for $X\oplus E$ if and only if it holds for $X$.
\formallabel{Formalism.L2.Pj.concl_xplus}
\formallabel{Formalism.L2.Pj.hypC_xplus}
\formallabel{GlobalTheory.FinalJoin.indexBoundConclusion_secondNorm_iff}
\end{lemma}

\begin{proof}
$\ontooneIB(\beta)$ involves only $\mathcal A$, $\iota$, the lattices $\mathcal M_\kappa$ and the theta lattice
(\cref{pf:onto1:rem-concl}), which are the same for $X$, $X\oplus E$ and $X^n$. $(\mathrm H)$ involves the forms, their
lattices and sup norms, the scalars and $\nu$, but neither the points nor the valuations, so it is the same for $X$ and
$X\oplus E$.
\end{proof}

\begin{lemma}[Depth bound at enlarged points]\label{pf:onto1:combBound_xplus}
Let $X$ be a valued form algebra, $E$ an enlargement, $G$ a global family datum, $S$ a scalar quadruple, $c\ge0$ and
$\nu$ a function on the primes. Assume that the $u_k$ span $W_{s_W}(w_L)$ over $\OO_K$, that $(\mathrm{DB})$ holds for
$X$, and that $E$ is \emph{regular} relative to the family $u$:
\begin{enumerate}[label=(\roman*)]
\item for every weight $w_E$, every $\eta'\in\mathcal F_{w_E}$ and every $c'\ge0$: if for all $\mathfrak l$ and
$y\in\ontoonePt(\mathfrak l)$ one has $v_{\mathfrak l,y}(\iota\eta')\ge0$ and, when this value is finite,
$\min_kv_{\mathfrak l,y}(\iota u_k)\le\nu(\mathfrak l)+c'v_{\mathfrak l,y}(\iota\eta')$, then the same holds at all points of $E$;
\item if $\min_kv_{\mathfrak l,y}(\iota u_k)\le\nu(\mathfrak l)$ at all points $y$ of $X$, then the same holds at all points of $E$.
\end{enumerate}
Then $(\mathrm{DB})$ holds for $X\oplus E$.
\formallabel{Formalism.L2.Pj.combBound_xplus}
\end{lemma}

\begin{proof}
At every point $y$ (of $X$ or of $E$) the depth equals $\min_kv_{\mathfrak l,y}(\iota u_k)$: an element of $W_{s_W}(w_L)$ is
a combination $\sum a_ku_k$ with $a_k\in\OO_K$, and $v_{\mathfrak l,y}(\iota\sum a_ku_k)\ge\min_k(\ontooneord_{\mathfrak
l}(a_k)+v_{\mathfrak l,y}(\iota u_k))\ge\min_kv_{\mathfrak l,y}(\iota u_k)$ since $\ontooneord_{\mathfrak l}(a_k)\ge0$, while the
$u_k$ themselves lie in $W_{s_W}(w_L)$. At points of $X$ the assertion is the hypothesis. At points of $E$ apply (i) with
$\eta'=\ontooneet$ and $c'=c/(2i_\eta)$ if $j<m$, and (ii) if $j\ge m$.
\end{proof}

\begin{lemma}[Lower bound for integral forms at the record]\label{pf:onto1:lamw}
There is $C_L>0$ such that for every datum $D$, top structure $T$ with laws $h$, size datum $L$ and $C_A\in\R$ satisfying
the parameter bounds (H1) with $C_A$, every $j$ and every weight $w$, the lower bound $(\mathrm A)$ holds for $B_T$ at $w$
with the constants $(C_L,A_L)=(C_L,2)$.
\formallabel{GlobalTheory.T.integralForms_logSupNorm_lower_bound_absolute}
\formallabel{GlobalTheory.RunJoin.lamwAt_absolute}
\end{lemma}

\begin{proof}
By \lref{GlobalTheory.T.meanLogSupNorm_lower_bound_of_coreLaws} there is $C_0>0$ such that for every core with
laws satisfying the congruence property described in Step 1 of the proof of \cref{pf:onto1:genok}, every size datum with
$\log g_i\le9\mathrm L+1$ for $i=0,1$, all $j,w$ and all $C\ge C_0$, $A\ge2$, the lower bound $(\mathrm A)$ holds for the form
provider of the core at $w$ with the constants $(C,A)$. The congruence property holds by the laws of $T$
(\lref{GlobalTheory.RunIn.gamAt_of_lawsCore}, which states exactly this), and $\log g_i\le9\mathrm L+1$ is part of
(H1). Take $C_L=C_0$ and $(C,A)=(C_0,2)$.
\end{proof}

\begin{theorem}[Net inputs at the record]\label{pf:onto1:inputsNet}
There are $C_L\in\R$ and functions $C_q,Z_0\colon\R^2\to\R$ with the following property. Let $D$ be a datum, $T$ a top
structure with laws $h$, $\mathrm{EX}$ an error datum, $p$ a prime, $j$ in range, $L$ a size datum, $\mathbf H$ a record of
constants, $C_2,A_2,C_T\in\R$, $\kappa,w_L,w_F$ weights, $T_n,n_P,n_W\in\N$ and $h_0\in\R$. Assume
$\mathcal C_{\mathrm{par}}(T,\mathrm{EX},p,j)$, $\mathcal C^+_{\mathrm{page}}$ for $X_{T,j}$, and $Z_0(C_A,C_K)\le z_B$. Then
$\Pi_{\mathrm{net}}$ holds for $X^+_{T,j}$ with $\mathbf K=\mathbf K_R(C_L)$, the constants
$(C_h,A_h,C_q,A_q,C_2,A_2,C_R,C_T)=(1,4,C_q(C_A,C_K),4,C_2,A_2,7,C_T)$, the scalars $S_{\mathrm{rec}}$, the integer $T_n$
and $\nu=\nu_T$; and $(\mathrm{PS})$ holds for $(B_T,S_{\mathrm{rec}},w_L)$.
\formallabel{GlobalTheory.InputsAtRecord.inputsNet_recPlus}
\formallabel{GlobalTheory.T.exists_indexBoundHypothesesWithNet_allValuedPoints}
\end{theorem}

\begin{proof}
We use the following results; all except \ref{pf:onto1:R10} and \ref{pf:onto1:R11} are proved in this section.
\begin{enumerate}[label=(R\arabic*)]
\item\label{pf:onto1:R1} \Cref{pf:onto1:lamw}: there is $C_L>0$ such that $(\mathrm A)$ holds for $B_T$ at every weight with the
constants $(C_L,2)$, whenever (H1) holds.
\item\label{pf:onto1:R2} \Cref{pf:onto1:qAt}: there is $C_q$ such that, under $\mathcal C^+_{\mathrm{page}}$, $j<m$ and $(\mathrm L)$,
the quotient bound $(\mathrm Q)$ holds with the constants $(7,C_K,1,4,C_q(C_A,C_K),4)$, the exponent $T_n$ and $c$.
\item\label{pf:onto1:R3} \Cref{pf:onto1:polyDom}: given $C_L>0$ and $C_q$ there is $Z_0$ such that $\mathcal C^+_{\mathrm{page}}$ and
$Z_0(C_A,C_K)\le z_B$ imply $G^{\max}\le\ontoonePP(1,4)$ and $(\mathrm P)$ with $\mathbf K_R(C_L)$, $(1,4,C_q(C_A,C_K),4)$, $C_R=7$ and the
auxiliary constants $(G_c,P_H,P_{rw},P_{et},P_d)=(G^{\max},\mathrm{ph}(C_L,w_L),38\mathrm L+\mathrm{PO}+1,\ontoonePP(200000,2),
\ontoonePP(1041,3)(1+\log\ontoonesize\kappa))$.
\item\label{pf:onto1:R4} \Cref{pf:onto1:runHyp}: $(\mathrm H)$ for $(X_{T,j},S_{\mathrm{rec}},\nu_T)$, $(\mathrm L)$, $W_{s(w)}(w)\subseteq\mathcal
M_w$, $(\mathrm{PS})$, and a spanning set $\mathrm{Sh}$ with $(\mathrm G2)$ for $G^{\max}$.
\item\label{pf:onto1:R5} \Cref{pf:onto1:exists_gen_eta}: an admissible generating family $G_1$ for $(S_{\mathrm{rec}},c)$ which, if
$j<m$, satisfies the auxiliary form condition.
\item\label{pf:onto1:R6} \Cref{pf:onto1:shares}: forms $\theta'_0,\dots,\theta'_{\Delta-1}\in\mathrm{Sh}$ such that the datum obtained from
$G_1$ by inserting them satisfies $(\mathrm B_{\mathfrak l})$, $(\mathrm M'_{\mathfrak l})$, $(\mathrm{Sh}_\Delta)$, $(\mathrm{Sh}_R)$ and
$(\mathrm{Sh}_q)$ for $\rho_P(7,T_n,c)$, whatever the remaining components.
\item\label{pf:onto1:R7} \Cref{pf:onto1:combBound}, applied with the Hodge--Tate inclusion for formal groups with jet charts at
every prime (\lref{GlobalTheory.ChF.hodgeTate_inclusion_chart}): $c>0$ and $(\mathrm{DB})$ for every global family datum
satisfying the auxiliary form condition if $j<m$.
\item\label{pf:onto1:R8} \Cref{pf:onto1:netrun}: the net condition $(\mathrm N)$ with $(C_G,A_G)=(1041,3)$ and
$P_d=\ontoonePP(1041,3)(1+\log\ontoonesize\kappa)$, $(\mathrm D2)$, and $\log\Delta\le1041z_B^3$ under $(\mathrm L)$ and $(\mathrm{Sh}_\Delta)$.
\item\label{pf:onto1:R9} \Cref{pf:onto1:primeData}: a family of local point data over $X^+_{T,j}$ with $(\mathrm C_{\mathfrak l})$,
$(\mathrm T_{\mathfrak l})$, $(\mathrm E'_{\mathfrak l})$ and $(\mathrm{Ht})$ with the constants of \ref{pf:onto1:R3}.
\item\label{pf:onto1:R10} \lref{Formalism.L2.Pj.regularity_R2_extended_points}: $E_{\mathrm{all}}$ is regular, in the
sense of \cref{pf:onto1:combBound_xplus}, relative to every family of forms without common zero.
\item\label{pf:onto1:R11} \lref{GlobalTheory.RunJoin.genOK_refill}: replacing the components other than $N$, the
$u_k$ and $m_x$ of an admissible generating family gives an admissible generating family.
\end{enumerate}
Choose $C_L$ by \ref{pf:onto1:R1}, $C_q$ by \ref{pf:onto1:R2} and $Z_0$ by \ref{pf:onto1:R3}, and let the data satisfy the
hypotheses. By \ref{pf:onto1:R4} we have $(\mathrm H)$, $(\mathrm L)$, $(\mathrm{PS})$ (the last assertion of the theorem)
and a set $\mathrm{Sh}$. By \ref{pf:onto1:R5} we get $G_1$, and by \ref{pf:onto1:R6} forms $\theta'_i\in\mathrm{Sh}$; let $G$ be
$G_1$ with these forms inserted. By \ref{pf:onto1:R11} $G$ is admissible; it satisfies the auxiliary form condition if
$j<m$ (same auxiliary form), and $(\mathrm G2)$ with $G^{\max}$ by \ref{pf:onto1:R4}. By \ref{pf:onto1:R7}, $(\mathrm{DB})$ holds
for $G$ and $c>0$; by \ref{pf:onto1:R8} we get the net, $(\mathrm D2)$ and $\log\Delta\le1041z_B^3$; by \ref{pf:onto1:R9} a
family $Q$ of local point data; and by \ref{pf:onto1:R3} the bound $G^{\max}\le\ontoonePP(1,4)$ and $(\mathrm P)$.

Now $\Pi_{\mathrm{base}}$ holds for $X_{T,j}$: $(\mathrm H)$ and $(\mathrm L)$ by \ref{pf:onto1:R4}, $(\mathrm A)$ at $\kappa$ by
\ref{pf:onto1:R1} (with (H1), a part of $\mathcal C^+_{\mathrm{page}}$), $(\mathrm Q)$ by \ref{pf:onto1:R2} and $(\mathrm D2)$
by \ref{pf:onto1:R8}. The family package $\Pi_{\mathrm{fam}}(G;G^{\max},P_H,P_{rw},P_{et},P_d)$ with the constants of
\ref{pf:onto1:R3} holds for $X_{T,j}$: $(\mathrm G1)$ with $(G^{\max},38\mathrm L+\mathrm{PO})$ is part of admissibility and
implies $(\mathrm G1)$ with $(G^{\max},38\mathrm L+\mathrm{PO}+1)$, since only the bound on $\log N$ changes; $(\mathrm G5)$ with
$(C_f,A_f)=(200000,2)$ is part of the auxiliary form condition; $(\mathrm D4)$ for $P_d$ is an equality for
$(C_G,A_G)=(1041,3)$; the growth bound holds because $G^{\max}$ is the maximum in its definition; and the remaining
conditions were obtained above.

Finally we pass to $X^+_{T,j}$. The packages $\Pi_{\mathrm{base}}$ and $\Pi_{\mathrm{fam}}$ involve the points only through
the integrality condition, which holds for $X^+_{T,j}$ with the same lattices, so they transfer. The net condition is
\ref{pf:onto1:R8}. $(\mathrm{DB})$ for $X^+_{T,j}$ follows from \cref{pf:onto1:combBound_xplus}: the $u_k$ span $W$ by
$(\mathrm G1)$, $c>0$ by \ref{pf:onto1:R7}, and $E_{\mathrm{all}}$ is regular relative to $u$ by \ref{pf:onto1:R10}, because the
$u_k$ have no common zero by admissibility. Together with the family $Q$ of \ref{pf:onto1:R9} this is $\Pi_{\mathrm{pt}}$,
and $\Pi_{\mathrm{net}}$ follows.
\end{proof}

\begin{lemma}[The numerical threshold]\label{pf:onto1:threshold}
For every function $Z_0\colon\R^2\to\R$ there is $x_0\colon\R^2\to\R$ such that for all $D$, $B$, all valued form algebras $X$
over $B$ at $j$ in range and all data as above, $\mathcal C_{\mathrm{page}}$ for $X$ and $x_0(C_A,C_K)\le x$ imply $\mathcal
C^+_{\mathrm{page}}$ for $X$ and $Z_0(C_A,C_K)\le z_B$.
\formallabel{GlobalTheory.T.runNumericalConditions_of_indexBound_beyond_threshold}
\end{lemma}

\begin{proof}
This is \lref{GlobalTheory.FinalJoin.pageOK_of_pageRun_beyond}, which states exactly this assertion, with
$x_0(a,b)=\max(Z_0(a,b),4b)$.
\end{proof}

\begin{theorem}[Index bound along runs]\label{pf:onto1:runholds}
There is a function $x_0\colon\R^2\to\R$ with the following property. Let $D$ be a datum, $T$ a top structure with laws $h$
such that the map $\OO\to\mathrm{LL}$ of the datum is injective, $\mathrm{EX}$ an error datum and $p\in\N$, and assume $\mathcal
C_{\mathrm{par}}(T,\mathrm{EX},p,j)$ for all $j$ in range. Let $j$ be in range, $L$ a size datum, $\mathbf H$ a record of
constants, $C_2,A_2,C_T\in\R$, $\kappa,w_L,w_F$ weights, $T_n,n_P,n_W\in\N$ and $h_0\in\R$, and assume $\mathcal
C_{\mathrm{page}}$ for $X_{T,j}$ and $x_0(C_A,C_K)\le x$. Put $\beta=\ontoonerk_\kappa\,f\,m\,b\,\delta_2$. Then:
\begin{enumerate}[label=(\alph*)]
\item there are a second norm $n$ of $B_T$, a record $\mathbf K$ and $C_q\in\R$ such that $\Pi_{\mathrm{full}}$ holds for
$X^{+,n}_{T,j}$ with $\mathbf K$, $\mathbf H$, $(C_h,A_h,C_q,A_q,C_2,A_2,C_R,C_T)=(1,4,C_q,4,C_2,A_2,7,C_T)$, the scalars
$\mathrm{scal}(B_T^n,\mathrm{EX},p,\kappa,w_L,w_F,T_n)$, $T_n$ and $\nu_T$, and $\ontooneIB(\beta)$ holds for $X^{+,n}_{T,j}$ with
these data;
\item $(\mathrm H)$ holds for $X_{T,j}$ with $\mathbf H$, $C_2,A_2$, $C_R=7$, $C_T$, $S_{\mathrm{rec}}$ and $\nu_T$;
\item the index bound $\ontooneIB(\beta)$ holds for $(X_{T,j},S_{\mathrm{rec}},\kappa,w_L,w_F,T_n,\nu_T)$.
\end{enumerate}
\formal{GlobalTheory.FinalJoin.RunHolds}
\formallabel{GlobalTheory.FinalJoin.runHolds_of_parOK}
\end{theorem}

\begin{proof}
Let $C_L,C_q,Z_0$ be given by \cref{pf:onto1:inputsNet} and $x_0$ by \cref{pf:onto1:threshold} applied to $Z_0$. Let the
data be as in the statement. By \cref{pf:onto1:threshold}, $\mathcal C^+_{\mathrm{page}}$ holds for $X_{T,j}$ and
$Z_0(C_A,C_K)\le z_B$. By \cref{pf:onto1:inputsNet}, applied with $\mathcal C_{\mathrm{par}}$ at $j$, $\Pi_{\mathrm{net}}$ holds for
$X^+_{T,j}$ with $\mathbf K_R(C_L)$ and $C_q(C_A,C_K)$, and $(\mathrm{PS})$ holds. The algebras $X^+_{T,j'}$, $1\le j'\le m$,
form a family of valued form algebras, and $\mathcal C_{\mathrm{page}}$ provides $C_1=C_c=C_2$, $A_1=A_c=A_2$, $C_{Rc}=7$,
$12\le A_2$ and $(\mathrm{PA})$. Hence \cref{pf:onto1:inputs4_of_net} gives a second norm $n$ such that $\Pi_{\mathrm{full}}$
holds for $X^{+,n}_{T,j}$ with $\mathbf K=\mathbf K_4(C_L)$, $C_q(C_A,C_K)$ and the scalars $(S_{\mathrm{rec}})^n=\mathrm{scal}
(B_T^n,\dots)$, and such that $\ontooneIB(\beta)$ holds for $X^{+,n}_{T,j}$; here $\ontoonerk_\kappa$ is the same for $B_T$ and
$B_T^n$. This is (a). The size hypotheses in $\Pi_{\mathrm{net}}$ concern $X^+_{T,j}$, and \cref{pf:onto1:transport} transfers
them to $X_{T,j}$; this is (b). By \cref{pf:onto1:transport}, $\ontooneIB(\beta)$ for $X^{+,n}_{T,j}$ is the same statement as
$\ontooneIB(\beta)$ for $X^+_{T,j}$ and for $X_{T,j}$; this is (c).
\end{proof}

\begin{remark}\label{pf:onto1:rem-runholds}
The conjunction of (a), (b), (c), for all the data quantified in \cref{pf:onto1:runholds}, is the predicate of the
development whose name is tagged first at the end of the statement. The parameter conditions
are assumed for all $j$ in range but used only at the given $j$. The injectivity of $\OO\to\mathrm{LL}$ enters only
through the scalars, and the prime $p$ only through the parameter conditions and the scalars. The existential
quantifier over $\mathbf K$ and $C_q$ in (a) is witnessed by $\mathbf K_4(C_L)$ and $C_q(C_A,C_K)$, where $C_L$ and the
function $C_q$ do not depend on the data.
\end{remark}

%----------------------------------------------------------------------------------------------------
\subsection{Admissible generating families of scaled theta lattices}\label{sec:pf:onto1:gen}

\begin{lemma}[Scaled families]\label{pf:onto1:scaledFamily}
Let $X$ be a valued form algebra over $B$, $S$ a scalar quadruple, and $w'_0,\dots,w'_{N-1}\in\mathcal F_{w_L}$ forms
spanning $\mathcal T(w_L)$ over $\OO_K$. Let $a,b\ge0$ and $G_c,P_{rw}\in\R$ with $\log|\sigma(s_W)|\le a$ and
$\log\|w'_k\|_\sigma\le b$ for all $\sigma,k$, $a+b\le G_c\ontoonesize(w_L)$, $W_{s_W}(w_L)\subseteq\mathcal M_{w_L}$ and
$\log N\le P_{rw}(\ontoonesize(w_L)-1)$. Let $G$ be the global family datum with $u_k=s_Ww'_k$, $m_x\equiv a+b$, and
all other components empty or zero. Then $G$ satisfies $(\mathrm G1)$ with $(G_c,P_{rw})$.
\formallabel{GlobalTheory.RunJoin.tiedW_scaledFamily}
\end{lemma}

\begin{proof}
The $\OO_K$-span of the $s_Ww'_k$ is $s_W\mathcal T(w_L)=W_{s_W}(w_L)$, which is contained in $\mathcal M_{w_L}$. The bound on
$\log N$ is assumed. If $s_Ww'_k\ne0$ then $\log\|s_Ww'_k\|_\sigma=\log|\sigma(s_W)|+\log\|w'_k\|_\sigma\le a+b$; if it is $0$
the left side is $\log0=0\le a+b$. Finally $[K:\Q]^{-1}\sum_\sigma(a+b)=a+b\le G_c\ontoonesize(w_L)$.
\end{proof}

\begin{theorem}[Admissible generating families]\label{pf:onto1:genok}
Let $D$ be a datum, $T$ a top structure with laws $h$ such that $\OO\to\mathrm{LL}$ is injective, $\mathrm{EX}$ an error
datum, $p\in\N$, $j$ in range, $L$ a size datum, $\mathbf H$ a record of constants, $C_2,A_2,C_T\in\R$,
$\kappa,w_L,w_F$ weights, $T_n,n_P,n_W\in\N$ and $h_0\in\R$. Assume that $T$ is oriented for $\mathrm{EX}$ and that $\mathcal
C^+_{\mathrm{page}}$ holds for $X_{T,j}$. Then there is an admissible generating family $G$ for $(X_{T,j},\kappa,w_L,w_F)$
with respect to $(S_{\mathrm{rec}},c)$, where $c=R(7)\ontoonePP(C_2,A_2)\ontooneDF(w_L)^2/\min(g_0,g_1)$.
\formallabel{GlobalTheory.RunJoin.exists_genOK}
\formallabel{GlobalTheory.T.exists_admissible_thetaLatticeGenerators}
\end{theorem}

\begin{proof}
\emph{Step 1: preliminary facts.} We use the following results proved elsewhere in the development.
\begin{itemize}
\item \lref{GlobalTheory.T.harmonic_germ_value_chart_independent}: values of harmonic polynomials on germs of
sections do not depend on the chosen chart (the statement denoted \emph{chart independence} below).
\item \lref{GlobalTheory.RunIn.gamAt_of_lawsCore}: by the laws of $T$, at every good embedding and every
component, every unitary element congruent to the identity modulo the level on the lattice of the component lies in
the arithmetic group of that component.
\item \lref{GlobalTheory.pageRows_of_pageAt}: under $\mathcal C^+_{\mathrm{page}}$ the \emph{page-row bounds} hold at
$\kappa$ ($\log g_i\le9\mathrm L+1$, $\log f\le2\mathrm L$, $\log p_{\mathrm{de}},\log p'_{\mathrm{de}}\le\mathrm L$, $\log(1+\ontooneDF
\kappa)\le2\mathrm L$, $\log n_\kappa\le\mathrm L$ and three similar bounds for invariants of the host), and $16\le f$.
\item \lref{GlobalTheory.page_two_le_L_and_signs}: under the page hypotheses and the page-row bounds,
$2\le\mathrm L$, $6\le r$ and $\lambda^*\ge0$.
\item \lref{GlobalTheory.RunWeights.weights_below}: under the page hypotheses, $\ontooneDF(w_L)\le\ontooneDF(\kappa)$,
$n_{w_L}\le n_\kappa$, $r\le n_{w_L}$ and $a^+_\tau(p_j-1)\ge1$ for $w_L$; and \lref{GlobalTheory.pageRows_of_le}: the
page-row bounds at $\kappa$ imply those at every weight $w$ with $\ontooneDF(w)\le\ontooneDF(\kappa)$ and $n_w\le n_\kappa$, in
particular at $w_L$.
\item \lref{GlobalTheory.RunIn.sW_eq_sOf} and \lref{GlobalTheory.RunIn.sOf_ne_zero_and_log_bounds}: the first
entry of $S_{\mathrm{rec}}$ is $s_W=s(w_L)$, it is non-zero, and for every $\sigma$:
$0\le\log|\sigma(s_W)|$, $\frac12\ontooneDF(w_L)\log\min(g_0,g_1)\le\log|\sigma(s_W)|$,
$\frac12\ontooneDF(w_L)\log(g_0g_1)\le\log|\sigma(s_W)|$ and $\log|\sigma(s_W)|\le\lambda^*\ontoonesize(w_L)$.
\item \lref{GlobalTheory.T.scaledThetaLattice_le_supLattice_every_weight}: if $T$ is oriented for $\mathrm{EX}$, then
$W_{s(w)}(w)\subseteq\mathcal M_w$ for every weight $w$.
\item \lref{GlobalTheory.RunIn.sup_bounds_le_Gc}: an explicit inequality which gives, with
$b_0=400z_B\ontoonesize(w_L)+3frn_{w_L}\mathrm L+5\log[K:\Q]+\log|d_K|/[K:\Q]+2$,
$(\lambda^*+2\log p_{\mathrm{de}})\ontoonesize(w_L)+b_0\le\mathrm{Gc}(p_{\mathrm{de}})\ontoonesize(w_L)$, provided $2\le\mathrm L$ and
$n_{w_L}>0$.
\end{itemize}
From these, $(\mathrm{PS})$ holds for $(B_T,S_{\mathrm{rec}},w_L)$ (its first, third, fourth and fifth items are part of
$\mathcal C^+_{\mathrm{page}}$, and the second is the third bound on $\log|\sigma(s_W)|$), and $(\mathrm K)$ holds for
$\mathbf K=(0,2,0,2,0,2,0,4,7)$, $(C_h,A_h,C_q,A_q)=(1,4,0,11)$ and $C_R=7$ (by the equalities of constants in
$\mathcal C_{\mathrm{page}}$). By \lref{GlobalTheory.RunParametersAtRecord.run_at_record}, which gives under the page
hypotheses, $(\mathrm K)$, $(\mathrm{PA})$ and $(\mathrm{PS})$ explicit bounds on the parameters of the theta runs, the
inequality $fr(r+1)\le2n_{w_L}$ holds.

\emph{Step 2: a short theta family.} By \lref{GlobalTheory.T.exists_short_thetaFamily_with_block_polynomial_rules},
which applies under chart independence, the congruence property of Step 1, the page-row bounds at $w_L$, $2\le\mathrm L$,
$r\le n_{w_L}$, $a^+_\tau(p_j-1)\ge1$ and $fr(r+1)\le2n_{w_L}$, there are $N\in\N$, forms
$w'_0,\dots,w'_{N-1}\in\mathcal F_{w_L}$ and a harmonic polynomial $P_1$ which is a product of blocks with disjoint columns,
such that: the family $(w'_k)$ contains the theta rules of $P_1$ (each theta rule attached to $P_1$ and to a tuple of level
representatives occurs as some $w'_k$, injectively in the tuple); the $w'_k$ span $\mathcal T(w_L)$ over $\OO_K$;
\[
N\le2(g_0g_1)^{frn_{w_L}}[K:\Q](1+\ontooneDF(w_L))^{frn_{w_L}};
\]
and $\log\|w'_k\|_\sigma\le b_0$ for all $k,\sigma$. In particular $N>0$, since the injective assignment of theta rules
has a non-empty source.

\emph{Step 3: the scaled family.} Put $a=\lambda^*\ontoonesize(w_L)\ge0$ and $b=b_0\ge0$. By Step 1,
$\log|\sigma(s_W)|\le a$ and $a+b\le(\lambda^*+2\log p_{\mathrm{de}})\ontoonesize(w_L)+b_0\le\mathrm{Gc}(p_{\mathrm{de}})
\ontoonesize(w_L)\le G^{\max}\ontoonesize(w_L)$, using $\log p_{\mathrm{de}}\ge0$. By
\lref{GlobalTheory.RunIn.log_two_mul_count_le_mul_size}, which derives $\log N\le P_{rw}(\ontoonesize(w)-1)$ from the
bound on $N$ of Step 2, the page-row bounds at $w$, $2\le\mathrm L$, $n_w>0$, the bounds $\log[K:\Q]\le\mathrm{PO}$ and
$\log|d_K|/[K:\Q]\le\mathrm{PO}$ contained in (H2), and $P_{rw}\ge38\mathrm L+\mathrm{PO}$, we have $\log N\le(38\mathrm
L+\mathrm{PO})(\ontoonesize(w_L)-1)$. Since $W_{s_W}(w_L)\subseteq\mathcal M_{w_L}$ by Step 1, \cref{pf:onto1:scaledFamily}
shows that the scaled family $G$ ($u_k=s_Ww'_k$, $m_x\equiv a+b$) satisfies $(\mathrm G1)$ with $(G^{\max},38\mathrm
L+\mathrm{PO})$. This is condition (i) of admissibility.

\emph{Step 4: the count.} $N>0$ by Step 2. Taking logarithms in the bound for $N$,
$\log N\le\log2+frn_{w_L}(\log(g_0g_1)+\log(1+\ontooneDF(w_L)))+\log[K:\Q]\le\mathrm{nb}(w_L)$, since $\log[K:\Q]\le\mathrm{PO}$.
This is (ii).

\emph{Step 5: no common zero.} By \lref{GlobalTheory.T.scaledThetaLattice_generators_noCommonZero}, which states
that under chart independence, the page hypotheses and $s_W\ne0$ the forms $u_k$ of a global family datum satisfying
$(\mathrm G1)$ have no common zero, condition (iii) holds.

\emph{Step 6: the lower bound at frames.} The family $u_k=s_Ww'_k$ contains the theta rules of $P_1$ scaled by $s_W$:
for each tuple $q$ of level representatives, $u_{k(q)}$ is $s_W$ times the theta rule attached to $P_1$ and $q$. By
\cref{pf:onto1:frames} (proved in \cref{sec:pf:onto1:frames}), which applies under chart independence, the page hypotheses,
$(\mathrm K)$, $(\mathrm{PA})$ and $(\mathrm{PS})$ and gives the lower bound at frames of \cref{pf:onto1:genok-def}(iv) with any
$t\ge p_{14}^\circ(w_L)$ in place of $p_{14}(w_L)$, where $p_{14}(w_L)=\max(1,p^\circ_{14}(w_L))$, we obtain (iv) with $t=p_{14}(w_L)$.
\end{proof}

\begin{lemma}[Generating family with auxiliary form]\label{pf:onto1:exists_gen_eta}
In the situation of \cref{pf:onto1:genok}, there is an admissible generating family $G$ for $(S_{\mathrm{rec}},c)$ which, if
$j<m$, satisfies the auxiliary form condition for $(i_\eta,\ontooneet)$.
\formallabel{GlobalTheory.InputsAtRecord.exists_gen_eta}
\end{lemma}

\begin{proof}
Take $G$ from \cref{pf:onto1:genok}. If $j\ge m$ there is nothing more to show. If $j<m$, by
\lref{GlobalTheory.EtaBoundForm.exists_etaOK}, which states that under $\mathcal C^+_{\mathrm{page}}$ and $j<m$ there are
$i\in\N$ and $\ontooneet\in\mathcal F_{\epsilon_i}$ satisfying the auxiliary form condition, we replace the auxiliary form of
$G$ by $(i,\ontooneet)$. Admissibility only involves $N$, the $u_k$ and $m_x$, so the new datum is still admissible
(\lref{GlobalTheory.RunJoin.genOK_refill}, which states exactly this).
\end{proof}

\begin{remark}\label{pf:onto1:rem-genok}
In \cref{pf:onto1:genok} the prime $p$ and the weight $w_F$ enter only through the scalar quadruple $S_{\mathrm{rec}}$, of
which only the first entry $s_W$ is used, and $\kappa$, $T_n$, $n_P$, $n_W$, $h_0$ enter only through the page conditions.
The substance of the theorem lies in the short theta family of Step 2 (cited) and in the lower bound at frames of Step 6,
which is proved in \cref{sec:pf:onto1:frames} from a lower bound for theta rules at orthogonal frames
(\cref{pf:onto1:thetaLower}) and an inequality for derivatives of translates of theta functions on principal tori
(\cref{pf:onto1:Iholds}).
\end{remark}

%----------------------------------------------------------------------------------------------------
\subsection{Proofs of the inputs at the record}\label{sec:pf:onto1:inputs}

In this subsection we prove the results \ref{pf:onto1:R2}--\ref{pf:onto1:R4} and \ref{pf:onto1:R6}--\ref{pf:onto1:R9} used in the
proof of \cref{pf:onto1:inputsNet}, together with two numerical statements used in the proof of the page conditions at
the record (\lref{GlobalTheory.FinalJoin.RecordRun.pageRun_record}). We first collect the numerical facts.

\begin{definition}[Parameter choices and size conditions]\label{pf:onto1:parchoice}
The ranks of a datum satisfy $f=2n$, $2\le m$, $m+3\le k$ and $r=k+m-1$. A \emph{parameter choice} for $(D,L)$ is a
structure of the development fixing auxiliary parameters of the run (among them the integers $T_O$, $n_P$, $n_W$ and a
constant $C_{nW}$); among its properties are $|\log\mathrm{Der}|\le x^{\eps}$, $\log p_{\mathrm{de}},\log p'_{\mathrm{de}}\le\mathrm L$ and
$n_W\le C_{nW}r\bigl(f+f^2(k-1)^3\mathrm L^2/m\bigr)$, and it implies $\log\mathrm{rd}_E\le\mathrm L$
(\lref{Formalism.BasicNotions.log_rdE_le_scL}). The \emph{size conditions} on $L$
are explicit windows for the ranks and parameters in terms of $x$ and $\eps$; among them are
$x^{2+8\eps-\eps/100}\le n\le x^{2+8\eps+\eps/100}$, $k\le x^{9/2+10\eps+\eps/100}$ and $x^{-\eps}\le\delta_2$. We write
$N_{\mathrm{cond}}=Q_Sqd$ for the conductor integer of the host, and recall $\log\Omega=\gamma_{\mathrm{EM}}+\log(2\pi/\mathrm{rd}_E)-\log
\mathrm{Der}$, where $\gamma_{\mathrm{EM}}\in(\frac12,\frac23)$ is Euler's constant.
\formal{Formalism.BasicNotions.ParameterChoice'}
\formal{Formalism.BasicNotions.ParamsT}
\formal{Formalism.MainNotions.SizesAt}
\formal{Formalism.BasicNotions.SizeData₂.logOm}
\end{definition}

Consequently $6\le r\le2k$ and $m\le r$, and, as $0<\eps<\frac1{52}$, the size conditions give $x\le n\le x^3$ and $k\le x^5$.
Moreover every parameter choice satisfies $\log N_{\mathrm{cond}}\le\mathrm L/72$, $\log f\le\mathrm L/72$ and $\log r\le\mathrm L/300$
(\lref{Formalism.MainNotions.Gaps.EtaleWeightGaps.Num.Pm.log_cond_le},
\lref{Formalism.MainNotions.Gaps.EtaleWeightGaps.Num.Pm.log_f_le} and
\lref{Formalism.MainNotions.Gaps.EtaleWeightGaps.Num.Pm.log_r_le}, which state these three bounds). We recall
$\mathrm L=x^{\eps}\ge1$ and $z_B=fr\mathrm L$.

\begin{proposition}[Parameter bounds]\label{pf:onto1:parA}
Let $L$ be a size datum satisfying the size conditions, $P$ a parameter choice for $(D,L)$, and assume $\log g_i\le9\mathrm
L+1$ for $i=0,1$ and $x\ge10^{60}$. Then the parameter bounds (H1) hold with $C_A=20$
(\lref{Formalism.MainNotions.WeightAlgebra.ParA}): $x\le fr\le x^{20}$, $\log|d_E|\le20fr\mathrm L$, $x^{-\eps}\le\delta_2$,
$\log g_i\le9\mathrm L+1$, $|\log\mathrm{Der}|\le x^{\eps}$, $|\log\Omega|\le20\mathrm L$, $\log p_{\mathrm{de}},\log p'_{\mathrm{de}}\le\mathrm
L$ and $2fm\le\mathrm{rd}_E$.
\formallabel{Formalism.MainNotions.Gaps.EtaleWeightGaps.Num.Pm.parA_of}
\formallabel{Formalism.MainNotions.Gaps.EtaleWeightGaps.Num.Pm.x_le_fr}
\formallabel{Formalism.MainNotions.Gaps.EtaleWeightGaps.Num.Pm.logOm_le}
\end{proposition}

\begin{proof}
$fr=2nr\ge12n\ge x$. Next $fr=2nr\le4nk\le4x^8\le x^9\le x^{20}$ (as $x\ge4$). Since $E$ is a CM field of degree $2f$,
$\log|d_E|=2f\log\mathrm{rd}_E\le2f\mathrm L\le20fr\mathrm L$. The bound $x^{-\eps}\le\delta_2$ is a size condition, $\log g_i\le9\mathrm
L+1$ is assumed, and the bounds on $\log\mathrm{Der}$, $p_{\mathrm{de}}$, $p'_{\mathrm{de}}$ are properties of parameter choices. For
$\log\Omega$: $0\le\log\mathrm{rd}_E\le\mathrm L$, $0<\gamma_{\mathrm{EM}}<1$, $0<\log2<1$, $0\le\log\pi\le2$ and $|\log\mathrm{Der}|\le\mathrm L$
give $|\log\Omega|\le\mathrm L+\mathrm L+3\le6\mathrm L\le20\mathrm L$. Finally $2fm\le\mathrm{rd}_E$ holds for every host
(\lref{UpperTheory.Ram.two_fm_le_rdE}, which states $2fm\le\mathrm{rd}_E$ and $4tr+1\le\mathrm{rd}_E$).
\end{proof}

\begin{proposition}[Domination at the last index]\label{pf:onto1:domAt}
Let $L$ be a size datum satisfying the size conditions, $P$ a parameter choice for $(D,L)$ with parameters $T_p$ whose
constant $C_{nW}$ equals $10$, and assume $\log g_i\le9\mathrm L+1$ and $x\ge10^{60}$. Then the domination conditions
(\lref{Formalism.MainNotions.WeightAlgebra.DomAt}) hold at the index $m$, for the integer $n_W$ of $T_p$ and $C_P=1000$:
$1000^{20}\le z_B$; $\lambda^*,\mathrm{PN},\log(N_{\mathrm{cond}}g_0g_1),\mathrm{PU}\le1000z_B$; $\mathrm{PO}\le1000z_B^2$;
$n_0(m)\le1000z_B^3$; $\mathrm{sc}(n_W)\le1000z_B^4$; and $D_W(m)\le1000z_B^{10}$. Here $\mathrm{PU}=fr(2\log\mathrm{rd}_E+2\log(fr)+3)$,
$n_0(m)=r+2g_0$ with $g_0=1250r$,
\begin{gather*}
\mathrm{sc}(n_W)=\max\bigl(90(fr+1),\,360fm^3,\,2(3r+3+r\log r),\,n_W,\,g_0,\,fr^2+6,\,11\bigr),\\
D_W(m)=72\cdot5^4\,\delta_m^2f^2r^2m^4 .
\end{gather*}
\formallabel{Formalism.MainNotions.Gaps.EtaleWeightGaps.Num.Pm.domAt_of}
\formallabel{Formalism.MainNotions.Gaps.EtaleWeightGaps.Num.Pm.polyConstant_le_zBase_pow_four}
\formallabel{Formalism.MainNotions.Gaps.EtaleWeightGaps.Num.Pm.liftConstant_le_zBase_pow_ten}
\formallabel{Formalism.MainNotions.Gaps.EtaleWeightGaps.Num.Pm.log_qdg_le}
\formallabel{Formalism.MainNotions.Gaps.EtaleWeightGaps.Num.Pm.PU_le}
\formallabel{Formalism.MainNotions.Gaps.EtaleWeightGaps.Num.Pm.PO_le}
\formallabel{Formalism.MainNotions.Gaps.EtaleWeightGaps.Num.Pm.lamStar_le}
\end{proposition}

\begin{proof}
We use $\mathrm L\ge1$, $f\ge2$, $r\ge6$, so $\mathrm L\le z_B$ and $f^ar^b\mathrm L^c\le z_B^{\,e}$ whenever $a,b,c\le e$. By
\cref{pf:onto1:parA}, $z_B\ge fr\ge x\ge10^{60}=1000^{20}$. From $\log g_i\le9\mathrm L+1$ we get $\log(g_0g_1)\le18\mathrm L+2$.

$\lambda^*$: one has $4(f/2)^2qd=4n^2qd\le N_{\mathrm{cond}}^2$ (\lref{UpperTheory.Ram.cond_le_four_n_sq}, which states this
inequality for every host), so $\lambda^*\le18\mathrm L+2+2\log N_{\mathrm{cond}}\le21\mathrm L$. Then $\mathrm{PN}=2\mathrm L+\lambda^*+4\le
27\mathrm L$, and $\log(N_{\mathrm{cond}}g_0g_1)\le\mathrm L/72+18\mathrm L+2\le21\mathrm L$. Next $2\log\mathrm{rd}_E+2\log(fr)+3\le2\mathrm
L+\mathrm L/36+\mathrm L/150+3\le7\mathrm L$, so $\mathrm{PU}\le7z_B$. For $\mathrm{PO}=19fr^2\mathrm L+198f^2\mathrm L$: $fr^2\mathrm L\le
z_B^2$ and $f^2\mathrm L\le z_B^2$, so $\mathrm{PO}\le217z_B^2$. Since $q_m=0$, $n_0(m)=r+2g_0=2501r\le2501z_B\le1000z_B^3$.

$\mathrm{sc}(n_W)$: $90(fr+1)\le180fr$, $360fm^3\le360fr^3$, $2(3r+3+r\log r)\le2(3r+3+r\mathrm L)$, $g_0=1250r$, $fr^2+6$
and $11$ are all at most $1000z_B^4$. For $n_W$, the parameter choice with $C_{nW}=10$ gives
$n_W\le10r\bigl(f+f^2(k-1)^3\mathrm L^2/m\bigr)\le10fr+10f^2r^4\mathrm L^2\le28z_B^4$, using $k-1\le r$.

$D_W(m)$: $\delta_m\le fr^2$ (\lref{Formalism.MainNotions.Gaps.EtaleWeightGaps.Num.Pm.dimY_m_le}), $m\le r$ and $f^6\ge64$ give
$D_W(m)\le45000f^4r^{10}\le1000f^{10}r^{10}\le1000z_B^{10}$.
\end{proof}

\begin{proposition}[Polynomial domination]\label{pf:onto1:polyDom}
For every $C_L>0$ and every function $C_q\colon\R^2\to\R$ there is $Z_0\colon\R^2\to\R$ such that the following holds. Let
$X$ be a valued form algebra over a form provider $B$ at $j$ in range, and let the data satisfy the strengthened page
conditions $\mathcal C^+_{\mathrm{page}}$ and $Z_0(C_A,C_K)\le z_B$. Then $G^{\max}\le\ontoonePP(1,4)$, and the domination
conditions $(\mathrm P)$ hold for $\mathbf K=\mathbf K_R(C_L)$, $(C_h,A_h,C_q,A_q)=(1,4,C_q(C_A,C_K),4)$, $C_R=7$ and the auxiliary
constants $G_c=G^{\max}$, $P_H=\mathrm{ph}(C_L,w_L)$, $P_{rw}=38\mathrm L+\mathrm{PO}+1$, $P_{et}=\ontoonePP(200000,2)$ and
$P_d=\ontoonePP(1041,3)(1+\log\ontoonesize\kappa)$. Here
\[
\mathrm{ph}(C_L,w)=2\bigl(\lambda_R(C_L,w)+\max(1,\mathrm{PN})+p_{14}(w)+200000z_B^2+\max(1,\Pi_4)+P_{rw}\bigr)+G^{\max}
\]
with explicit quantities $\lambda_R(C_L,w)$ and $\Pi_4$.
\formallabel{GlobalTheory.RunNum.polyDom_of_pageOK}
\formal{Formalism.MainNotions.Gaps.EtaleWeightGaps.PolyDom}
\formal{GlobalTheory.RunJoin.phR}
\end{proposition}

\begin{proof}
Put $Z_0(a,b)=C_L+|b|+|C_q(a,b)|$. Write $z=z_B$ and $\mathsf P_2=\ontoonePP(C_2,A_2)$. By $\mathcal C_{\mathrm{page}}$ the domination
conditions hold for $C_P=C_1=C_2$, so $1000\le C_2$, $C_2^{20}\le z$ (hence $C_2\le z$ and $z\ge1000^{20}\ge10^9$), and $A_2\ge12$;
thus $\mathsf P_2=C_2z^{A_2}\ge1000z^{12}$. By the choice of $Z_0$, $C_K,C_q(C_A,C_K),C_L\le z$. The following bounds are
consequences of $\mathcal C^+_{\mathrm{page}}$ and these facts:
\begin{itemize}
\item $n_{w_L}=n_W\le z^5$, $\log\ontoonesize\kappa\le3z$ and $\log\ontoonesize w_L\le3z$
(\lref{GlobalTheory.RunNum.log_size_le}, which derives $\log\ontoonesize w\le3z$ from $\ontoonesize w\le C_Kkfb$ and $C_K\le z$);
\item $1\le G^{\max}\le2000z^2$ (\lref{GlobalTheory.RunNum.Gc_le}, which bounds the growth constant $\mathrm{Gc}(p)$ by
$2000z^2$ when $\log p\le\mathrm L$, using the field bounds and $\log g_i\le9\mathrm L+1$), and $1\le P_{rw}\le256z^2$ by the bound
$\mathrm{PO}\le217z^2$ of the proof of \cref{pf:onto1:domAt};
\item $0\le\ontoonePPz(C_L,2)\le900z^3$ (\lref{GlobalTheory.RunNum.runBaseConstant_two_bounds}, which states this for
$0\le C_L\le z$ and every parameter choice);
\item $\lambda_R(C_L,w_L)\le2\cdot10^6z^{10}$, $p_{14}(w_L)\le35z^7$, $\max(1,\mathrm{PN})\le27z^2$ and $\Pi_4\le6000z^2$
(\lref{GlobalTheory.RunNum.lamR_le}, \lref{GlobalTheory.RunNum.thetaGeneratorArchLowerConst_le} and
\lref{GlobalTheory.RunNum.runIndexConstant_le_zBase_sq}, which state the first, second and fourth of these bounds under
the preceding ones, and \cref{pf:onto1:domAt} for $\mathrm{PN}$); hence $P_H\le5\cdot10^6z^{10}$, and $P_{rw}\le P_H$;
\item $P_d=1041z^3(1+\log\ontoonesize\kappa)\le4164z^4$, $\ontoonePPz(C_L,2)(1+\log\ontoonesize\kappa)\le3600z^4$ and
$\ontoonePP(C_q,4)\le z^5$.
\end{itemize}
The conditions $(\mathrm P)$ are then verified one by one: $C_{Rm}=7\ge2$, $C_R=7\ge0$, $G_c,P_{rw}\ge0$, $\ontoonePPz(C_L,2)\ge0$; the
field bounds for $(C_h,A_h)=(1,4)$ follow from the field bounds of $\mathcal C^+_{\mathrm{page}}$ (\lref{GlobalTheory.RunNum.obound_of_opo});
$R(C_{Rc})\ontoonePP(C_c,A_c)=R(7)\mathsf P_2$ by the ties of constants; $R(7)\mathrm{rd}_E^2=\mathrm{rd}_E^9\le\mathrm{rd}_E^{14}=R(7)^2$;
$2\le\mathsf P_2^2$; $4C_K\le\mathsf P_2^2f$ is part of $\mathcal C^+_{\mathrm{page}}$; and the remaining inequalities, namely
$25(2G_c+6P_H+P_{rw})\le\mathsf P_2$, $P_{et}\le\mathsf P_2$, $\Xi\le\mathsf P_2$, $60G_c+180P_H+44\,\Xi+44\ontoonePP(C_q,4)\le
\mathsf P_2^2\delta_j$ (with $\Xi=2G_c+P_{rw}+\ontoonePPz(C_L,2)(1+\log\ontoonesize\kappa)+\frac32(P_d+\mathrm{PO})$), $\frac12\le P_{rw}\le P_H$,
$0\le P_{et}$ and $P_{rw}\le\mathsf P_2$, follow from the bounds above, $\mathsf P_2\ge1000z^{12}$, $z\ge10^9$ and $\delta_j\ge1$. Finally
$G^{\max}\le2000z^2\le z^5=\ontoonePP(1,4)$.
\end{proof}
\begin{theorem}[Nets at the record]\label{pf:onto1:netrun}
Let $D,T,h$, $j$ in range and the data be as in \cref{pf:onto1:inputsNet}, and assume $\mathcal C^+_{\mathrm{page}}$ for $X_{T,j}$.
Then the net condition $(\mathrm N)$ holds for $(B_T,j,\kappa)$ with the constants $(C_G,A_G)=(1041,3)$ and
$P_d=\ontoonePP(1041,3)(1+\log\ontoonesize\kappa)$; $(\mathrm D2)$ holds for $X_{T,j}$; $\log\dim_K\mathcal F_\kappa\le1041z_B^3$; and for
every $\Delta\in\N$, $(\mathrm L)$ and $(\mathrm{Sh}_\Delta)$ (with $T_n$) imply $\log\Delta\le1041z_B^3$.
\formallabel{GlobalTheory.RunJoin.net_at_run}
\formallabel{GlobalTheory.RunIn.NetAtRecord.netData_exists}
\formallabel{GlobalTheory.exists_scalar_readings_on_net}
\formallabel{GlobalTheory.RunJoin.net_count_dim}
\end{theorem}

\begin{proof}
\emph{Step 1: preliminary bounds.} By $\mathcal C^+_{\mathrm{page}}$, the page-row bounds hold at $\kappa$ and $16\le f$
(\lref{GlobalTheory.pageRows_of_pageAt}), $2\le\mathrm L$ (\lref{GlobalTheory.EtaAn.page_rows}, which derives it from the
parameter bounds and the size conditions), and $8\le r$; the weight conditions give $n_\kappa=n_V=n_P+T_nn_W\ge8fr+4\ge r$; and
$z_B=fr\mathrm L\ge1$.

\emph{Step 2: evaluations at a good embedding.} Let $\sigma$ be a good embedding of the coefficient field. By
\lref{GlobalTheory.RunIn.exists_net_max_at_goodEmb}, which states that under the congruence property, the page-row
bounds, $2\le\mathrm L$, $r\le n_\kappa$ and $16\le f$ there is a non-empty finite set $\mathcal N$ of frame points $p$ of the
unfolding at $\sigma$ such that every holomorphic section $P$ of the relevant line bundle satisfies $\sup|P|\le4\,|P(p)|_p$
for some $p\in\mathcal N$ (pointwise norm at $p$), and $\log(16|\mathcal N|)\le1041z_B^3$, we obtain such a set. The sup norm
of a form $F\in\mathcal F_\kappa$ at $\sigma$ is a constant $L_\sigma>0$ times the sup norm of its realisation, and the pointwise norm
at $p$ is $a_p$ times the absolute value of the value of the realisation at $p$, for an explicit factor $a_p\ge0$. Put
$\mathrm{ev}_{\sigma,p}(F)=L_\sigma a_p\cdot(\text{value of the realisation of }F\text{ at }p)$; this is $\sigma$-semilinear, and
$\|F\|_\sigma/4\le|\mathrm{ev}_{\sigma,p}(F)|$ for some $p$, while $|\mathrm{ev}_{\sigma,p}(F)|\le\|F\|_\sigma$ for all $p$. If $F_1,\dots,F_n$ are
linearly independent over $K$, then their realisations at $\sigma$ are linearly independent over $\C$
(\lref{GlobalTheory.cusp_readings_linearIndependent}, which states that realisation commutes with base change and is
injective); if $\sum_ic_i\mathrm{ev}_{\sigma,p}(F_i)=0$ for all $p$, the section $\sum_ic_i\cdot(\text{realisation of }F_i)$ vanishes on
$\mathcal N$, hence has sup norm $0$ by the net property, so all $c_i=0$. Thus the vectors $(\mathrm{ev}_{\sigma,p}(F_i))_p$ are
linearly independent.

\emph{Step 3: all embeddings.} An embedding $\sigma$ is good if and only if its conjugate is not, and sup norms are
invariant under conjugation. Defining the evaluations at a non-good $\sigma$ as the complex conjugates of those at
$\bar\sigma$ (\lref{GlobalTheory.T.exists_netData_from_half_conjugate_embeddings}, which carries out this construction)
gives a net with the first two items of $(\mathrm N)$, $\log(16|\mathcal N_\sigma|)\le1041z_B^3$ for all $\sigma$, and the
independence at every $\sigma$.

\emph{Step 4: counting.} Since $1041\le C_G$, $3\le A_G$ and $z_B\ge1$, $1041z_B^3\le\ontoonePP(1041,3)(1+\log\ontoonesize\kappa)$; this
gives the second item of $(\mathrm N)$, and $\log|\mathcal N_{\sigma_0}|\le\log(16|\mathcal N_{\sigma_0}|)\le P_d$ for any embedding $\sigma_0$.
Applying the independence to a $K$-basis of $\mathcal F_\kappa$ gives $\dim_K\mathcal F_\kappa\le|\mathcal N_{\sigma_0}|$, hence $\log
\dim_K\mathcal F_\kappa\le1041z_B^3$. Under $(\mathrm L)$, $(\mathrm{Sh}_\Delta)$ reads $\binom{T_n-1}{\delta}\Delta\le\dim_K\mathcal F_\kappa$
with $\delta+1\le T_n$, so $\Delta\le\dim_K\mathcal F_\kappa$ and $\log\Delta\le1041z_B^3$.

\emph{Step 5.} $(\mathrm D2)$ for $X_{T,j}$ is \lref{GlobalTheory.deepSubMul_rec_all}, which states that the sup norms of the
record are submultiplicative on products of forms.
\end{proof}

\begin{theorem}[Size hypotheses at the record]\label{pf:onto1:runHyp}
Let $D,T,h$ with $\OO\to\mathrm{LL}$ injective, $\mathrm{EX}$, $p$, $j$ in range and the data be as in \cref{pf:onto1:inputsNet},
and assume $\mathcal C_{\mathrm{par}}(T,\mathrm{EX},p,j)$ and $\mathcal C^+_{\mathrm{page}}$ for $X_{T,j}$. Then $(\mathrm H)$ holds for
$(X_{T,j},S_{\mathrm{rec}},\nu_T)$ with $C_R=7$; $(\mathrm L)$ holds; $W_{s(w)}(w)\subseteq\mathcal M_w$ for every weight $w$; $(\mathrm{PS})$
holds for $(B_T,S_{\mathrm{rec}},w_L)$; and there is a finite set $\mathrm{Sh}\subseteq\mathcal F_{w_F}$ spanning $\mathcal F_{w_F}$ such that
$(\mathrm G2)$ with $G_c=G^{\max}$ holds for every global family datum whose forms $\theta'_i$ lie in $\mathrm{Sh}$.
\formallabel{GlobalTheory.T.runHypotheses_formsScalars_of_parameterConditions}
\formallabel{GlobalTheory.RunJoin.hypC_record}
\end{theorem}

\begin{proof}
\emph{Step 1: numerical consequences.} By $\mathcal C^+_{\mathrm{page}}$: the page-row bounds at $\kappa$
(\lref{GlobalTheory.pageRows_of_pageAt}); $\ontooneDF(w_L)\le\ontooneDF(\kappa)$
(\lref{GlobalTheory.RunWeights.weights_below}), hence $\log\ontooneDF(w_L)\le2\mathrm L$; and $2fr\le n_{w_L}$, since
$n_{w_L}=n_W\ge fr(r+1)$ and $r\ge8$.

\emph{Step 2: the prime sum bound (H4).} By \lref{GlobalTheory.DepthSum.depth_sum_record}, which states that for
$2fr\le n_w$ and $\log\ontooneDF(w)\le2\mathrm L$ the function $\nu_T$ attached to $w$ is non-negative, finitely supported and
satisfies $\sum_{\mathfrak l}\nu_T(\mathfrak l)\ell(\mathfrak l)\le\mathrm{PN}\,\ontoonesize(w)$ (on $X_{T,j}$ and on its enlargements),
(H4) holds for $\nu_T$.

\emph{Step 3: integrality (H5).} The parameter conditions give that $p$ is the prime of the shortening data at $j$, and that
$T$ is oriented for $\mathrm{EX}$. By \lref{GlobalTheory.T.scaledRunLattices_in_supLattices}, which states that under these
conditions the scaled theta lattices $W_{s_W}(w_L)$, $\mathrm{Th}_{s_F}(w_F)$, $\mathrm{Th}_{s_0}(w_0)$ and $\mathrm{Th}_{s_K}(\kappa)$ of the
quadruple $S_{\mathrm{rec}}$ consist of integral forms, (H5) holds, and by
\lref{GlobalTheory.T.scaledThetaLattice_le_supLattice_every_weight}, $W_{s(w)}(w)\subseteq\mathcal M_w$ for every $w$.

\emph{Step 4: the remaining size hypotheses.} By \lref{GlobalTheory.T.thetaIndexBound_hypotheses_at_instance}, which
states that under chart independence, the congruence property, the laws of the shortening data, the positivity of the
bad integer, the triviality of the levels, the sharing of group and character data, the injectivity of $\OO\to\mathrm{LL}$,
the page conditions, the integrality of Step 3 and the bounds of the growth constant, the hypotheses (H1)--(H3) and
(H6)--(H7) hold for $(X_{T,j},S_{\mathrm{rec}})$ with $C_R=7$ as soon as (H4) holds, and that there is a finite spanning set
$\mathrm{Sh}\subseteq\mathcal F_{w_F}$ as in the statement, Step 2 gives $(\mathrm H)$ and the set $\mathrm{Sh}$.

\emph{Step 5: the lattice.} By \lref{GlobalTheory.T.shortThetaSeries_span_cuspForms_weight}, which states that under chart
independence, the congruence property and the parameter conditions the short theta series span $\mathcal F_\kappa$ (here
$\kappa$ has positive last entry and $8fr+4\le n_\kappa$ by the weight conditions), and
\lref{GlobalTheory.T.weightModule_isLattice_of_shortSeries_span}, which deduces from this spanning, $s_K\ne0$, the
parameter bounds and (H2) that $\mathcal M_\kappa$ is a lattice, $(\mathrm L)$ holds.

\emph{Step 6: the page support.} The first, third, fourth and fifth items of $(\mathrm{PS})$ are part of $\mathcal C^+_{\mathrm{page}}$,
and the second is the bound $\frac12\ontooneDF(w_L)\log(g_0g_1)\le\log|\sigma(s_W)|$ of
\lref{GlobalTheory.RunIn.sOf_ne_zero_and_log_bounds}, as $s_W=s(w_L)$ (\lref{GlobalTheory.RunIn.sW_eq_sOf}).
\end{proof}

\begin{theorem}[Depth bound at the record]\label{pf:onto1:combBound}
Assume the Hodge--Tate inclusion for formal groups with jet charts at every prime. Let $D,T,h$ with $\OO\to\mathrm{LL}$ injective,
$\mathrm{EX}$, $p$, $j$ in range and the data be as in \cref{pf:onto1:inputsNet}, assume that $T$ is oriented for $\mathrm{EX}$ and
$\mathcal C^+_{\mathrm{page}}$ for $X_{T,j}$, and let $c=R(7)\ontoonePP(C_2,A_2)\ontooneDF(w_L)^2/\min(g_0,g_1)$. Then $c>0$, and every global
family datum $G$ for $X_{T,j}$ which, if $j<m$, satisfies the auxiliary form condition for $(i_\eta,\ontooneet)$, satisfies
$(\mathrm{DB})$ for $(X_{T,j},S_{\mathrm{rec}},c,\nu_T)$.
\formallabel{GlobalTheory.T.combinedBound_for_runScalars}
\formallabel{GlobalTheory.RunJoin.combBound_run}
\formallabel{GlobalTheory.PtTie.combBound_alg}
\formallabel{GlobalTheory.PtTie.combBound_of_depth_le}
\end{theorem}

\begin{proof}
\emph{Step 1: run numerics.} The page conditions imply the \emph{run numerical conditions} for $(w_L,c)$: $w_L$ has positive
last entry, $2fr\le n_{w_L}$, $4\ontooneDF(w_L)+4\le g_i$, $c\ge0$, and, if $j<m$, $\ontooneDF(w_L)\cdot2\sqrt{2fm}<g_i$ and an explicit
upper bound $r_0\bigl((2fm)^4n_{w_L}+2(2fm)^3\bigr)\ontooneDF(w_L)^2/g_i\le c$; moreover $c>0$. These are numerical consequences of
$\mathcal C^+_{\mathrm{page}}$ (\lref{GlobalTheory.RunJoin.PageOK.runNum}, which states them).

\emph{Step 2: the depth of the scaled theta lattice.} Let $\phi_C(\mathfrak l,y)\ge0$ be the chart function of
\cref{pf:onto1:PJ}(v); it vanishes if $j\ge m$. By \lref{GlobalTheory.T.thetaLattice_depth_le_unscaled_add_boundary}, for every
prime $\mathfrak l$ and every point $y$,
\begin{equation}\label{pf:onto1:depthle}
\ontoonedepth_{\mathfrak l,y}\bigl(W_{s(w_L)}(w_L)\bigr)\le\nu_T(\mathfrak l)+c\,\phi_C(\mathfrak l,y).
\end{equation}
This cited result states \eqref{pf:onto1:depthle} under the Hodge--Tate inclusion, the orientation and the run numerical
conditions; its proof distinguishes the primes $\mathfrak l$ dividing $g_0g_1$
(\lref{GlobalTheory.T.thetaLatticeDepth_le_at_places_over_level}, where the Hodge--Tate inclusion is used) from the other
primes (\lref{GlobalTheory.T.thetaLatticeDepth_le_at_places_prime_to_level}).

\emph{Step 3: $(\mathrm{DB})$.} Let $\mathfrak l$ and $y$ be given. If $j<m$, the auxiliary form condition gives $i_\eta>0$ and
$v_{\mathfrak l,y}(\iota\ontooneet)\ge2i_\eta\phi_C(\mathfrak l,y)\ge0$; if this value is finite, then $\phi_C(\mathfrak l,y)\le
v_{\mathfrak l,y}(\iota\ontooneet)/(2i_\eta)$, and \eqref{pf:onto1:depthle} with $c\ge0$ gives $\ontoonedepth_{\mathfrak l,y}\le\nu_T(\mathfrak
l)+\frac c{2i_\eta}v_{\mathfrak l,y}(\iota\ontooneet)$. If $j\ge m$, $\phi_C=0$ and \eqref{pf:onto1:depthle} is the depth bound. As
$S_{\mathrm{rec}}$ has first entry $s(w_L)$, this is $(\mathrm{DB})$.
\end{proof}

\begin{theorem}[Quotient bound at the record]\label{pf:onto1:qAt}
There is a function $C_q\colon\R^2\to\R$ such that for all $D,T,h$, $j$ in range with $j<m$ and data as in
\cref{pf:onto1:inputsNet} satisfying $\mathcal C^+_{\mathrm{page}}$ for $X_{T,j}$ and $(\mathrm L)$, the quotient bound $(\mathrm Q)$
holds for $B_T$ at $\kappa$ with the constants $(C_{Rm},C_K,C_h,A_h,C_q,A_q)=(7,C_K,1,4,C_q(C_A,C_K),4)$, the exponent $T_n$ and
$c=R(7)\ontoonePP(C_2,A_2)\ontooneDF(w_L)^2/\min(g_0,g_1)$.
\formallabel{GlobalTheory.T.recursive_page_lattice_bound}
\formallabel{GlobalTheory.RunJoin.qAt_run}
\end{theorem}

\begin{proof}
\emph{Step 1: room.} By \lref{GlobalTheory.T.exists_polynomial_bound_runConstants_degree_four}, there is a function
$C_0\ge0$ of three real variables such that, under the parameter bounds, $\log\Nm\mathfrak c_K\le\mathrm L$, $j<m$, $r\ge3$, $C_R\ge0$,
$\ontoonesize\kappa\le C_Kkfb$, $T_n>0$, $c>0$ and $T_nc\le1$ (the premises of $(\mathrm Q)$), one has $z_B\ge1$ and the \emph{room
bounds} with constants $(C_0(C_A,7,C_K),4)$: the archimedean constant and the degree constant of the CM functionals of Step 3
are at most $\ontoonePP(C_0,4)\ontoonesize\kappa$, and $\log|d_E|$ plus the logarithm of their number of indices is at most
$\ontoonePP(C_0,4)$. Put $C_q(a,b)=1+3C_0(a,7,b)+1041$.

\emph{Step 2: dimension.} By \cref{pf:onto1:netrun}, $\log\dim_K\mathcal F_\kappa\le1041z_B^3\le\ontoonePP(1041,3)\ontoonesize\kappa$.

\emph{Step 3: detection by CM functionals.} For every number field $K_1$ over the coefficient field and every CM arc datum
$d$ with values in $K_1$ in the family attached to $(j,\kappa,7,T_n,c)$, let $\phi_d$ be the coefficient functional of $d$ on
$\mathcal F_\kappa$, with values in $K_1$. These functionals vanish on $\mathcal M^{\mathrm{gg}}_\kappa(7,T_n,c)$
(\lref{GlobalTheory.ArcCM.killQ_of_gamAt}), and a form on which all of them vanish lies in $\mathcal M^{\mathrm{gg}}_\kappa(7,T_n,c)$
(\lref{GlobalTheory.T.cmFamilyFunctionals_detect_lowCoefficients}); both statements use the congruence property and known
theorems on CM abelian varieties (\lref{GlobalTheory.T.arithmeticCMKnownTheorems_proof}). Each $\phi_d$ is $K$-linear and
satisfies explicit integrality bounds on $\mathcal M_\kappa$ at the finite places and explicit bounds at the archimedean places in
terms of the sup norms (\lref{GlobalTheory.T.cmArcDatumFamily_functional_bounded_linear_on_cuspForms}).

\emph{Step 4: the quotient bound.} By \lref{GlobalTheory.T.quotientBound_of_kill_detect_constantRoom}, which states that
the vanishing, the detection, the boundedness of the CM functionals, $(\mathrm L)$, the room bounds with constants
$(C_0,A_0)$ and a dimension bound $\log\dim\mathcal F_\kappa\le\ontoonePP(C_d,A_d)\ontoonesize\kappa$ imply $(\mathrm Q)$ with the constants
$(C_h+3C_0+C_d,\max(A_h,A_0,A_d))$, the quotient bound holds with $(1+3C_0(C_A,7,C_K)+1041,\max(4,4,3))=(C_q(C_A,C_K),4)$. The
proof of the cited result embeds the quotient $Q_\kappa$ by finitely many bounded CM functionals into a product of number
fields and bounds the lengths of $Q_\kappa/\sum_i\OO_K\bar\theta_i$ by comparing the integrality of the functionals at the
finite places with their archimedean bounds via the product formula.
\end{proof}

\begin{theorem}[Shares at the record]\label{pf:onto1:shares}
Let $D,T,h$, $\mathrm{EX}$, $p$, $j$ in range and the data be as in \cref{pf:onto1:inputsNet}, let $S$ be a scalar quadruple and
$\nu$ a function on the primes, and let $c=R(7)\ontoonePP(C_2,A_2)\ontooneDF(w_L)^2/\min(g_0,g_1)$. Assume
$\mathcal C_{\mathrm{par}}(T,\mathrm{EX},p,j)$, $\mathcal C^+_{\mathrm{page}}$ for $X_{T,j}$, $(\mathrm H)$ for $(X_{T,j},S,\nu)$ with $C_R=7$, and
$(\mathrm L)$; let $G_1$ be an admissible generating family for $(S,c)$ which, if $j<m$, satisfies the auxiliary form
condition, and let $\mathrm{Sh}\subseteq\mathcal F_{w_F}$ span $\mathcal F_{w_F}$. Then there are $\Delta\in\N$ and
$\theta'_0,\dots,\theta'_{\Delta-1}\in\mathrm{Sh}$ such that, for all values of the components $r_z,\alpha,\theta^0,z,r_q,\theta,i_\eta,\ontooneet$,
the global family datum $G$ obtained from $G_1$ by inserting the $\theta'_i$ and these components satisfies $(\mathrm B_{\mathfrak l})$
and $(\mathrm M'_{\mathfrak l})$ (with $(7,c)$) for all $\mathfrak l$, $(\mathrm{Sh}_\Delta)$, $(\mathrm{Sh}_R)$ and $(\mathrm{Sh}_q)$ for
$\rho_P(7,T_n,c)$ (all with $T_n$).
\formallabel{GlobalTheory.InputsAtRecord.sharesAt_all}
\formallabel{GlobalTheory.T.cuspFormBasisShares_indefinite_liftIndex}
\formallabel{GlobalTheory.RunJoin.sharesAt_of_lt}
\formallabel{GlobalTheory.Shares.exists_inBasis_shares_inContain}
\end{theorem}

\begin{proof}
Both cases rest on an \emph{analytic model} of the zero-dimensional section at a fixed good embedding $\sigma_0$: a complex
model of the moduli space with spaces $\mathrm{sp}(u)$ of sections of weight $u$, complex coefficients $t$ for the generic
sections $s_0,\dots,s_\delta$ built from the $u_k$, and points $x_1,\dots,x_\Delta$. For a weight $u$ and $a\in\N$ let
$\mathrm{deep}(u,a)$ be the span of the products $s^\alpha g$ with $|\alpha|=a$ (monomials in $s_1,\dots,s_\delta$) and $g$ a section of
weight $u$. The \emph{analytic facts} include: $s_0(x_q)\ne0$ and $s_i(x_q)=0$ for $i\ge1$; every vector of values at the
$x_q$ is attained by a section of weight $w_F$; a section $g$ of weight $w_F+kw_L$ vanishing at all $x_q$ satisfies
$s_0^{m+1}g\in\mathrm{deep}(w_F+(k+m)w_L,1)$ for some $m$; for sections $\theta_1,\dots,\theta_n$ of weight $w_F$ with linearly independent
value vectors, the leading products $s^\alpha s_0^{T_n-|\alpha|}\theta_i$ ($|\alpha|<a_0$) are linearly independent and their span meets
$\mathrm{deep}(w_0,a_0)$ trivially; $\dim\mathrm{sp}(\kappa)\le\Delta\binom{T_n+2\delta+6}\delta$; $\Delta\binom{T_n+\delta}\delta\le\dim
\mathrm{sp}(\kappa,\mathcal M^{\mathrm{gg}})$; and, if the value vectors of $\theta_1,\dots,\theta_n$ span $\C^\Delta$, the sections of weight
$\kappa$ coming from $\mathcal M^{\mathrm{gg}}_\kappa$ lie in the span of the leading products plus $\mathrm{deep}(w_0,a_0)$.

\emph{Case $j<m$.} The weight conditions give $\kappa=w_F+T_nw_L$, $\delta\ge2$ and $11\delta^2\le T_n$
(\lref{GlobalTheory.ExtraWeight.numbers}). By \lref{GlobalTheory.T.thetaArchimedeanData_indefiniteModels}, the
admissible family $G_1$ (with its auxiliary form) provides archimedean theta data on the indefinite models, and by
\lref{GlobalTheory.T.cuspReadings_span_model_forms_shifted_weights} the realisations of the forms of every weight
$w_F+kw_L$ span the corresponding model sections, and those of $\mathcal M^{\mathrm{gg}}_\kappa(7,T_n,c)$ span the model sections with
vanishing low Fourier--Jacobi coefficients. From these, \lref{GlobalTheory.T.analyticFacts_of_thetaData_on_models} gives the
analytic facts. By \lref{GlobalTheory.Shares.exists_valueBasis_in_spanning}, since $\mathrm{Sh}$ spans $\mathcal F_{w_F}$ and the
values at the $x_q$ can be prescribed, there are $\theta'_0,\dots,\theta'_{\Delta-1}\in\mathrm{Sh}$ whose value vectors form a basis
of $\C^\Delta$. For the datum $G$ obtained by inserting them:
\begin{itemize}
\item $(\mathrm B_{\mathfrak l})$ holds for every $\mathfrak l$, because it is equivalent to the linear independence and spanning
of the value vectors (\lref{GlobalTheory.T.inBasis_iff_pointValues_linearIndependent_span});
\item $(\mathrm{Sh}_\Delta)$ follows from $(\mathrm B_{\mathfrak l})$ at one prime, the independence of the leading products and
$(\mathrm L)$ (\lref{GlobalTheory.Shares.shareRd_of_inBasis}); note that there are $\binom{T_n-1}\delta\Delta$ leading products;
\item $(\mathrm{Sh}_R)$ and $(\mathrm{Sh}_q)$ follow from the dimension bounds of the analytic facts, $\delta\ge2$, $11\delta^2\le T_n$ and
$(\mathrm L)$ (\lref{GlobalTheory.Shares.shareRhoR_of_anFacts} and \lref{GlobalTheory.Shares.shareRhoP_of_anFacts});
\item $(\mathrm M'_{\mathfrak l})$ follows from the containment of the analytic facts, transported to the generic section at
$\mathfrak l$ (\lref{GlobalTheory.Shares.Contain.inContainMR_of_anFacts}).
\end{itemize}
These statements are first obtained for the valued form algebra with all valued points and the own norm of the record,
and transported to $X_{T,j}$ (the conditions involved do not depend on the points or on the norm).

\emph{Case $j=m$.} Here $q_m=0$. By \lref{GlobalTheory.T.cuspFormBasisShares_definite_liftIndex}, the same conclusion
holds: its proof obtains archimedean theta data on the definite models
(\lref{GlobalTheory.T.thetaArchimedeanData_definiteModels_nonempty}), the spanning of the readings as above, the analytic
facts for the definite models (\lref{GlobalTheory.SharesDefinite.exists_anFacts_definite}), and concludes as in the first
case (\lref{GlobalTheory.SharesDefinite.sharesAt_of_anFacts}).
\end{proof}

\begin{theorem}[Local point data and heights at the record]\label{pf:onto1:primeData}
Let $D,T,h$, $j$ in range and the data be as in \cref{pf:onto1:inputsNet}, let $S$ be a scalar quadruple, $\nu$ a function on
the primes, $C_L>0$ and $c$ as in \cref{pf:onto1:shares}. Assume that the coefficient field is algebraically closed in the
function fields of the moduli spaces at $j$, $\mathcal C^+_{\mathrm{page}}$ for $X_{T,j}$, $(\mathrm H)$ for $(X_{T,j},S,\nu)$ with $C_R=7$,
and $(\mathrm A)$ with $(C_L,2)$ at all weights. Let $G$ be an admissible generating family for $(S,c)$ which, if $j<m$,
satisfies the auxiliary form condition, and assume $(\mathrm G2)$ with $G^{\max}$, $(\mathrm B_{\mathfrak l})$ for all $\mathfrak l$,
$\log\Delta\le1041z_B^3$ and $(\mathrm{DB})$ for $(S,c,\nu)$. Then there is a family $Q$ of local point data for $G$, regarded over
$X^+_{T,j}$, satisfying $(\mathrm C_{\mathfrak l})$, $(\mathrm T_{\mathfrak l})$ and $(\mathrm E'_{\mathfrak l})$ for all $\mathfrak l$, and $(\mathrm{Ht})$
with $(G_c,P_H,P_{rw},P_{et})=(G^{\max},\mathrm{ph}(C_L,w_L),38\mathrm L+\mathrm{PO}+1,\ontoonePP(200000,2))$.
\formallabel{GlobalTheory.T.exists_primeData_curve_lifting_heightBounds}
\formallabel{GlobalTheory.RunPrimeData.exists_primeData_heights}
\formallabel{GlobalTheory.RunPrimeData.exists_primeData_heights_of}
\end{theorem}

\begin{proof}
\emph{Step 1: the generic section.} The generic frames of the moduli space at $j$ exist and the set of components is finite.
By \lref{GlobalTheory.T.page_hypotheses_at_cusp}, there is a form $g_\circ\in\mathcal F_{w_F}$ such that the \emph{curve
hypotheses} hold for the generic section attached to $G$ and $g_\circ$; by \lref{GlobalTheory.Reg.genValInj_record} and
\lref{GlobalTheory.Reg.compSplit_record} the valuations at generic points are injective and the components split; by
\lref{GlobalTheory.Reg.birSt_of_ratios_trdeg} the ratios of forms define a birational map (for every family satisfying the
curve hypotheses).

\emph{Step 2: realisation.} By \lref{GlobalTheory.T.exists_realisationPrimeData_curve_lifting}, which states that under
the curve hypotheses, the properties of Step 1, $n_{w_L}>0$ and, if $j<m$, the non-vanishing of the auxiliary form on every
component, there are, for every prime $\mathfrak l$, valued points $y_\zeta$ of $X^+_{T,j}$ attached to the points $\zeta$ of the
generic section such that the local point data $Q_{\mathfrak l}$ with these points and the form $g_\circ$ satisfy $(\mathrm C_{\mathfrak l})$,
$(\mathrm T_{\mathfrak l})$ and $(\mathrm E'_{\mathfrak l})$ with $(G^{\max},38\mathrm L+\mathrm{PO}+1,T_n)$, we obtain $Q$.

\emph{Step 3: heights.} If $\Delta=0$, then $\mathrm{SQ}(w_F)=0$ by $(\mathrm B_{\mathfrak l})$, hence $\OO_Z=0$ by $(\mathrm C_{\mathfrak l})$,
there are no points, all $I_{\mathfrak l}$ and $\mathrm{Net}_{\mathfrak l}$ vanish, and $(\mathrm{Ht})$ reduces to $0\le0$
(\lref{GlobalTheory.RunJoin.NoPoint.heights_no_point}). Let $\Delta>0$. The page conditions provide the numerical
inputs: the bounds for $\lambda_R(C_L,w_L)$ and the constant of the lower bound $(\mathrm A)$ at the weight
$8((\delta+2)N)^2(\delta+1)\Delta(1+\lceil\mathrm{PO}\rceil)w_L$ (\lref{GlobalTheory.RunJoin.lamConst_lamR_rows}), the field bounds,
$\log N+\frac12\le P_{rw}(\ontoonesize w_L-1)$ (from admissibility), the integral forms $t_1,t_2$ of (H6) with their sup norm
bounds, the prime sum bound (H4), and the algebraic form of $(\mathrm{DB})$ (\lref{GlobalTheory.RunJoin.combAlg_of_tiedW}). If
$j<m$, \lref{GlobalTheory.T.heightInequalities_indefinite_liftIndex} gives $(\mathrm{Ht})$ from these inputs, the auxiliary
form condition and the lower bound at frames (iv) of admissibility; if $j=m$,
\lref{GlobalTheory.T.heightInequality_definite_liftIndex} gives it, the lower bound at frames being used in the form
$-\log(\sum_k\|u_k\|_x^2)^{1/2}\le p_{14}(w_L)\ontooneDF(w_L)$, since the frame function $\phi_{FD}$ vanishes at the last index
(\lref{GlobalTheory.phiFD_eq_zero_at_last_index}).
\end{proof}

%----------------------------------------------------------------------------------------------------
\subsection{Lower bounds at frames}\label{sec:pf:onto1:frames}

In this subsection we prove the lower bound at frames used in Step 6 of the proof of \cref{pf:onto1:genok}, together with the
two results on theta functions on complex tori on which it rests. Throughout, $E$ is the host CM field.

\begin{definition}[Principal tori and theta sections]\label{pf:onto1:ptorus}
A \emph{principal torus} $X$ over $E$ consists of a complex vector space $\mathrm{V}=\C^{d}$ with an action of $\OO_E$, a lattice
$\Lambda\subseteq\mathrm V$, a hermitian form $H$ and a semicharacter $\varepsilon\colon\Lambda\to\C^\times$ (with further
data and compatibilities of the development), together with an integer $r$. For $q,n\ge1$, $\mathrm{Sec}(q,n)$ is the space of
holomorphic functions $F$ on $\mathrm V^n$ with $F(v+\gamma)=\prod_k\varepsilon(\gamma_k)^q\,e^{q\pi H(v,\gamma)+q\pi H(\gamma,\gamma)/2}F(v)$ for
$\gamma\in\Lambda^n$; $\|F\|_q^2$ denotes the mean over a fundamental domain of the pointwise squared norm of level $q$, and
$\langle\cdot,\cdot\rangle_q$ the corresponding inner product. For $b\in\mathrm V^n$ put $(U_q(b)F)(v)=e^{q\pi H(v,b)-q\pi H(b,b)/2}F(v-b)$, and
for a $\C$-linear endomorphism $\pi$ of $\mathrm V$ and $s\in\mathrm{Sec}(1,1)$ put $s_0(\pi,s)(v)=\prod_ks(\pi v_k)$. We say that $\pi$ is \emph{of level $q$} if it commutes
with the action of $\OO_E$, maps $\Lambda$ into $\Lambda$, satisfies $H(\pi v,\pi w)=qH(v,w)$ and $\varepsilon(\pi x)=\varepsilon(x)^q$. A
finite set $K\subseteq\mathrm V^n$ is a \emph{set of representatives} for $(q,n,\pi)$ if $qb_k\in\Lambda$ for $b\in K$ and every
$b\in\mathrm V^n$ with $qb_k\in\Lambda$ for all $k$ is congruent to a unique $b'\in K$ in the sense that $\pi(b_k-b'_k)\in\Lambda$ for
all $k$.

A \emph{block} is a polynomial in the coordinates of $\mathrm V^n$ attached to a set of columns; for a family $\mathbf b$ of
blocks with disjoint columns, $P_1(\mathbf b)$ is their product, $\mathrm{DF}(\mathbf b)$ its degree, and $\partial_{\mathbf b}F$
denotes the derivative of $F$ along $P_1(\mathbf b)$ in the fixed direction of the development. \emph{Parameters} $P$ of degree
$D$ for $X$ consist of a level $q(P)\ge1$, an integer $m(P)$ and a real number $\mathrm{ph}(P)$ (among others), and
\[
J_P(n)=4^{-n}\bigl(6(2d^*)^{d^*/2}\bigr)^{-n}\,2^{-(r+2)D}\exp\bigl(-2\pi nC_\chi\,\mathrm{ph}(P)/q(P)\bigr),
\]
where $d^*$ and $C_\chi$ are explicit constants attached to $m(P)$.
\formal{ReciprocityTheory.ChartRestriction.G.I2.PT}
\formal{ReciprocityTheory.ChartRestriction.G.I2.IsPi}
\formal{ReciprocityTheory.ChartRestriction.G.I2.IsReps}
\formal{ReciprocityTheory.ChartRestriction.G.I2.Par}
\formal{ReciprocityTheory.ChartRestriction.G.I2.Par.Jbound}
\end{definition}

\begin{theorem}[Derivatives of translates of theta functions]\label{pf:onto1:Iholds}
Let $X$ be a principal torus over $E$, $n\in\N$, $\mathbf b$ a family of blocks with disjoint columns, $P$ parameters of degree
$\mathrm{DF}(\mathbf b)$ for $X$ with level $q=q(P)$, $\pi$ of level $q$, $0\ne s\in\mathrm{Sec}(1,1)$ and $K$ a set of representatives
for $(q,n,\pi)$. Then
\[
J_P(n)\,\|s\|_1^{2n}\le\sum_{b\in K}\bigl|\partial_{\mathbf b}\bigl(U_q(b)s_0(\pi,s)\bigr)(0)\bigr|^2 .
\]
\formallabel{GlobalTheory.ThetaTorus.I_holds}
\formallabel{ReciprocityTheory.ChartRestriction.G.I2.iOf}
\end{theorem}

\begin{proof}
We use two results proved in the development:
\begin{enumerate}[label=(\alph*)]
\item (\lref{GlobalTheory.ThetaTorus.N_holds}) the sections $U_b:=U_q(b)s_0(\pi,s)$, $b\in K$, lie in $\mathrm{Sec}(q,n)$, are
pairwise orthogonal, satisfy $\|U_b\|_q^2=\|s\|_1^{2n}$, and span $\mathrm{Sec}(q,n)$;
\item (\lref{GlobalTheory.ThetaTorus.J_holds}) there is $F\in\mathrm{Sec}(q,n)$ with $\|F\|_q=1$ and $J_P(n)\le|\partial_{\mathbf b}F(0)|^2$.
\end{enumerate}
By (a) write $F=\sum_{b\in K}\lambda_bU_b$. By orthogonality, $1=\|F\|_q^2=\|s\|_1^{2n}\sum_b|\lambda_b|^2$; in particular $\|s\|_1>0$. The
sections $U_b$ are holomorphic, so $\partial_{\mathbf b}F(0)=\sum_b\lambda_b\,\partial_{\mathbf b}U_b(0)$, and by the Cauchy--Schwarz inequality
\[
J_P(n)\le|\partial_{\mathbf b}F(0)|^2\le\Bigl(\sum_b|\lambda_b|^2\Bigr)\sum_b|\partial_{\mathbf b}U_b(0)|^2=\|s\|_1^{-2n}\sum_b|\partial_{\mathbf
b}U_b(0)|^2 .\qedhere
\]
\end{proof}

\begin{remark}\label{pf:onto1:rem-Iholds}
The existence statement (b) is the analytic core of the lower bounds at frames. Its formal proof constructs the section
$F$ by splitting the principal torus greedily into a product of tori of small dimension
(\lref{GlobalTheory.T.greedy_splitting_principal_torus_exists}) and treats the factors according to their type, using a
character statement (\lref{GlobalTheory.charStatement_holds}) and three column constructions; it is not reproduced here.
\end{remark}

\begin{theorem}[Theta rules at orthogonal frames]\label{pf:onto1:thetaLower}
Assume chart independence. Let $\sigma$ be a good embedding of the coefficient field such that the unfolding at $\sigma$
satisfies the axioms of unfoldings and the test property of the development, $k$ a component of the unfolding at
$\sigma$ and $x$ a frame in its domain which is orthogonal, and let $x_R>0$ be the common absolute value of the conjugates of the
element $\xi_0$ of the component. Let $\alpha\in\OO_E$ with $\alpha\bar\alpha=M$ (the level) and $M_b=(\bar\alpha)$, let $P_1$ be a
harmonic block product of weight $(a,n)$, and let $R_0\in\R$ and $m\in\N$ satisfy the run conditions for $(r,\deg P_1,M,R_0,m)$ and
$p\le m$ or $q\le m$ (for the signature $(p,q)$ of the datum). Then there are a principal torus $X$ with integer $r$ and
parameters $P$ of degree $\deg P_1$ with $q(P)=M$, $m(P)=m$ and $\mathrm{ph}(P)=\phi_{FD}(x)$ (for the threshold function attached to
$(M,R_0)$), such that for every family $(f_{\mathbf q})$ indexed by the tuples $\mathbf q$ of level representatives, where $f_{\mathbf q}$
is a theta rule attached to $P_1$ and $\mathbf q$,
\[
\Bigl(\frac{2x_R}M\Bigr)^{\deg P_1}J_P(n)\le\sum_{\mathbf q}|f_{\mathbf q}(x)|^2 .
\]
\formallabel{GlobalTheory.ThetaLower.thetaRule_sum_lower}
\formallabel{GlobalTheory.T.unfoldedThetaRule_archimedeanValue_bound}
\end{theorem}

\begin{proof}
\emph{Step 1: the torus at $x$.} At the orthogonal frame $x$ the hermitian matrix of the component is the scalar matrix
$2x_R\cdot\mathrm{Id}$ (\lref{GlobalTheory.HermDiag.hm_eq_diagonal}). From the period data of the component one obtains a
principal torus $X$ of dimensions given by the multiplicities of the datum, with a linear isomorphism $\psi$ from the space of
the component onto $\mathrm V$ carrying the periods onto $\Lambda$ (\lref{GlobalTheory.ScalarTorus.exists_torusOf_rows} and
\lref{GlobalTheory.OrthoPeriodTorus.isTorusAt_of_isTorusOf}); the run conditions provide parameters $P$ with the stated
values (the phase $\mathrm{ph}(P)$ being $\phi_{FD}(x)$). After renaming the variables, $P_1$ becomes a product $P_1(\mathbf b)$ of blocks with
disjoint columns, with $\mathrm{DF}(\mathbf b)=\deg P_1$ (\lref{GlobalTheory.BlockRen.blockProd_rename}).

\emph{Step 2: the theta section.} The theta function of the component, transported by $\psi$, is a section $0\ne s\in\mathrm{Sec}(1,1)$
with $\|s\|_1=1$ (\lref{GlobalTheory.EtaAn.ThetaTorus.theta_section_l2_torus}); multiplication by $\alpha$ is of level $M$ on $X$
(\lref{GlobalTheory.AnVal.IsPiAct.isPi_actL_of_isTorusOf}), and the translates by the level representatives give a set of
representatives $K$ for $(M,n,\alpha)$ together with a surjection from the tuples of level representatives onto $K$
(\lref{GlobalTheory.isReps_level_translates}).

\emph{Step 3: values of theta rules.} For a tuple $\mathbf q$ with image $b\in K$, the value $f_{\mathbf q}(x)$ is computed through the
periods of the component (\lref{GlobalTheory.AnVal.thetaRule_value_periods}) and transported to the torus
(\lref{GlobalTheory.ThetaLower.value_torus}); comparing the normalisations (\lref{GlobalTheory.ThetaLower.exists_kappa}) one
obtains $(2x_R/M)^{\deg P_1}\,|\partial_{\mathbf b}(U_M(b)s_0(\alpha,s))(0)|^2\le|f_{\mathbf q}(x)|^2$.

\emph{Step 4.} By \cref{pf:onto1:Iholds} with $\|s\|_1=1$, $J_P(n)\le\sum_{b\in K}|\partial_{\mathbf b}(U_M(b)s_0(\alpha,s))(0)|^2$. Multiplying by
$(2x_R/M)^{\deg P_1}$, using Step 3 and the surjectivity of $\mathbf q\mapsto b$ (each $b\in K$ is the image of some tuple, and the other
terms are non-negative) gives the assertion.
\end{proof}

\begin{theorem}[Lower bound at frames]\label{pf:onto1:frames}
Assume chart independence. Let $D,T,h$, $\mathrm{EX}$, $j$ in range and the data be as in \cref{pf:onto1:inputsNet}, let $S$ be a
scalar quadruple, and assume the page hypotheses $\mathcal H_{\mathrm{page}}$ (with $C_R$), $(\mathrm K)$, $(\mathrm{PA})$ and $(\mathrm{PS})$ for
$(S,w_L)$. Let $P_1$ be a harmonic block product of weight $w_L$, let $w_0,\dots,w_N\in\mathcal F_{w_L}$ be a family containing the
theta rules of $P_1$ scaled by $s_W$ (for every tuple $\mathbf q$ of level representatives some $w_{k(\mathbf q)}$ equals $s_W$ times a
theta rule attached to $P_1$ and $\mathbf q$, injectively in $\mathbf q$), and let $t\ge p^\circ_{14}(w_L)$. Then for every good embedding
$\sigma$, every component and every frame $x$ in its domain,
\[
-\log\Bigl(\sum_k\|w_k\|_x^2\Bigr)^{1/2}\le t\,\ontooneDF(w_L)+c\,\phi_{FD}(x),\qquad c=R(C_R)\ontoonePP(C_2,A_2)\ontooneDF(w_L)^2/\min(g_0,g_1).
\]
\formallabel{GlobalTheory.T.lower_bound_at_frames_family_with_scaled_thetaRules}
\formallabel{GlobalTheory.DegreeInterpolationBounds.i2iAt_record}
\formallabel{GlobalTheory.AnVal.i2iAt_of_containsTh}
\end{theorem}

\begin{proof}
\emph{Step 1: inputs.} The laws provide the frame data at the index $j$, and there is $\alpha\in\OO_E$ with $\alpha\bar\alpha=M=g_0g_1$ and
$M_b=(\bar\alpha)$ (\lref{GlobalTheory.RunIn.half_level_generator_norm}). By
\lref{GlobalTheory.T.thetaRunParameters_at_typePrimes_product}, the page hypotheses, $(\mathrm K)$, $(\mathrm{PA})$ and $(\mathrm{PS})$
imply the run conditions for $(r,\ontooneDF(w_L),M,R_0,m)$ with $R_0=\sqrt{90(fr+\ontooneDF(w_L))}$ and an explicit bound relating $c$ to
the run parameters. Moreover $\deg P_1=\ontooneDF(w_L)$ (\lref{GlobalTheory.WeightExponents.weight_exponents}), $q_j\le m$, and by
$(\mathrm{PS})$, $\frac12\ontooneDF(w_L)\log(g_0g_1)\le\log|\sigma(s_W)|$, i.e.\ $M^{\deg P_1}\le|\sigma(s_W)|^2$ for every $\sigma$.

\emph{Step 2: reduction to an orthogonal frame.} Fix $\sigma$, a component and $x$. There is a frame change $b$ such that $xb$ is
orthogonal, $\phi_{FD}(xb)=\phi_{FD}(x)$ and the normalised frame norms at $x$ are computed at $xb$
(\lref{GlobalTheory.AnVal.exists_ortho_frame}): $\|w\|_x=\Lambda\cdot|w(xb)|$ with $\Lambda^2=(2\pi x_R)^{-\deg P_1}$.

\emph{Step 3: the estimate.} By \cref{pf:onto1:thetaLower} at the orthogonal frame $xb$ we get $X$, $P$ and
$(2x_R/M)^{\deg P_1}J_P(n)\le\sum_{\mathbf q}|f_{\mathbf q}(xb)|^2$ for the theta rules $f_{\mathbf q}$ with $w_{k(\mathbf q)}=s_Wf_{\mathbf q}$. Hence
\begin{multline*}
\sum_k\|w_k\|_x^2\ge\sum_{\mathbf q}\Lambda^2|\sigma(s_W)|^2|f_{\mathbf q}(xb)|^2\\
\ge(2\pi x_R)^{-\deg P_1}M^{\deg P_1}\Bigl(\frac{2x_R}M\Bigr)^{\deg P_1}J_P(n)=\pi^{-\deg P_1}J_P(n).
\end{multline*}
Finally, inserting the formula for $J_P(n)$, the bound $t\ge p^\circ_{14}(w_L)$ and the bound relating $c$ to the run parameters
gives $-\frac12\log\bigl(\pi^{-\deg P_1}J_P(n)\bigr)\le t\,\ontooneDF(w_L)+c\,\mathrm{ph}(P)$ (\lref{GlobalTheory.AnVal.neg_log_sqrt_le_of_jbound}), and
$\mathrm{ph}(P)=\phi_{FD}(xb)=\phi_{FD}(x)$.
\end{proof}
